% Leçon 3 — Expression écrite : situation de la femme — Chinois Tle A — Cours signé OnBuch+
% Programme officiel MINESEC. Compiler avec : tectonic ZH03-ecrit-situation-de-la-femme.tex
\newcommand{\DOCMATIERE}{Chinois}
\newcommand{\DOCNIVEAU}{Tle A}
\newcommand{\DOCMODULE}{Module 1 : Vie familiale et intégration sociale}
\newcommand{\DOCLECON}{ZH03}
\newcommand{\DOCTITRE}{Leçon 3 — Expression écrite : situation de la femme}
% OnBuch+ — préambule « cours signés » ENRICHI (compilé avec tectonic / XeLaTeX)
% Système visuel « Pop » d'OnBuch+ : fond crème (jamais blanc), bordures encre
% épaisses, ombres « dures » décalées, titres Archivo Black en capitales.
% Version MATHÉMATIQUES : graphes (pgfplots/tikz), boîtes pédagogiques
% (définition, propriété, méthode, attention, le savais-tu, exercice résolu,
% exercice, TP, bilan), page de garde et signature OnBuch+ ancrée en bas.
%
% Macros attendues dans main.tex AVANT \input{preamble} :
%   \DOCMATIERE  \DOCNIVEAU  \DOCMODULE  \DOCLECON (numéro)  \DOCTITRE
\documentclass[11pt]{article}

\usepackage{fontspec}
\usepackage[french]{babel} % français = langue principale ; le chinois via \zh{..} / textebox (xeCJK)
\usepackage{xeCJK}
\usepackage{pifont} % \ding{51} \ding{55} utilisés par les modèles
\usepackage[a4paper,margin=2cm,top=3.4cm,bottom=2.8cm,headheight=1.4cm,footskip=1.3cm]{geometry}
\usepackage{amsmath,amssymb,amsfonts}
\usepackage{mathtools} % \xmapsto, \coloneqq... souvent utilisés par les modèles
\usepackage{cancel} % simplifications barrées (bilans de forces)
% \abs{z} (module, valeur absolue) et \norm{u} (norme), utilisés par les modèles
\providecommand{\abs}{}\let\abs\relax\DeclarePairedDelimiter{\abs}{\lvert}{\rvert}
\providecommand{\norm}{}\let\norm\relax\DeclarePairedDelimiter{\norm}{\lVert}{\rVert}
\usepackage[table]{xcolor} % [table] : fournit aussi \rowcolors (alternance de lignes)
\usepackage{pagecolor}
\usepackage{tikz}
\usetikzlibrary{arrows.meta,positioning,calc,shapes.geometric,shapes.misc,shapes.arrows,decorations.pathmorphing,decorations.pathreplacing,decorations.markings,patterns,shadows,mindmap,trees,fit,backgrounds,matrix,intersections,angles,quotes}
\usepackage{pgfplots}
\usepackage[version=4]{mhchem} % équations nucléaires \ce{^{235}_{92}U}
\usepackage[european,siunitx]{circuitikz} % schémas électriques
\pgfplotsset{compat=1.18}
\usepgfplotslibrary{fillbetween,groupplots} % aires entre courbes (calcul intégral)
\usepackage{enumitem}
\usepackage{fancyhdr}
% [most] charge toutes les bibliothèques tcolorbox (listings, minted, raster,
% documentation...) et gonfle inutilement la mémoire TeX sur un document long
% (~300 boîtes) : on ne charge que ce qui est réellement utilisé.
\usepackage[skins,breakable]{tcolorbox}
\usepackage{graphicx}
\usepackage{colortbl}
\usepackage{booktabs}
\usepackage{tabularx}
\usepackage{diagbox} % case d'en-tête diagonale des tableaux à double entrée
% Colonnes extensibles centrée / gauche / droite (C, L, R), souvent utilisées par les modèles
\newcolumntype{C}{>{\centering\arraybackslash}X}
\newcolumntype{L}{>{\raggedright\arraybackslash}X}
\newcolumntype{R}{>{\raggedleft\arraybackslash}X}
\usepackage{array}
\usepackage{multicol}
\usepackage{multirow}
\usepackage{siunitx}
% parse-numbers=false : accepte aussi \SI{2,5 \times 10^{-3}}{...}
\sisetup{locale=FR,output-decimal-marker={,},group-minimum-digits=4,parse-numbers=false}
\DeclareSIUnit\ohm{\ensuremath{\symup{\Omega}}} % U+2126 absent de la police
\DeclareSIUnit\division{div} % réglages d'oscilloscope (ms/div, V/div)
\DeclareSIUnit\tr{tr} % tours (vitesse de rotation, ex. \SI{1500}{\tr\per\minute})
\DeclareSIUnit\molar{M} % concentration molaire (ex. \SI{3}{\molar}, \SI{1}{\micro\molar})
\DeclareSIUnit\an{an} % année (le modèle confond parfois avec \year, un compteur TeX) — ex. \SI{5}{\kilo\meter\per\an}
\DeclareSIUnit\FCFA{FCFA} % franc CFA (ex. \SI{500}{\FCFA\per\kilo\gram})
\DeclareSIUnit\degreeBrix{\textdegree Brix} % teneur en sucre d'un sirop/jus (réfractométrie)
\DeclareSIUnit\month{mois} % le modèle utilise parfois \month, un compteur TeX, comme unité de durée
\usepackage{qrcode}
% \degree : commande fréquemment utilisée par les modèles pour un angle ou une
% température en dehors de \SI{}{} (non fournie par les paquets chargés).
\providecommand{\degree}{^{\circ}}
% \qed (fin de démonstration), utilisé par les modèles sans amsthm
\providecommand{\qed}{\hfill$\square$}
% \barre : double barre des valeurs interdites dans un tableau de variations (tkz-tab)
\providecommand{\barre}{\ensuremath{\|}}
% environnement proof (amsthm non chargé) utilisé par les modèles
\ifdefined\proof\else\newenvironment{proof}[1][Démonstration]{\par\noindent\textit{#1.}\ }{\qed\par}\fi
\usepackage{lastpage}
\usepackage{hyperref}
\hypersetup{hidelinks}

% ============================================================
% Palette « Pop » OnBuch+ — mêmes hex que la classe P (plus_ui.dart)
% ============================================================
\definecolor{poporange}{HTML}{F59321}
\definecolor{popdark}{HTML}{E07A0C}
\definecolor{popcream}{HTML}{FFF6E8}
\definecolor{popink}{HTML}{1C1714}
\definecolor{poppurple}{HTML}{7A5AE0}
\definecolor{popgreen}{HTML}{2E9E6A}
\definecolor{poppink}{HTML}{E0457B}
\definecolor{popblue}{HTML}{2F7DE1}
\definecolor{popgold}{HTML}{C98A12}
\definecolor{popmuted}{HTML}{8A7F76}
\definecolor{popline}{HTML}{EADFD2}
\definecolor{popgray}{HTML}{8A7F76}
\colorlet{popgrey}{popgray}
% Alias tolérants (le modèle nomme parfois les couleurs différemment)
\colorlet{popwhite}{white}
\colorlet{popblack}{popink}
\colorlet{popred}{poppink}
\colorlet{popyellow}{popgold}
% Teintes claires (fonds de tableaux/encadrés) — le modèle nomme parfois
% les couleurs avec un suffixe L (ex. popblueL) sans les avoir définies.
\colorlet{poporangeL}{poporange!20}
\colorlet{popdarkL}{popdark!20}
\colorlet{poppurpleL}{poppurple!20}
\colorlet{popgreenL}{popgreen!20}
\colorlet{poppinkL}{poppink!20}
\colorlet{popblueL}{popblue!20}
\colorlet{popgoldL}{popgold!20}
\colorlet{popredL}{popred!20}
\colorlet{popgrayL}{popgray!20}
% Teintes claires pour fonds de tableaux / zones
\colorlet{popblueL}{popblue!10!white}
\colorlet{poporangeL}{poporange!14!white}
\colorlet{popgreenL}{popgreen!12!white}
\colorlet{poppurpleL}{poppurple!12!white}
\colorlet{poppinkL}{poppink!12!white}
\colorlet{popgoldL}{popgold!14!white}

\pagecolor{popcream}

% ============================================================
% Typographie
% ============================================================
\defaultfontfeatures{Ligatures=TeX}
\setmainfont{PlusJakartaSans-Medium.ttf}[Path=../../../fonts/,BoldFont=PlusJakartaSans-Bold.ttf]
% Symboles, API et emojis absents de la police principale : repli sur DejaVu Sans (caractères actifs)
\newfontface\symfont{DejaVuSans.ttf}[Path=../../../fonts/]
\makeatletter
\newcommand\symactive[1]{\catcode#1=13 \begingroup\lccode`\~=#1 \lowercase{\endgroup\def~}{{\symfont\char#1}}}
\newcommand\symblank[1]{\catcode#1=13 \begingroup\lccode`\~=#1 \lowercase{\endgroup\def~}{}}
\newcommand\symhyph[1]{\catcode#1=13 \begingroup\lccode`\~=#1 \lowercase{\endgroup\def~}{\mbox{-}}}
\newcount\symc
\newcommand\symrange[2]{\symc=#1 \@whilenum\symc<#2 \do{\expandafter\symactive\expandafter{\the\symc}\advance\symc by 1 }}
\newcommand\symblankrange[2]{\symc=#1 \@whilenum\symc<#2 \do{\expandafter\symblank\expandafter{\the\symc}\advance\symc by 1 }}
\makeatother
\symrange{"0250}{"02FF}\symactive{"2713}\symactive{"2714}\symactive{"2717}\symactive{"2718}\symactive{"2610}\symactive{"2611}\symactive{"2612}
\symrange{"2600}{"27BF}
\catcode"2705=13 \begingroup\lccode`\~="2705 \lowercase{\endgroup\def~}{{\symfont\char"2713}}\symblank{"26BD}\symblank{"26A1}
\symrange{"2460}{"257F}\symactive{"223C}\symactive{"2248}\symactive{"2260}
\symhyph{"2011}\symhyph{"2E1A}\symblankrange{"1F300}{"1FAFF}\symblank{"FE0F}
% Caractères chinois : WenQuanYi Zen Hei (sous-ensemble GB2312 + pinyin), gras/italique simulés
\setCJKmainfont{WenQuanYiZenHei-subset.ttf}[Path=../../../fonts/,AutoFakeBold=3,AutoFakeSlant=0.2]
\xeCJKsetup{CJKspace=false}
% Pinyin : voyelles à caron (ǎ ǐ ǒ ǔ ǖ ǘ ǚ ǜ) absentes de la police principale -> caractères actifs
\newfontface\pyfont{WenQuanYiZenHei-subset.ttf}[Path=../../../fonts/]
\newcommand\pyactive[1]{\catcode#1=13 \begingroup\lccode`\~=#1 \lowercase{\endgroup\def~}{{\pyfont\char#1}}}
\pyactive{"01CD}\pyactive{"01CE}\pyactive{"01CF}\pyactive{"01D0}\pyactive{"01D1}\pyactive{"01D2}\pyactive{"01D3}\pyactive{"01D4}\pyactive{"01D5}\pyactive{"01D6}\pyactive{"01D7}\pyactive{"01D8}\pyactive{"01D9}\pyactive{"01DA}\pyactive{"01DB}\pyactive{"01DC}
% Maths en sans-serif (Fira Math) avec chiffres et lettres droites de Plus
% Jakarta, pour que les formules se fondent dans le texte.
\usepackage[math-style=ISO,bold-style=ISO]{unicode-math}
\setmathfont{FiraMath-Regular.otf}[Path=../../../fonts/]
\setmathfont{PlusJakartaSans-Medium.ttf}[Path=../../../fonts/,range={up/num,up/{Latin,latin},\mathup}]
\usepackage{icomma} % virgule décimale sans espace en mode math
% Caractères mathématiques Unicode absents des polices de texte (Plus Jakarta,
% Archivo Black) : rendus par leur équivalent mathématique, en texte comme en
% formule — sinon ils s'affichent en carré vide (ex. « Arithmétique dans ℤ »).
\usepackage{newunicodechar}
\newunicodechar{ℝ}{\ensuremath{\mathbb{R}}}
\newunicodechar{ℤ}{\ensuremath{\mathbb{Z}}}
\newunicodechar{ℂ}{\ensuremath{\mathbb{C}}}
\newunicodechar{ℕ}{\ensuremath{\mathbb{N}}}
\newunicodechar{ℚ}{\ensuremath{\mathbb{Q}}}
\newunicodechar{𝔻}{\ensuremath{\mathbb{D}}}
\newunicodechar{α}{\ensuremath{\alpha}}
\newunicodechar{β}{\ensuremath{\beta}}
\newunicodechar{γ}{\ensuremath{\gamma}}
\newunicodechar{δ}{\ensuremath{\delta}}
\newunicodechar{ε}{\ensuremath{\varepsilon}}
\newunicodechar{θ}{\ensuremath{\theta}}
\newunicodechar{λ}{\ensuremath{\lambda}}
\newunicodechar{μ}{\ensuremath{\mu}}
\newunicodechar{σ}{\ensuremath{\sigma}}
\newunicodechar{τ}{\ensuremath{\tau}}
\newunicodechar{φ}{\ensuremath{\varphi}}
\newunicodechar{ψ}{\ensuremath{\psi}}
\newunicodechar{ω}{\ensuremath{\omega}}
\newunicodechar{Σ}{\ensuremath{\Sigma}}
% majuscules produites par \MakeUppercase dans les titres (α-aminés -> Α)
\newunicodechar{Α}{\ensuremath{\alpha}}
\newunicodechar{Β}{\ensuremath{\beta}}
\newunicodechar{Γ}{\ensuremath{\Gamma}}
\newunicodechar{Δ}{\ensuremath{\Delta}}
\newunicodechar{Ε}{\ensuremath{\varepsilon}}
\newunicodechar{Θ}{\ensuremath{\Theta}}
\newunicodechar{Λ}{\ensuremath{\Lambda}}
\newunicodechar{Μ}{\ensuremath{\mu}}
\newunicodechar{Τ}{\ensuremath{\tau}}
\newunicodechar{Φ}{\ensuremath{\Phi}}
\newunicodechar{Ψ}{\ensuremath{\Psi}}
\newunicodechar{Ω}{\ensuremath{\Omega}}
\newunicodechar{↦}{\ensuremath{\mapsto}}
\newunicodechar{∈}{\ensuremath{\in}}
\newunicodechar{∉}{\ensuremath{\notin}}
\newunicodechar{∧}{\ensuremath{\wedge}}
\newunicodechar{⇔}{\ensuremath{\Leftrightarrow}}
\newunicodechar{⟺}{\ensuremath{\Longleftrightarrow}}
\newunicodechar{⇒}{\ensuremath{\Rightarrow}}
\newunicodechar{⟹}{\ensuremath{\Longrightarrow}}
\newunicodechar{∘}{\ensuremath{\circ}}
\newunicodechar{∪}{\ensuremath{\cup}}
\newunicodechar{∩}{\ensuremath{\cap}}
\newunicodechar{≡}{\ensuremath{\equiv}}
\newunicodechar{⊂}{\ensuremath{\subset}}
\newunicodechar{⊥}{\ensuremath{\perp}}
\newunicodechar{∃}{\ensuremath{\exists}}
\newunicodechar{∀}{\ensuremath{\forall}}
\newunicodechar{∛}{\ensuremath{\sqrt[3]{\,}}}
\newunicodechar{∜}{\ensuremath{\sqrt[4]{\,}}}
\newunicodechar{⃗}{\ensuremath{{}^{\to}}}
\newunicodechar{ⁿ}{\ifmmode^{n}\else\textsuperscript{\ensuremath{n}}\fi}
\newunicodechar{ˣ}{\ifmmode^{x}\else\textsuperscript{\ensuremath{x}}\fi}
\newunicodechar{ᵇ}{\ifmmode^{b}\else\textsuperscript{\ensuremath{b}}\fi}
\newunicodechar{ʳ}{\ifmmode^{r}\else\textsuperscript{\ensuremath{r}}\fi}
\newunicodechar{ᵘ}{\ifmmode^{u}\else\textsuperscript{\ensuremath{u}}\fi}
\newunicodechar{ʸ}{\ifmmode^{y}\else\textsuperscript{\ensuremath{y}}\fi}
\newunicodechar{ᵃ}{\ifmmode^{a}\else\textsuperscript{\ensuremath{a}}\fi}
\newunicodechar{ᵉ}{\ifmmode^{e}\else\textsuperscript{\ensuremath{e}}\fi}
\newunicodechar{ᵏ}{\ifmmode^{k}\else\textsuperscript{\ensuremath{k}}\fi}
\newunicodechar{ᵖ}{\ifmmode^{p}\else\textsuperscript{\ensuremath{p}}\fi}
\newunicodechar{ᴺ}{\ifmmode^{N}\else\textsuperscript{\ensuremath{N}}\fi}
\newunicodechar{ᶜ}{\ifmmode^{c}\else\textsuperscript{\ensuremath{c}}\fi}
\newunicodechar{ʰ}{\ifmmode^{h}\else\textsuperscript{\ensuremath{h}}\fi}
\newunicodechar{ᵠ}{\ifmmode^{q}\else\textsuperscript{\ensuremath{q}}\fi}
\newunicodechar{ₙ}{\ifmmode_{n}\else\textsubscript{\ensuremath{n}}\fi}
\newunicodechar{ₐ}{\ifmmode_{a}\else\textsubscript{\ensuremath{a}}\fi}
\newunicodechar{ᵢ}{\ifmmode_{i}\else\textsubscript{\ensuremath{i}}\fi}
\newunicodechar{ᵣ}{\ifmmode_{r}\else\textsubscript{\ensuremath{r}}\fi}
\newunicodechar{ₖ}{\ifmmode_{k}\else\textsubscript{\ensuremath{k}}\fi}
\newunicodechar{ᵦ}{\ifmmode_{\beta}\else\textsubscript{\ensuremath{\beta}}\fi}
\newunicodechar{ₓ}{\ifmmode_{x}\else\textsubscript{\ensuremath{x}}\fi}
\newfontfamily\popdisplay{ArchivoBlack-Regular.ttf}[Path=../../../fonts/]
\newfontfamily\poplogo{SpaceGrotesk-Bold.ttf}[Path=../../../fonts/]
\newfontfamily\popsemi{PlusJakartaSans-SemiBold.ttf}[Path=../../../fonts/]
\newfontfamily\popxbold{PlusJakartaSans-ExtraBold.ttf}[Path=../../../fonts/]
\newfontfamily\popxboldls{PlusJakartaSans-ExtraBold.ttf}[Path=../../../fonts/,LetterSpace=3.0]

\setlist[enumerate]{leftmargin=1.7em,itemsep=5pt,topsep=4pt}
\setlist[itemize]{leftmargin=1.7em,itemsep=4pt,topsep=4pt,label={\color{poporange}\textbullet}}
\setlength{\parindent}{0pt}
\setlength{\parskip}{5pt}
\renewcommand{\arraystretch}{1.25}

% pgfplots : style Pop par défaut
\pgfplotsset{
  every axis/.append style={axis line style={popink,line width=1pt},tick style={popink},
    grid style={popline,line width=0.5pt},label style={font=\small},tick label style={font=\footnotesize}},
  popcurve/.style={poporange,line width=1.8pt,no markers,smooth},
}
% Même style disponible dans un \draw TikZ simple (hors pgfplots)
\tikzset{popcurve/.style={poporange,line width=1.8pt}}
\tikzset{popgrid/.style={popline,very thin}}
\colorlet{popcurve}{poporange} % le nom du style est parfois utilisé comme couleur

% ============================================================
% Pill / squiggle
% ============================================================
\newcommand{\poppill}[3]{%
  \tikz[baseline=-0.55ex]\node[draw=popink,line width=1.2pt,rounded corners=2.6pt,%
    fill=#2,inner xsep=6.5pt,inner ysep=3pt]{%
    \popxboldls\fontsize{7}{7}\selectfont\color{#3}\MakeUppercase{#1}};%
}
\newcommand{\popsquiggle}[1]{%
  \tikz[baseline=-0.4ex]{
    \draw[#1,line width=2.6pt,line cap=round]
      plot [smooth] coordinates {(0,0.09) (0.34,0.01) (0.68,0.21) (1.02,0.01) (1.36,0.10)};
  }%
}
% Mot-clé surligné (marqueur orange sous le texte)
\newcommand{\cle}[1]{{\popxbold\color{popdark}#1}}

% ============================================================
% En-tête / pied
% ============================================================
\pagestyle{fancy}
\fancyhf{}
\renewcommand{\headrulewidth}{0pt}
\renewcommand{\footrulewidth}{0pt}
\lhead{\raisebox{-0.15ex}{\poplogo\fontsize{13}{13}\selectfont\color{popink}OnBuch\color{poporange}+}}
\rhead{\poppill{\DOCMATIERE\ \textbullet\ \DOCNIVEAU\ \textbullet\ Leçon \DOCLECON}{white}{popink}}
\lfoot{\popxbold\fontsize{7.6}{7.6}\selectfont\color{popmuted}\MakeUppercase{Cours signé OnBuch+}\ \textbullet\ {\popsemi\DOCTITRE}}
\rfoot{\tikz[baseline=-0.6ex]\node[rounded corners=8.5pt,fill=popink,minimum height=17pt,minimum width=17pt,inner xsep=5pt,inner sep=0]{\popxbold\fontsize{8}{8}\selectfont\color{popcream}\thepage\,/\,\pageref*{LastPage}};}

% ============================================================
% Titres
% ============================================================
\newcommand{\chaptitre}[2]{%
  \begin{tcolorbox}[enhanced,colback=poporange,colframe=popink,boxrule=2.6pt,arc=11pt,
      drop shadow=popink,left=15pt,right=15pt,top=12pt,bottom=11pt,before skip=3pt,after skip=18pt]
    \poppill{#2}{popink}{popcream}\par\vspace{9pt}
    {\raggedright\hyphenpenalty=10000\exhyphenpenalty=10000\popdisplay\fontsize{22}{24}\selectfont\color{popcream}\MakeUppercase{#1}\par}
  \end{tcolorbox}
}
% Section numérotée : pastille orange avec le numéro + titre Archivo
\newcounter{popsec}
\newcommand{\coursec}[1]{%
  \stepcounter{popsec}\setcounter{popsub}{0}%
  \par\vspace{16pt}\noindent
  \begin{minipage}{\linewidth}
  \tikz[baseline=-0.7ex]\node[draw=popink,line width=1.6pt,fill=poporange,rounded corners=4pt,minimum size=22pt,inner sep=0]{\popdisplay\fontsize{12}{12}\selectfont\color{popcream}\thepopsec};%
  \hspace{8pt}{\popdisplay\fontsize{15}{17}\selectfont\color{popink}\MakeUppercase{#1}}\par
  \vspace{4pt}\noindent{\color{poporange}\rule{36pt}{3.6pt}}
  \end{minipage}\par\nopagebreak\vspace{8pt}
}
\newcounter{popsub}
\newcommand{\courssub}[1]{%
  \stepcounter{popsub}%
  \par\vspace{9pt}\noindent{\popxbold\fontsize{11.5}{13}\selectfont\color{popink}{\color{poporange}\thepopsec.\thepopsub}\hspace{6pt}#1}\par\nopagebreak\vspace{4pt}
}

% ============================================================
% Boîtes pédagogiques — même recette : fond blanc, bordure épaisse colorée,
% ombre dure encre, badge plein. Titre optionnel [#1].
% ============================================================
\tcbset{popbase/.style={breakable,enhanced,colback=white,boxrule=2.3pt,arc=9pt,
  shadow={3pt}{-3pt}{0pt}{popink},left=14pt,right=14pt,top=11pt,bottom=11pt,before skip=11pt,after skip=15pt,
  fontupper=\popsemi\fontsize{10.6}{14.5}\selectfont\color{popink},
  fontlower=\popsemi\fontsize{10.6}{14.5}\selectfont\color{popink}}}
% Coche dessinée (le glyphe ✓ n'existe pas dans les polices du document)
\newcommand{\popcheck}[1]{\tikz[baseline=-0.5ex]\draw[#1,line width=1.6pt,line cap=round,line join=round](-0.13,0.0)--(-0.04,-0.09)--(0.13,0.1);}
\providecommand{\coche}{\popcheck{popgreen}}
\providecommand{\resultat}[1]{\par\smallskip\noindent\fcolorbox{popgreen}{popgreenL}{\parbox{\dimexpr\linewidth-2\fboxsep-2\fboxrule\relax}{\textbf{Résultat.} #1}}\par\smallskip}
\newcommand{\popboxhead}[3]{\poppill{#1}{#2}{white}\ifx\relax#3\relax\else\hspace{6pt}{\popxbold\fontsize{10.6}{12}\selectfont\color{popink}#3}\fi\par\vspace{7pt}}

\newtcolorbox{aretenir}[1][]{popbase,colframe=popgreen,before upper={\popboxhead{À retenir}{popgreen}{#1}}}
% Texte en chinois (document de lecture, dialogue, extrait) : cadre
% d'étude. Un titre optionnel (source, auteur) s'affiche dans l'en-tête.
\newtcolorbox{textebox}[1][]{popbase,colframe=popblue,before upper={\popboxhead{文本 · Texte}{popblue}{#1}},after upper={}}
% Marques ✓ / ✗ (forme correcte / fautive) souvent utilisées par les modèles
\providecommand{\cross}{\textcolor{poppink}{\ensuremath{\times}}}
\providecommand{\xmark}{\cross}\providecommand{\crossmark}{\cross}\providecommand{\cmark}{\ensuremath{\checkmark}}
% Ligne de réponse / justification à compléter (inventée par les modèles)
\providecommand{\justification}[1]{\par\noindent\textit{Justification~:} \dotfill\par}
% \textipa (package tipa non chargé) : la police ne couvre pas l'API -> simple passage
\providecommand{\textipa}[1]{#1}
% Soulignement (ulem non chargé) : \uline, \ul
\providecommand{\uline}[1]{\underline{#1}}\providecommand{\ul}[1]{\underline{#1}}
% \glossaire{\item ...} inventée par les modèles : liste de glossaire
\providecommand{\glossaire}[1]{\begin{itemize}#1\end{itemize}}
% \alert (beamer) inventée par les modèles : texte mis en évidence
\providecommand{\alert}[1]{\textbf{\textcolor{poppink}{#1}}}
% variantes ASCII d'ordinaux écrites en macro
\providecommand{\iere}{\textsuperscript{re}}\providecommand{\ieme}{\textsuperscript{e}}\providecommand{\ere}{\textsuperscript{re}}\providecommand{\eme}{\textsuperscript{e}}\providecommand{\er}{\textsuperscript{er}}\providecommand{\ie}{\textsuperscript{e}}
% Ordinaux français écrits en macro par les modèles : 1\ière, 3\ième
\newcommand{\ière}{\textsuperscript{re}}\newcommand{\ième}{\textsuperscript{e}}
% \exemple (inventée par les modèles) : étiquette « Exemple : »
\providecommand{\exemple}{\textbf{Exemple~:}\ }
% \square utilisable aussi hors mode math (cases à cocher)
\AtBeginDocument{\let\squareMath\square\renewcommand{\square}{\ensuremath{\squareMath}}}
% Mot ou phrase en chinois à l'intérieur d'un paragraphe français
\newcommand{\zh}[1]{\ifmmode\text{#1}\else{#1}\fi}
\let\eng\zh \let\esp\zh \let\all\zh % alias tolérés
% flèches d'intonation inventées par les modèles
\providecommand{\textsearrow}{}\renewcommand{\textsearrow}{\ensuremath{\searrow}}\providecommand{\textnearrow}{}\renewcommand{\textnearrow}{\ensuremath{\nearrow}}
% \fr{..} : mot français (inventé par les modèles)
\providecommand{\fr}[1]{\textit{#1}}
% zhbox : encadré de texte chinois inventé par les modèles
\newenvironment{zhbox}{\begin{quote}}{\end{quote}}
% \vs : « versus » inventé par les modèles
\providecommand{\vs}{\textit{vs}\ }
% \blank : case à compléter inventée par les modèles
\providecommand{\blank}[1][]{\underline{\hspace{1.6em}}}
\newtcolorbox{exemplebox}[1][]{popbase,colframe=poporange,before upper={\popboxhead{Exemple}{poporange}{#1}}}
\newtcolorbox{definition}[1][]{popbase,colframe=popblue,before upper={\popboxhead{Définition}{popblue}{#1}}}
\newtcolorbox{propriete}[1][]{popbase,colframe=poppurple,before upper={\popboxhead{Propriété}{poppurple}{#1}}}
\newtcolorbox{methode}[1][]{popbase,colframe=popgold,before upper={\popboxhead{Méthode}{popgold}{#1}}}
\newtcolorbox{attention}[1][]{popbase,colframe=poppink,before upper={\popboxhead{Attention}{poppink}{#1}}}
\newtcolorbox{experience}[1][]{popbase,colframe=popblue,colback=popblueL,before upper={\popboxhead{Expérience}{popblue}{#1}}}
% « Le savais-tu ? » : carte violette pleine (contexte, culture, Cameroun)
\newtcolorbox{savaistu}[1][]{popbase,colback=poppurple,colframe=popink,
  fontupper=\popsemi\fontsize{10.4}{14.2}\selectfont\color{white},
  before upper={\poppill{Le savais-tu\,?}{white}{poppurple}\ifx\relax#1\relax\else\hspace{6pt}{\popxbold\color{white}#1}\fi\par\vspace{7pt}}}
% Exercice résolu : énoncé en haut, \tcblower puis la solution sur fond crème
\newcounter{popexo}\newcounter{popexr}
\newtcolorbox{exoresolu}[1][]{popbase,colframe=poporange,colbacklower=popcream,
  segmentation style={popink,line width=1.2pt,dashed},
  before upper={\stepcounter{popexr}\popboxhead{Exercice résolu \thepopexr}{poporange}{#1}},
  before lower={\poppill{Solution}{popink}{popcream}\par\vspace{6pt}}}
% Exercice (non résolu en place) — la correction est dans la partie Corrigés
\newtcolorbox{exercice}[1][]{popbase,colframe=popink,
  before upper={\stepcounter{popexo}\popboxhead{Exercice \thepopexo}{popink}{#1}}}
% Corrigé
\newtcolorbox{corrige}[1][]{popbase,colframe=popgreen,colback=popgreenL,
  before upper={\popboxhead{Corrigé}{popgreen}{#1}}}
% Objectifs / prérequis / bilan
\newtcolorbox{objectifs}{popbase,colframe=popink,colback=poporangeL,
  before upper={\popboxhead{Objectifs de la leçon}{popink}{}}}
\newtcolorbox{prerequis}{popbase,colframe=popink,colback=popblueL,
  before upper={\popboxhead{Prérequis}{popblue}{}}}
\newtcolorbox{bilan}{popbase,colframe=popink,colback=white,boxrule=2.8pt,
  before upper={\popboxhead{Fiche bilan}{poporange}{L'essentiel à connaître pour l'examen}}}
% Figure encadrée avec légende
\newtcolorbox{popfigure}[1][]{enhanced,breakable=false,colback=white,colframe=popline,boxrule=1.4pt,arc=8pt,
  left=10pt,right=10pt,top=10pt,bottom=8pt,before skip=10pt,after skip=12pt,halign=center,
  lower separated=false,halign lower=center,
  fontlower=\popsemi\fontsize{9}{11.5}\selectfont\color{popmuted},
  #1}
\newcommand{\legende}[1]{\par\vspace{6pt}{\popsemi\fontsize{9}{11.5}\selectfont\color{popmuted}#1}}

% ============================================================
% Signature OnBuch+
% ============================================================
\newcommand{\poplogoinline}[1]{{\poplogo\fontsize{#1}{#1}\selectfont\color{popink}OnBuch\color{poporange}+}}
\newcommand{\signatureonbuch}{%
  \begin{tcolorbox}[enhanced,colback=popink,colframe=popink,boxrule=2.2pt,arc=10pt,
      drop shadow=poporange,left=14pt,right=14pt,top=10pt,bottom=10pt,before skip=0pt,after skip=0pt]
    \tikz[baseline=-0.7ex]\node[circle,fill=popgreen,draw=popcream,line width=1.3pt,minimum size=22pt,inner sep=0]{\popcheck{white}};%
    \hspace{9pt}\begin{minipage}[c]{0.62\linewidth}
      {\popxbold\fontsize{12}{13}\selectfont\color{popcream}Signé\ \textbullet\ {\poplogo OnBuch\color{poporange}+}}\par\vspace{2pt}
      {\raggedright\popsemi\fontsize{8}{10.5}\selectfont\color{popline}\MakeUppercase{Équipe pédagogique OnBuch+}\\\MakeUppercase{Conforme au programme officiel MINESEC}\par}
    \end{minipage}\hfill
    {\poplogo\fontsize{22}{22}\selectfont\color{popcream}OnBuch\color{poporange}+}
  \end{tcolorbox}%
}

% ============================================================
% Page de garde
% #1 titre · #2 matière · #3 niveau · #4 module · #5 n° leçon · #6 durée ·
% #7 description / objectifs (texte ou itemize)
% ============================================================
\newcommand{\pagegarde}[7]{%
  \thispagestyle{empty}
  \newgeometry{a4paper,margin=1.7cm,top=1.6cm,bottom=1.4cm}%
  \begin{tikzpicture}[remember picture,overlay]
    \fill[poporange!18] ([xshift=-3.2cm,yshift=-3.6cm]current page.north east) circle (4.6cm);
    \fill[poppurple!14] ([xshift=2.4cm,yshift=7.6cm]current page.south west) circle (3.8cm);
  \end{tikzpicture}%
  {\poplogo\fontsize{30}{30}\selectfont\color{popink}OnBuch\color{poporange}+}\hfill
  \raisebox{0.6ex}{\poppill{Cours signé \textbullet\ Premium}{popink}{popcream}}\par
  \vspace{1pt}\popsquiggle{poppurple}\par
  \vspace{8pt}
  \begin{tcolorbox}[enhanced,colback=poporange,colframe=popink,boxrule=2.8pt,arc=13pt,
      drop shadow=popink,left=18pt,right=18pt,top=12pt,bottom=13pt,after skip=10pt]
    \poppill{#2 \textbullet\ #3}{popink}{popcream}\hspace{5pt}\poppill{#4}{white}{popink}\par\vspace{8pt}
    {\popdisplay\fontsize{14}{14}\selectfont\color{popink}LEÇON #5}\par\vspace{4pt}
    {\raggedright\hyphenpenalty=10000\exhyphenpenalty=10000\popdisplay\fontsize{25}{27}\selectfont\color{popcream}\MakeUppercase{#1}\par}
    \vspace{8pt}
    {\popxbold\fontsize{9.5}{9.5}\selectfont\color{popink}\MakeUppercase{Durée conseillée : #6}}
  \end{tcolorbox}
  \begin{tcolorbox}[enhanced,colback=white,colframe=popink,boxrule=2.2pt,arc=10pt,
      drop shadow=popink,left=16pt,right=16pt,top=10pt,bottom=10pt,after skip=8pt]
    \poppill{Dans cette leçon}{poppurple}{white}\par\vspace{5pt}
    \popsemi\fontsize{9.3}{12.6}\selectfont\color{popink}#7
  \end{tcolorbox}
  \noindent\begin{minipage}[t]{0.487\linewidth}
    \begin{tcolorbox}[enhanced,colback=popgreen,colframe=popink,boxrule=2.2pt,arc=10pt,
        drop shadow=popink,left=11pt,right=11pt,top=10pt,bottom=10pt,height=3.1cm,valign=center]
      \tcbox[colback=white,colframe=popink,boxrule=1.3pt,arc=5pt,boxsep=0pt,nobeforeafter,
        left=3pt,right=3pt,top=3pt,bottom=3pt,valign=center]{\qrcode[height=1.9cm]{https://play.google.com/store/apps/details?id=com.luvvix.onbuch&hl=fr}}%
      \hspace{8pt}\begin{minipage}[c]{0.5\linewidth}
        {\popdisplay\fontsize{11}{12.5}\selectfont\color{white}\MakeUppercase{Télécharge OnBuch}}\par\vspace{4pt}
        {\raggedright\popsemi\fontsize{8.4}{11}\selectfont\color{white}Gratuit \textbullet\ Google Play\\Tuteur IA, annales, concours, résultats\par}
      \end{minipage}
    \end{tcolorbox}
  \end{minipage}\hfill
  \begin{minipage}[t]{0.487\linewidth}
    \begin{tcolorbox}[enhanced,colback=poppurple,colframe=popink,boxrule=2.2pt,arc=10pt,
        drop shadow=popink,left=11pt,right=11pt,top=10pt,bottom=10pt,height=3.1cm,valign=center]
      \tikz[baseline=-0.7ex]\node[circle,fill=white,draw=popink,line width=1.3pt,minimum size=1.6cm,inner sep=0]{\popdisplay\fontsize{24}{24}\selectfont\color{poppurple}+};%
      \hspace{8pt}\begin{minipage}[c]{0.6\linewidth}
        {\popdisplay\fontsize{11}{12.5}\selectfont\color{white}\MakeUppercase{Passe à OnBuch+}}\par\vspace{4pt}
        {\raggedright\popsemi\fontsize{8.4}{11}\selectfont\color{white}Tous les cours signés, Tuteur IA illimité, fiches PDF — onglet OnBuch+ de l'app\par}
      \end{minipage}
    \end{tcolorbox}
  \end{minipage}
  \vspace{8pt}
  \popcontenu
  \vfill
  \signatureonbuch
  \restoregeometry
  \newpage
}

% Rangée « Ce cours contient » (page de garde)
\newcommand{\poptile}[3]{% #1 couleur #2 gros texte #3 légende
  \begin{tcolorbox}[enhanced,colback=#1,colframe=popink,boxrule=2pt,arc=9pt,shadow={2.5pt}{-2.5pt}{0pt}{popink},
      left=6pt,right=6pt,top=9pt,bottom=9pt,halign=center,valign=center,height=1.75cm,width=\linewidth,nobeforeafter]
    {\popdisplay\fontsize{12.5}{13}\selectfont\color{white}\MakeUppercase{#2}}\par\vspace{3pt}
    {\centering\popsemi\fontsize{7.8}{9.5}\selectfont\color{white}#3\par}
  \end{tcolorbox}}
\newcommand{\popcontenu}{%
  \par\noindent{\popxboldls\fontsize{7.5}{7.5}\selectfont\color{popmuted}CE COURS SIGNÉ CONTIENT}\par\vspace{7pt}
  \noindent\begin{minipage}[t]{0.235\linewidth}\poptile{popblue}{Cours}{complet, illustré, pas à pas}\end{minipage}\hfill
  \begin{minipage}[t]{0.235\linewidth}\poptile{poporange}{Méthodes}{et exercices résolus}\end{minipage}\hfill
  \begin{minipage}[t]{0.235\linewidth}\poptile{poppink}{Exercices}{type examen + corrigés}\end{minipage}\hfill
  \begin{minipage}[t]{0.235\linewidth}\poptile{popgreen}{Fiche bilan}{l'essentiel à retenir}\end{minipage}\par}

% Sommaire
\newtcolorbox{sommaire}{enhanced,colback=white,colframe=popink,boxrule=2.2pt,arc=10pt,
  drop shadow=popink,left=18pt,right=18pt,top=13pt,bottom=13pt,before skip=6pt,after skip=20pt,
  before upper={\poppill{Sommaire}{poppurple}{white}\par\vspace{9pt}\popsemi\fontsize{10.6}{14.5}\selectfont\color{popink}}}

% Signature de fin de cours — poussée tout en bas de la dernière page
\newcommand{\findecours}{%
  \par\vfill\nopagebreak
  \begin{center}\popsquiggle{poporange}\end{center}\vspace{4pt}
  \signatureonbuch
}
\AtBeginDocument{\def\llbracket{\mathopen{[\mkern-3.5mu[}}\def\rrbracket{\mathclose{]\mkern-3.5mu]}}} % intervalles d'entiers (glyphe ⟦ absent de la police)
% Glyphes absents de Fira Math : redéfinis à partir de glyphes disponibles
\AtBeginDocument{%
  \def\setminus{\mathbin{\backslash}}\let\smallsetminus\setminus
  \def\vdots{\mathord{\vbox{\baselineskip3.2pt\lineskiplimit0pt\kern4pt\hbox{.}\hbox{.}\hbox{.}}}}%
  \def\ddots{\mathinner{\mkern1mu\raise6.5pt\hbox{.}\mkern2mu\raise3.6pt\hbox{.}\mkern2mu\raise0.7pt\hbox{.}\mkern1mu}}%
  \def\triangle{\mathord{\scalebox{0.9}{$\symup{\Delta}$}}}%
  \def\nabla{\mathord{\raisebox{\depth}{\scalebox{1}[-1]{$\symup{\Delta}$}}}}%
  \def\checkmark{\popcheck{popdark}}%
  \def\star{\mathbin{\ast}}\def\bigstar{\mathord{\ast}}%
  \def\diamond{\mathbin{\scalebox{0.6}{$\Diamond$}}}%
  \def\bigcup{\mathop{\mathchoice{\raisebox{-0.3em}{\scalebox{1.7}{$\cup$}}}{\scalebox{1.25}{$\cup$}}{\cup}{\cup}}\displaylimits}%
  \def\bigcap{\mathop{\mathchoice{\raisebox{-0.3em}{\scalebox{1.7}{$\cap$}}}{\scalebox{1.25}{$\cap$}}{\cap}{\cap}}\displaylimits}%
  \def\lesseqqgtr{\mathrel{\lessgtr}}\def\bigoplus{\mathop{\oplus}\displaylimits}%
}
\providecommand{\ssi}{si et seulement si} % abréviation fréquente
\AtBeginDocument{\providecommand{\wideparen}{\overparen}} % arc de cercle

\begin{document}

\pagegarde{Leçon 3 — Expression écrite : situation de la femme}{Chinois}{Tle A}{Module 1 : Vie familiale et intégration sociale}{ZH03}{2 h}{Leçon d'expression écrite consacrée à la situation de la femme : étude d'un texte modèle commenté, maîtrise des connecteurs et des caractères du thème (traits, radicaux, caractères proches), puis entraînement au thème et à la version. L'élève produit et traduit des textes simples sur la place de la femme dans la famille et la société.
\vspace{4pt}
\begin{multicols}{2}\small
\begin{itemize}
\item À la fin de cette leçon, je sais utiliser le vocabulaire essentiel sur la situation de la femme (femme, travail, famille, égalité, droits)
\item À la fin de cette leçon, je sais écrire correctement les caractères du thème en respectant les radicaux et l'ordre des traits
\item À la fin de cette leçon, je sais distinguer les caractères proches faciles à confondre (她/他, 妇/扫, 好/如)
\item À la fin de cette leçon, je sais analyser la structure d'un texte modèle simple (introduction, développement, conclusion)
\item À la fin de cette leçon, je sais utiliser les connecteurs utiles : 因为, 所以, 但是, 而且, 首先, 最后
\item À la fin de cette leçon, je sais produire un texte simple de 5 à 8 phrases en chinois sur la situation de la femme
\end{itemize}\end{multicols}}

\begin{sommaire}
\begin{multicols}{2}
\begin{enumerate}
\item Texte modèle : la femme d'aujourd'hui
\item Lexique du thème et glossaire
\item Écriture des caractères : traits, radicaux, caractères proches
\item Structure du texte et connecteurs
\item Thème et version guidés
\item Production écrite et grille d'évaluation
\item Activité d'intégration
\item Méthodes et exercices résolus
\item Exercices
\item Corrigés des exercices
\item Fiche bilan
\end{enumerate}
\end{multicols}
\end{sommaire}

\begin{objectifs}
\begin{itemize}
\item À la fin de cette leçon, je sais utiliser le vocabulaire essentiel sur la situation de la femme (femme, travail, famille, égalité, droits)
\item À la fin de cette leçon, je sais écrire correctement les caractères du thème en respectant les radicaux et l'ordre des traits
\item À la fin de cette leçon, je sais distinguer les caractères proches faciles à confondre (她/他, 妇/扫, 好/如)
\item À la fin de cette leçon, je sais analyser la structure d'un texte modèle simple (introduction, développement, conclusion)
\item À la fin de cette leçon, je sais utiliser les connecteurs utiles : 因为, 所以, 但是, 而且, 首先, 最后
\item À la fin de cette leçon, je sais produire un texte simple de 5 à 8 phrases en chinois sur la situation de la femme
\item À la fin de cette leçon, je sais traduire des phrases simples du français vers le chinois (thème)
\item À la fin de cette leçon, je sais traduire un court texte du chinois vers le français (version)
\item À la fin de cette leçon, je sais m'auto-évaluer à l'aide d'une grille d'évaluation inspirée du Bac
\end{itemize}
\end{objectifs}

\begin{prerequis}
\begin{itemize}
\item Vocabulaire de la famille et de la vie quotidienne (Première) : 家, 爸爸, 妈妈, 工作, 孩子
\item Structure de base de la phrase chinoise Sujet-Verbe-Complément
\item Emploi de 是, 有, 在 et des adjectifs attributs avec 很
\item Notions sur les radicaux et l'ordre général des traits (haut-bas, gauche-droite)
\item Connecteurs simples déjà rencontrés : 和, 也, 都
\end{itemize}
\end{prerequis}

\newpage

\coursec{Texte modèle : la femme d'aujourd'hui}

Au Cameroun comme en Chine, le rôle de la femme dans la famille et la société a profondément changé : de plus en plus de femmes sont chefs d'entreprise à Douala, ministres à Yaoundé ou ingénieures à Pékin. Pourtant, beaucoup restent aussi responsables des tâches domestiques. Comment raconter cette évolution en chinois ? Dans cette leçon, vous apprendrez à écrire un texte simple sur la situation de la femme, à traduire des phrases du français vers le chinois (thème) et du chinois vers le français (version), et à écrire correctement les caractères clés du thème, comme \zh{女} (femme), \zh{她} (elle) et \zh{务} (affaire, tâche).

\courssub{Lecture et compréhension du texte modèle}

Lisez attentivement le texte ci-dessous. Il est composé de phrases courtes, au vocabulaire HSK 2-3, structuré en trois temps : le passé, le présent, l'opinion. Observez les marqueurs temporels et les connecteurs logiques. Le pinyin ci-dessous reflète la \textbf{prononciation réelle} (avec changements de tons / sandhi), pas seulement la forme de citation.

\begin{textebox}{现在的女人  —  La femme d'aujourd'hui}
以前，很多女人不工作，只在家里做饭、看孩子。\\
现在的女人不一样了：她们可以上学，可以工作，也可以当医生、老师或者经理。\\
但是，很多女人在家里还要做很多家务。\\
我觉得男人和女人应该一样，家务也应该一起做。
\tcblower
\textbf{Pinyin (prononciation réelle) :}\\
Yǐqián, hěn duō nǚrén bù gōngzuò, zhǐ zài jiālǐ zuòfàn, kàn háizi.\\
Xiànzài de nǚrén bú yíyàng le: tāmen kěyǐ shàngxué, kěyǐ gōngzuò, yě kěyǐ dāng yīshēng, lǎoshī huòzhě jīnglǐ.\\
Dànshì, hěn duō nǚrén zài jiālǐ hái yào zuò hěn duō jiāwù.\\
Wǒ juéde nánrén hé nǚrén yīnggāi yíyàng, jiāwù yě yīnggāi yíqǐ zuò.

\textbf{Traduction :}\\
Autrefois, beaucoup de femmes ne travaillaient pas, elles restaient à la maison pour cuisiner et garder les enfants.\\
Les femmes d'aujourd'hui sont différentes : elles peuvent étudier, travailler, et aussi devenir médecins, enseignantes ou directrices.\\
Mais, beaucoup de femmes doivent encore faire beaucoup de tâches ménagères à la maison.\\
Je pense que l'homme et la femme devraient être égaux, et que les tâches ménagères devraient aussi être faites ensemble.
\end{textebox}

\courssub{Analyse de la structure du texte : trois temps, une progression}

Le texte suit une structure classique de l'argumentation simple en chinois : \textbf{Contexte passé} $\rightarrow$ \textbf{Constat présent} $\rightarrow$ \textbf{Avis personnel}. Repérez les trois blocs introduits par les marqueurs temporels \zh{以前} (autrefois), \zh{现在} (maintenant) et le verbe d'opinion \zh{觉得} (penser / trouver).

\begin{popfigure}
\begin{tikzpicture}[
    node distance=1.2cm and 2.5cm,
    block/.style={rectangle, rounded corners, draw=popblue, thick, fill=popblueL, text width=3.8cm, align=center, minimum height=2.2cm, font=\small},
    arrow/.style={-Stealth, thick, popdark},
    labelnode/.style={font=\footnotesize\bfseries, text=popblue, anchor=south}
]
\node[block, align=center] (past){
    \textbf{1. PASSÉ (\zh{以前})}\\
    \zh{很多女人不工作}\\
    (Beaucoup de femmes ne travaillent pas)\\
    \zh{只在家里做饭、看孩子}\\
    (Restent seulement à la maison cuisiner/garder enfants)
};
\node[block, right=of past, align=center] (present){
    \textbf{2. PRÉSENT (\zh{现在})}\\
    \zh{女人不一样了}\\
    (Femmes différentes)\\
    \zh{可以上学、工作、当医生/老师/经理}\\
    (Peuvent étudier, travailler, être médecin/prof/manager)\\
    \zh{也} (aussi) marque l'addition
};
\node[block, right=of present, align=center] (opinion){
    \textbf{3. OPINION (\zh{我觉得})}\\
    \zh{但是} (mais) : contraste\\
    \zh{还要做很多家务}\\
    (Doivent encore faire ménage)\\
    \zh{应该一样、一起做}\\
    (Devraient être égaux, faire ensemble)
};
\draw[arrow] (past.east) -- node[above, labelnode] {Évolution} (present.west);
\draw[arrow] (present.east) -- node[above, labelnode] {Nuance / Réalité} (opinion.west);
\node[labelnode] at (past.north) {\zh{以前} \textit{Yǐqián}};
\node[labelnode] at (present.north) {\zh{现在} \textit{Xiànzài}};
\node[labelnode] at (opinion.north) {\zh{我觉得} \textit{Wǒ juéde}};
\end{tikzpicture}
\legende{Organigramme de la structure du texte modèle : Passé $\rightarrow$ Présent $\rightarrow$ Opinion}
\end{popfigure}

\begin{methode}{Analyser un texte modèle en chinois (niveau Terminale)}
Appliquez ces trois étapes systématiquement à tout texte de compréhension ou modèle d'expression écrite :
\begin{enumerate}
    \item \textbf{Identifiez les marqueurs temporels et logiques} : soulignez \zh{以前}、\zh{现在}、\zh{以后}、\zh{但是}、\zh{因为}...\zh{所以}... Ils découpent le texte en paragraphes logiques.
    \item \textbf{Soulignez les connecteurs et mots-outils} : \zh{也}、\zh{还}、\zh{只}、\zh{可以}、\zh{应该}、\zh{了} (aspect accompli / changement d'état). Notez leur position dans la phrase (souvent avant le verbe).
    \item \textbf{Repérez la phrase d'opinion finale} : elle contient souvent \zh{我觉得}、\zh{我认为}、\zh{应该}、\zh{建议}. C'est le modèle à réutiliser pour votre propre production écrite.
\end{enumerate}
\end{methode}

\courssub{Lexique clé du thème : tableaux de vocabulaire}

Maîtrisez ces mots (écriture, pinyin de citation, sens, nature). Ils constituent le socle lexical du thème « Situation de la femme ». \textbf{Note :} Les tons indiqués sont les tons de citation (dictionnaire) ; la prononciation en phrase suit les règles de sandhi (voir exercice de lecture).

\begin{center}
\begin{tabularx}{\linewidth}{|X|X|X|X|}
\hline
\rowcolor{popblueL} \textbf{Caractères} & \textbf{Pinyin (citation)} & \textbf{Nature} & \textbf{Traduction} \\ \hline
\zh{以前} & yǐqián & nom / adv. temps & autrefois, avant \\ \hline
\zh{现在} & xiànzài & nom / adv. temps & maintenant, actuellement \\ \hline
\zh{女人} & nǚrén & nom & femme (adulte) \\ \hline
\zh{男人} & nánrén & nom & homme (adulte) \\ \hline
\zh{工作} & gōngzuò & verbe / nom & travailler / travail, emploi \\ \hline
\zh{只} & zhǐ & adv. & seulement, ne ... que \\ \hline
\zh{在家里} & zài jiālǐ & locution & à la maison, chez soi \\ \hline
\zh{做饭} & zuòfàn & verbe & cuisiner, faire la cuisine \\ \hline
\zh{看孩子} & kàn háizi & verbe + obj. & garder les enfants, s'occuper des enfants \\ \hline
\zh{不一样} & bù yíyàng & adj. & différent, pas pareil \\ \hline
\zh{了} & le & particule & marque le changement d'état / aspect accompli \\ \hline
\zh{可以} & kěyǐ & verbe modal & pouvoir (permission, possibilité) \\ \hline
\zh{上学} & shàngxué & verbe & aller à l'école, étudier \\ \hline
\zh{当} & dāng & verbe & être (fonction), devenir, exercer le métier de \\ \hline
\zh{医生} & yīshēng & nom & médecin, docteur \\ \hline
\zh{老师} & lǎoshī & nom & professeur, enseignant \\ \hline
\zh{经理} & jīnglǐ & nom & directeur, manager, gérant \\ \hline
\zh{或者} & huòzhě & conjonction & ou bien, ou (alternative) \\ \hline
\zh{但是} & dànshì & conjonction & mais, cependant \\ \hline
\zh{还要} & hái yào & adv. + verbe & devoir encore, avoir encore à \\ \hline
\zh{家务} & jiāwù & nom & tâches ménagères, travail domestique \\ \hline
\zh{觉得} & juéde & verbe & penser, trouver, avoir l'impression \\ \hline
\zh{应该} & yīnggāi & verbe modal & devoir (obligation morale, conseil) \\ \hline
\zh{一样} & yíyàng & adj. & pareil, identique, égal \\ \hline
\zh{一起} & yìqǐ & adv. & ensemble \\ \hline
\end{tabularx}
\end{center}

\courssub{Écriture des caractères : traits, radicaux, caractères proches}

L'écriture correcte (ordre des traits, proportion, radicaux) est évaluée à l'examen. Travaillez ces caractères fréquents du thème.

\begin{propriete}{Rappel des règles d'ordre des traits (révision)}
\begin{enumerate}
    \item De haut en bas (\zh{三}).
    \item De gauche à droite (\zh{川}).
    \item Horizontal avant vertical (\zh{十}).
    \item Diagonale droite avant diagonale gauche (\zh{人}).
    \item Le centre avant les ailes (\zh{小}).
    \item Fermer en dernier (\zh{口} \rightarrow \zh{日}).
    \item Le trait transversal final (\zh{王} \leftrightarrow \zh{主}).
\end{enumerate}
\end{propriete}

\begin{exemplebox}{Fiches caractères du thème (à recopier 5 fois chacun)}
\begin{center}
\begin{tabularx}{\linewidth}{|c|X|X|X|}
\hline
\rowcolor{popblueL} \textbf{Car.} & \textbf{Ordre des traits (description)} & \textbf{Radical (clé)} & \textbf{Caractères proches / Pièges} \\ \hline
\zh{女} & 1. Trait courbe gauche (haut→bas) 2. Trait courbe droit (haut→bas) 3. Trait horizontal (gauche→droite) & \zh{女} (nǚ, femme) & \zh{妈} (mère), \zh{好} (bon), \zh{她} (elle), \zh{姐} (sœur aînée), \zh{妹} (sœur cadette) \\ \hline
\zh{她} & \zh{女} (3 traits) + \zh{也} (3 traits : horizontale, verticale, horizontale) = 6 traits & \zh{女} & \zh{他} (lui, radical \zh{亻}), \zh{它} (ça, radical \zh{宀}), \zh{牠} (animal, radical \zh{牛}) \\ \hline
\zh{务} & 1. Horizontal 2. Vertical 3. Horizontal 4. Crochet vertical 5. Point (force \zh{力}) = 5 traits & \zh{力} (lì, force) & \zh{务} (simplifié) vs \zh{務} (traditionnel). Ne pas confondre avec \zh{事} (shì, affaire, radical \zh{亅}) \\ \hline
\zh{当} & 1. Horizontal 2. Vertical 3. Horizontal 4. Vertical 5. Horizontal 6. Horizontal 7. Crochet vertical = 7 traits & \zh{彐} (jí, tête de cochon) ou \zh{小} (petit) selon dictionnaires & \zh{尚} (shàng, encore), \zh{党} (dǎng, parti) \\ \hline
\zh{医} & \zh{殳} (shū, arme) + \zh{酉} (yǒu, vase/vin) = 7 traits (simplifié) & \zh{酉} (yǒu) & \zh{药} (yào, médicament, radical \zh{艹}), \zh{毉} (traditionnel) \\ \hline
\zh{家} & \zh{宀} (mián, toit) + \zh{豕} (shǐ, porc) = 10 traits & \zh{宀} (mián, toit) & \zh{贾} (gǔ, marchand), \zh{豪} (háo, héroïque) \\ \hline
\end{tabularx}
\end{center}
\end{exemplebox}

\begin{attention}{Pièges d'écriture fréquents}
\begin{itemize}
    \item \zh{女} : Le 3ème trait (horizontal) \textbf{ne traverse pas} les deux premiers traits. Il s'arrête à droite du second trait.
    \item \zh{她} / \zh{他} / \zh{它} : Même phonétique \zh{也} (simplifié) ou \zh{也} (traditionnel pour 他/她), mais radicaux différents : \zh{女} (femme), \zh{亻} (homme/humain), \zh{宀} (chose/animal/toit). \textbf{Ne les confondez pas à l'écrit.}
    \item \zh{务} : Version simplifiée très différente du traditionnel \zh{務}. Apprenez la forme simplifiée \zh{务} (radical \zh{力}).
    \item \zh{当} : Le bas est \zh{彐} (deux horizontales, une verticale, un crochet), pas \zh{尤} (yóu) ni \zh{尢} (wāng).
\end{itemize}
\end{attention}

\courssub{Focus sur les connecteurs logiques et mots-outils du texte}

Le texte modèle utilise trois connecteurs essentiels pour articuler les idées. Observez leur place syntaxique : \textbf{en position préverbale (avant le verbe ou l'adjectif)}.

\begin{propriete}{Position des connecteurs \zh{也}、\zh{还}、\zh{但是} et des modaux \zh{可以}、\zh{应该}}
En chinois, les adverbes de portée (aussi, encore), les conjonctions de coordination (mais) et les verbes modaux (pouvoir, devoir) se placent \textbf{avant le verbe principal} (ou l'adjectif s'il n'y a pas de verbe).
\medskip
\textbf{Schéma :} Sujet + (Temps) + \textbf{Connecteur / Modal} + Verbe / Adjectif + Complément.
\end{propriete}

\begin{exemplebox}{Analyse des connecteurs du texte}
\begin{itemize}
    \item \zh{也可以当医生} \textit{(yě kěyǐ dāng yīshēng)} : \zh{也} (aussi) + \zh{可以} (pouvoir) + \zh{当} (être). \textbf{Ordre :} aussi > pouvoir > verbe.
    \item \zh{还要做很多家务} \textit{(hái yào zuò hěn duō jiāwù)} : \zh{还} (encore) + \zh{要} (devoir/falloir) + \zh{做} (faire). \textbf{Ordre :} encore > devoir > verbe.
    \item \zh{应该一样} \textit{(yīnggāi yíyàng)} : \zh{应该} (devoir) + \zh{一样} (adj. pareil). \textbf{Ordre :} devoir > adjectif.
    \item \zh{家务也应该一起做} \textit{(jiāwù yě yīnggāi yíqǐ zuò)} : Sujet \zh{家务} + \zh{也} (aussi) + \zh{应该} (devoir) + \zh{一起} (ensemble) + \zh{做} (faire). \textit{Note :} \zh{一起} se place après le modal et avant le verbe.
\end{itemize}
\end{exemplebox}

\begin{attention}{Pièges fréquents pour francophones : \zh{也} vs \zh{还}, \zh{只}, \zh{但是}}
\begin{itemize}
    \item \textbf{\zh{也} (yě) vs \zh{还} (hái)} : \zh{也} = « aussi / également » (addition simple, même sujet ou sujet différent). \zh{还} = « en plus / encore / davantage » (souvent même sujet, idée de progression ou quantité). \textit{Exemple :} \zh{我喜欢喝茶，也喜欢喝咖啡} (J'aime le thé, j'aime \textbf{aussi} le café) vs \zh{我喝了茶，还想喝咖啡} (J'ai bu du thé, je veux \textbf{encore/en plus} du café).
    \item \textbf{Place de \zh{只} (zhǐ) « seulement »} : Il se place \textbf{immédiatement avant le verbe} (ou l'élément qu'il limite). \zh{只在家里} (seulement à la maison) ; \zh{只做饭} (ne faire que cuisiner). \textbf{Erreur :} *\zh{在只家里} ou *\zh{做只饭}.
    \item \textbf{\zh{但是} (dànshì) en tête de phrase} : Contrairement au français « Mais, ... » qui peut être suivi d'une virgule et du sujet, en chinois \zh{但是} introduit souvent une nouvelle phrase complète. \textit{Correct :} \zh{但是，很多女人...} \textbf{Erreur :} *\zh{但是很多女人...} (sans pause, style oral uniquement).
    \item \textbf{\zh{一} (yī) et \zh{不} (bù) : Sandhi tonaux obligatoires} : \zh{不} $\rightarrow$ \zh{bu} (2ème ton) devant un 4ème ton (\zh{不一样} \rightarrow \textit{bú yíyàng}). \zh{一} $\rightarrow$ \zh{yi} (2ème ton) devant un 4ème ton (\zh{一样} \rightarrow \textit{yíyàng}) ; \zh{yi} (4ème ton) devant 1, 2, 3ème tons (\zh{一起} \rightarrow \textit{yìqǐ} dictionnaire, mais \textit{yíqǐ} à l'oral devant ton 3).
\end{itemize}
\end{attention}

\courssub{Point de grammaire : la structure \zh{当} + métier}

Pour exprimer « être / devenir [métier / fonction] », on utilise le verbe \zh{当} (dāng). C'est la structure standard à ce niveau (HSK 3), plus naturelle que \zh{是} (shì) pour une fonction professionnelle.

\begin{propriete}{Structure : Sujet + \zh{当} + Métier / Fonction}
\zh{当} (dāng) signifie « exercer la fonction de », « être (en tant que) ». Il s'utilise sans article (pas de « un/une »).
\medskip
\textbf{Exemples :} \zh{当医生} (être médecin), \zh{当老师} (être prof), \zh{当经理} (être manager), \zh{当工程师} (être ingénieur), \zh{当部长} (être ministre).
\end{propriete}

\begin{exemplebox}{Variantes avec \zh{当} et modaux}
\begin{tabular}{lll}
\zh{她可以当医生。} & Tā kěyǐ dāng yīshēng. & Elle peut / a le droit d'être médecin. \\
\zh{她想当老师。} & Tā xiǎng dāng lǎoshī. & Elle veut être professeure. \\
\zh{我姐姐当经理。} & Wǒ jiějie dāng jīnglǐ. & Ma sœur aînée est manager. \\
\zh{将来我想当工程师。} & Jiānglái wǒ xiǎng dāng gōngchéngshī. & Plus tard, je veux être ingénieur. \\
\zh{她当了三年医生了。} & Tā dāng le sān nián yīshēng le. & Elle est médecin depuis 3 ans (durée).
\end{tabular}
\end{exemplebox}

\courssub{Le saviez-vous ? : La fête des femmes en Chine (\zh{三八妇女节})}

\begin{savaistu}{Le 8 mars en Chine : \zh{三八妇女节} (Sānbā Fùnǚjié)}
En Chine, la Journée internationale des droits des femmes (\zh{三八国际妇女节}) est un jour férié \textbf{spécifique aux femmes}. Selon la loi, les femmes bénéficient d'une \textbf{demi-journée de congé} (souvent l'après-midi). C'est l'occasion pour les entreprises d'offrir des petits cadeaux (fleurs, produits de beauté) ou d'organiser des activités. Au Cameroun, le 8 mars est aussi très célébré (défilés, tenues traditionnelles \zh{印花布} / \zh{卡巴} \textit{kaba}), mais c'est un jour férié pour tous. Notez le vocabulaire : \zh{妇女} (femme mariée / femme adulte, terme formel, administratif) vs \zh{女人} (femme, terme neutre/courant) vs \zh{女生} (jeune fille, étudiante) vs \zh{女士} (dame, terme poli).
\end{savaistu}

\courssub{Exercice résolu : Thème guidé (Français $\rightarrow$ Chinois)}

Appliquons la méthode : identifiez le temps, le sujet, le modal/connecteur, le verbe.

\begin{exoresolu}{Traduire en chinois (caractères + pinyin)}
\textbf{Énoncé :} « Autrefois, ma mère ne travaillait pas dehors. Elle restait à la maison pour faire le ménage. Maintenant, elle est médecin. Mais elle doit encore cuisiner le soir. Je pense qu'elle est très fatiguée. »

\textbf{Solution détaillée :}
\begin{enumerate}
    \item \textit{« Autrefois, ma mère ne travaillait pas dehors. »} \\
    Marqueur temps : \zh{以前}. Sujet : \zh{我妈妈} (ma mère). Négation état habituel passé : \zh{不}. Verbe : \zh{工作}. Complément lieu : \zh{在外面} (dehors / à l'extérieur). \\
    $\rightarrow$ \zh{以前，我妈妈不在外面工作。} \textit{(Yǐqián, wǒ māma bù zài wàimiàn gōngzuò.)}
    \item \textit{« Elle restait à la maison pour faire le ménage. »} \\
    Sujet omis (elle). Structure lieu + action. « Faire le ménage » = \zh{做家务} ou \zh{打扫卫生}. \zh{只} (seulement) marque la restriction. \\
    $\rightarrow$ \zh{她只在家里做家务。} \textit{(Tā zhǐ zài jiālǐ zuò jiāwù.)}
    \item \textit{« Maintenant, elle est médecin. »} \\
    Marqueur : \zh{现在}. Structure \zh{当} + métier. Changement d'état / nouvelle situation : \zh{了} en fin de phrase. \\
    $\rightarrow$ \zh{现在她当医生了。} \textit{(Xiànzài tā dāng yīshēng le.)}
    \item \textit{« Mais elle doit encore cuisiner le soir. »} \\
    Connecteur : \zh{但是}. Sujet : \zh{她}. Temps : \zh{晚上} (placé avant le modal). « Doit encore » : \zh{还要}. Verbe : \zh{做饭}. \\
    $\rightarrow$ \zh{但是，她晚上还要做饭。} \textit{(Dànshì, tā wǎnshang hái yào zuòfàn.)}
    \item \textit{« Je pense qu'elle est très fatiguée. »} \\
    Opinion : \zh{我觉得}. Complément d'objet : proposition \zh{她很累}. \zh{很} + adj. (pas de \zh{是}). \\
    $\rightarrow$ \zh{我觉得她很累。} \textit{(Wǒ juéde tā hěn lèi.)}
\end{enumerate}
\textbf{Texte complet :} \zh{以前，我妈妈不在外面工作。她只在家里做家务。现在她当医生了。但是，她晚上还要做饭。我觉得她很累。}
\end{exoresolu}

\courssub{Exercice résolu : Version guidée (Chinois $\rightarrow$ Français)}

\begin{exoresolu}{Traduire en français}
\textbf{Énoncé :} \zh{我奶奶以前不识字，不上学。现在她学会了用微信，也会看新闻。但是她觉得眼睛不太好。建议她少看手机。}

\textbf{Analyse et solution :}
\begin{enumerate}
    \item \zh{我奶奶以前不识字，不上学。} \textit{(Wǒ nǎinai yǐqián bù shízì, bù shàngxué.)} \\
    \zh{奶奶} (grand-mère maternelle), \zh{识字} (savoir lire/écrire, litt. connaître les caractères), \zh{不} double négation = « ne savait ni lire ni écrire, n'allait pas à l'école ».
    \item \zh{现在她学会了用微信，也会看新闻。} \textit{(Xiànzài tā xuéhuì le yòng Wēixìn, yě huì kàn xīnwén.)} \\
    \zh{学会} (apprendre à + résultat acquis), \zh{用} (utiliser), \zh{微信} (WeChat), \zh{也会} (sait aussi), \zh{看新闻} (regarder les infos/lire l'actualité).
    \item \zh{但是她觉得眼睛不太好。} \textit{(Dànshì tā juéde yǎnjing bù tài hǎo.)} \\
    \zh{眼睛} (yeux), \zh{不太好} (pas très bien / mauvais).
    \item \zh{建议她少看手机。} \textit{(Jiànyì tā shǎo kàn shǒujī.)} \\
    \zh{建议} (suggérer/conseiller), \zh{少} (peu/moins), \zh{手机} (téléphone portable).
\end{enumerate}
\textbf{Traduction :} « Ma grand-mère autrefois ne savait ni lire ni écrire, n'allait pas à l'école. Maintenant elle a appris à utiliser WeChat, elle sait aussi lire les actualités. Mais elle trouve que ses yeux ne vont pas trop bien. Je lui conseille de moins regarder son téléphone. »
\end{exoresolu}

\courssub{Entraînement à la lecture à voix haute}

\begin{exercice}{Lecture fluide et expressive}
Lisez le texte modèle \zh{现在的女人} à voix haute en respectant :
\begin{enumerate}
    \item Les \textbf{tons} (notamment les tons 3 \zh{女人 nǚrén}, \zh{可以 kěyǐ}, \zh{经理 jīnglǐ} et le ton neutre \zh{了 le}, \zh{一起 yìqǐ}).
    \item Les \textbf{pauses} (marquées par les virgules chinoises \zh{、} \zh{，} et le point \zh{。}) : groupe \zh{以前} / \zh{很多女人不工作} / \zh{只在家里做饭、看孩子}。
    \item L'\textbf{intonation} : ton narratif pour les deux premières phrases, ton contrastif sur \zh{但是}, ton affirmatif sur \zh{我觉得}.
\end{enumerate}
\textit{En classe : lecture collective chorale, puis lecture individuelle par volontaires. L'enseignant corrige les tons 3 enchaînés (règle du sandhi : 3+3 $\rightarrow$ 2+3, ex: \zh{可以} \textit{kéyǐ}).}
\end{exercice}

\begin{corrige}{Exercice : Lecture fluide et expressive}
\textbf{Points de vigilance corrigés par l'enseignant :}
\begin{itemize}
    \item \zh{女人} \textit{nǚrén} : 3+2 (sandhi sur premier 3) $\rightarrow$ \textit{núren}.
    \item \zh{可以} \textit{kěyǐ} : 3+3 $\rightarrow$ \textit{kéyǐ}.
    \item \zh{经理} \textit{jīnglǐ} : 1+3 (attention au ton 1 haut et plat).
    \item \zh{不一样} \textit{bù yíyàng} : \zh{不} seul devant ton 4 $\rightarrow$ ton 2 \textit{bú}. \zh{一样} \textit{yíyàng} : ton 1 (sandhi sur \zh{一} devant ton 4) + ton 4.
    \item \zh{一起做} \textit{yìqǐ zuò} : \zh{一} devant ton 3 $\rightarrow$ ton 2 \textit{yíqǐ} (souvent prononcé ton 4 \textit{yìqǐ} en débit rapide, accepter les deux).
    \item Liaison \zh{觉得} \textit{juéde} (ton 2 + ton neutre) sans pause avant la proposition.
    \item \zh{家务} \textit{jiāwù} : 1+4. \zh{医生} \textit{yīshēng} : 1+1.
\end{itemize}
\end{corrige}

\courssub{Production écrite guidée : « Portrait d'une femme d'aujourd'hui »}

\begin{methode}{Rédiger un paragraphe structuré (60-80 caractères chinois)}
Suivez le plan en 3 étapes du texte modèle. Utilisez la grille d'auto-évaluation ci-dessous avant de rendre votre copie.
\begin{enumerate}
    \item \textbf{Étape 1 - Le passé (\zh{以前})} : Présentez la situation d'une femme (mère, grand-mère, voisine, figure publique) dans le passé. Utilisez \zh{不工作}、\zh{只在家里}、\zh{做饭}、\zh{看孩子} ou \zh{不识字}...
    \item \textbf{Étape 2 - Le présent (\zh{现在}... \zh{了})} : Décrivez le changement : études, travail, métier (\zh{当} + métier), nouvelles compétences (\zh{学会}...). Marquez le changement avec \zh{了}.
    \item \textbf{Étape 3 - La réalité / L'opinion (\zh{但是} / \zh{我觉得} / \zh{应该})} : Nuancez (tâches ménagères \zh{还要做家务}, fatigue \zh{很累}) et donnez votre avis (\zh{应该一样}、\zh{一起做}、\zh{建议}...).
\end{enumerate}
\end{methode}

\begin{exercice}{Sujet de rédaction (Devoir maison ou composition en classe)}
\textbf{Consigne :} Écrivez un texte de \textbf{60 à 80 caractères chinois} (ponctuation comprise) sur le thème : \zh{我认识的一位女性} (Une femme que je connais).
\textbf{Contraintes obligatoires :}
\begin{itemize}
    \item Utiliser la structure en 3 temps : \zh{以前} -- \zh{现在} -- \zh{我觉得/建议}.
    \item Employer au minimum \textbf{5 mots du lexique} du tableau (ex: \zh{工作}、\zh{当}、\zh{家务}、\zh{可以}、\zh{应该}).
    \item Utiliser correctement les connecteurs : \zh{也}、\zh{还}、\zh{但是}、\zh{只}.
    \item Écrire les caractères soigneusement (ordre des traits, proportions).
    \item Ajouter le pinyin au-dessus de chaque caractère (facultatif mais recommandé pour la révision).
\end{itemize}
\textbf{Amorce suggérée :} \zh{以前，我的... 现在... 但是... 我觉得...}
\end{exercice}

\begin{center}
\begin{tabularx}{\linewidth}{|X|c|c|c|c|}
\hline
\rowcolor{popblueL} \textbf{Critères d'évaluation (Grille sur 20 pts)} & \textbf{5 pts} & \textbf{4 pts} & \textbf{2 pts} & \textbf{0-1 pt} \\ \hline
\textbf{Vocabulaire (5 pts)} & 5+ mots du thème utilisés correctement, variés. & 3-4 mots du thème utilisés correctement. & 1-2 mots du thème, ou répétitions. & Vocabulaire hors thème ou absent. \\ \hline
\textbf{Grammaire / Structures (5 pts)} & Structures \zh{当}+métier, \zh{可以/应该}+V, \zh{了}(changement), \zh{只}+V parfaites. & 1-2 erreurs mineures (place modal, \zh{了}). & Erreurs majeures (ordre mots, confusion \zh{是}/\zh{当}, pas de \zh{了}). & Phrases incompréhensibles, pas de structure. \\ \hline
\textbf{Caractères / Écriture (4 pts)} & Caractères corrects, proportions respectées, ordre traits OK. & 1-2 fautes d'écriture (traits manquants, radical faux). & Plusieurs fautes, écriture illisible par endroits. & Caractères inventés, radical systématiquement faux. \\ \hline
\textbf{Cohérence / Connecteurs (3 pts)} & 3 temps respectés, connecteurs (\zh{以前/现在/但是/也/还}) fluides. & Structure présente mais connecteurs maladroits. & Structure confuse, connecteurs absents ou faux. & Pas de structure logique. \\ \hline
\textbf{Respect consignes / Longueur (3 pts)} & 60-80 caractères, pinyin si demandé, 3 temps clairs. & Longueur \textasciitilde ok, un temps manquant. & Trop court (<50) ou trop long (>100), 2 temps manquants. & Non rendu ou hors sujet total. \\ \hline
\end{tabularx}
\end{center}

\courssub{Activité orale : Débat « Partage des tâches »}

\begin{experience}{Jeu de rôle : « À la maison, qui fait quoi ? »}
\textbf{Objectif :} Réinvestir \zh{应该}、\zh{一起做}、\zh{家务}、\zh{男人/女人} à l'oral.
\textbf{Déroulement (15 min) :}
\begin{enumerate}
    \item \textbf{Préparation (5 min) :} Par binôme (A = homme, B = femme). Chaque élève prépare 3 arguments pour défendre : « Les tâches doivent être partagées » (B) vs « La tradition veut que la femme gère la maison » (A - rôle difficile, jouer le diable). Vocabulaire d'aide au tableau : \zh{公平} (équitable), \zh{累} (fatigué), \zh{赚钱} (gagner de l'argent), \zh{照顾孩子} (s'occuper des enfants), \zh{做饭} (cuisiner).
    \item \textbf{Débat (7 min) :} Dialogue libre. L'enseignant circule, note les bonnes structures (\zh{我觉得应该...}、\zh{但是...也...}) et les erreurs de tons/ordre.
    \item \textbf{Restitution (3 min) :} 2-3 binômes jouent devant la classe. Correction collective bienveillante : focus sur \zh{应该一起做} vs \zh{应该谁做}.
\end{enumerate}
\end{experience}

\begin{corrige}{Exercice : Production écrite (Exemple de production attendue niveau Terminale)}
\textbf{Sujet :} \zh{我认识的一位女性} (Ma mère)
\zh{以前，我妈妈不工作，只在家里做饭、看孩子，很辛苦。\\
现在她当护士了，也会用手机看新闻，不一样了。\\
但是，她晚上还要做家务，觉得很累。\\
我觉得男人和女人应该一样，家务一起做，这样比较公平。}
\\
\textbf{Analyse :} 72 caractères. 3 temps clairs. Vocabulaire : \zh{工作}、\zh{只}、\zh{做饭}、\zh{看孩子}、\zh{当}、\zh{护士}、\zh{手机}、\zh{新闻}、\zh{不一样}、\zh{还要}、\zh{家务}、\zh{累}、\zh{应该}、\zh{一起}、\zh{公平}. Connecteurs : \zh{以前}、\zh{现在}、\zh{也}、\zh{但是}、\zh{觉得}、\zh{应该}. \zh{了} utilisé 2 fois (changement \zh{当...了}, état \zh{不一样了}). Structure \zh{当护士} correcte. Niveau HSK 3 solide.
\end{corrige}

\coursec{Lexique du thème et glossaire}

Après avoir découvert le texte modèle « La femme d'aujourd'hui » qui illustre le parcours d'une femme active à Yaoundé, nous allons maintenant bâtir l'arsenal lexical indispensable pour comprendre, traduire et rédiger sur ce thème. Comme le dit l'adage chinois : \zh{工欲善其事，必先利其器} (Gōng yù shàn qí shì, bì xiān lì qí qì) — « Pour bien faire son travail, il faut d'abord aiguiser ses outils ». Le vocabulaire est votre premier outil.

\courssub{Glossaire principal : les dix mots-clés du thème}

Observez le tableau ci-dessous. Ce sont les \cle{mots-outils} (vocabulaire de fréquence HSK 2-3) qui structureront toute votre expression écrite sur la situation de la femme. Notez la nature grammaticale (n. = nom, v. = verbe, adj. = adjectif, pron. = pronom, v. mod. = verbe modal) : elle dicte la place du mot dans la phrase.

\begin{center}
\begin{tabularx}{\linewidth}{|X|X|X|X|}
\hline
\rowcolor{popblueL} \textbf{Caractères} & \textbf{Pinyin} & \textbf{Nature} & \textbf{Traduction} \\
\hline
女人 & nǚrén & n. & femme (individu, courant) \\
\hline
妇女 & fùnǚ & n. & femmes (terme collectif, officiel, administratif) \\
\hline
她 & tā & pron. & elle (3e pers. sg. fém.) \\
\hline
工作 & gōngzuò & v./n. & travailler / travail, emploi \\
\hline
家务 & jiāwù & n. & tâches ménagères, travail domestique \\
\hline
孩子 & háizi & n. & enfant, enfants \\
\hline
平等 & píngděng & adj./n. & égal / égalité \\
\hline
权利 & quánlì & n. & droit, droits \\
\hline
地位 & dìwèi & n. & statut, position, rang social \\
\hline
应该 & yīnggāi & v. mod. & devoir, avoir l'obligation morale \\
\hline
\end{tabularx}
\end{center}

\begin{propriete}[Observations grammaticales et emploi]
\begin{itemize}
\item \textbf{女人 nǚrén vs 妇女 fùnǚ} : \zh{女人} désigne une femme (individu) ou les femmes en général dans le langage courant (\zh{这个女人很聪明} — « Cette femme est intelligente »). \zh{妇女} est un terme collectif, formel, utilisé dans les lois, les statistiques, les noms d'organisations (\zh{妇女联合会} — Fédération des femmes) et pour la fête : \zh{妇女节} (Fête des femmes, 8 mars). \textbf{Ne dites pas} \zh{*一个妇女} pour « une femme ».
\item \textbf{她 tā} : S'écrit avec le radical femme \zh{女} à gauche. À l'oral, indistinguable de \zh{他} (lui) et \zh{它} (ça/animal/objet). À l'écrit, la distinction est obligatoire.
\item \textbf{工作 gōngzuò} : Verbe \textbf{ou} nom. Ex. v. : \zh{我在医院工作} (Je travaille à l'hôpital). Ex. n. : \zh{找工作} (Chercher du travail). Le classifieur pour « un travail/un emploi » est souvent \zh{份 fèn} ou \zh{个 gè}.
\item \textbf{家务 jiāwù} : Nom collectif, \textbf{pas de classifieur} pour « faire le ménage ». On dit \zh{做家务} (faire les tâches ménagères), jamais \zh{*做一个家务}.
\item \textbf{应该 yīnggāi} : Verbe modal placé \textbf{avant} le verbe principal. Structure : Sujet + 应该 + Verbe + Complément. Ex. : \zh{女人应该有权利} (Les femmes devraient avoir des droits). Négation : \zh{不应该 bù yīnggāi}.
\end{itemize}
\end{propriete}

\courssub{Compléments lexicaux et collocations fréquentes}

Pour enrichir vos rédactions et varier les sujets (métiers, vie quotidienne, société), maîtrisez ces compléments et surtout leurs \cle{collocations} (associations mots-à-mots figées). En chinois, « faire le ménage » ne se dit pas \zh{做干净} mais \zh{做家务}.

\begin{center}
\begin{tabularx}{\linewidth}{|X|X|X|X|}
\hline
\rowcolor{poporangeL} \textbf{Caractères} & \textbf{Pinyin} & \textbf{Nature} & \textbf{Traduction} \\
\hline
经理 & jīnglǐ & n. & directeur, directrice, gérant(e) \\
\hline
医生 & yīshēng & n. & médecin, docteur \\
\hline
照顾 & zhàogù & v. & s'occuper de, prendre soin de \\
\hline
社会 & shèhuì & n. & société \\
\hline
\end{tabularx}
\end{center}

\begin{center}
\begin{tabularx}{\linewidth}{|X|X|X|}
\hline
\rowcolor{popgreenL} \textbf{Collocation (Structure V+O ou N+N)} & \textbf{Exemple en caractères + pinyin} & \textbf{Traduction} \\
\hline
做家务 zuò jiāwù (v.+n.) & 她每天做很多家务。\newline Tā měitiān zuò hěn duō jiāwù. & Elle fait beaucoup de tâches ménagères chaque jour. \\
\hline
看孩子 kàn háizi (v.+n.) & 奶奶帮忙看孩子。\newline Nǎinai bāngmáng kàn háizi. & Grand-maman aide à garder les enfants. \\
\hline
男女平等 nánnǚ píngděng (n.+n.+adj.) & 我们要实现男女平等。\newline Wǒmen yào shíxiàn nánnǚ píngděng. & Nous devons réaliser l'égalité homme-femme. \\
\hline
社会地位 shèhuì dìwèi (n.+n.) & 女性的社会地位提高了。\newline Nǚxìng de shèhuì dìwèi tígāo le. & Le statut social des femmes s'est amélioré. \\
\hline
\end{tabularx}
\end{center}

\begin{definition}[男女平等 nánnǚ píngděng]
Expression figée (结构类成语) signifiant « égalité entre l'homme et la femme ». Littéralement : « homme-femme égal/égalité ». C'est le terme standard dans les discours officiels, les médias, les rédactions scolaires et les textes de loi des deux pays. \textbf{À connaître par cœur pour l'écrit.}
\end{definition}

\begin{exemplebox}[Exemples contextualisés Cameroun/Chine]
\begin{itemize}
\item \zh{喀麦隆的许多女性当经理。} (Kāmàilóng de xǔduō nǚxìng dāng jīnglǐ.) — « De nombreuses femmes camerounaises sont directrices. »
\item \zh{男医生和女医生都应该照顾病人。} (Nán yīshēng hé nǚ yīshēng dōu yīnggāi zhàogù bìngrén.) — « Les médecins hommes et femmes doivent tous s'occuper des patients. »
\item \zh{做家务不是女人的专利。} (Zuò jiāwù bù shì nǚrén de zhuānlì.) — « Faire le ménage n'est pas l'apanage des femmes. » (Phrase utile pour un débat).
\end{itemize}
\end{exemplebox}

\begin{attention}[Piège majeur : 权利 quánlì vs 权力 quánlì]
Ces deux mots sont \textbf{homophones} (même pinyin, mêmes tons) mais leurs sens et contextes sont radicalement différents :
\begin{itemize}
\item \textbf{权利 quánlì} (droit) : Radical \zh{木} (bois) en bas à droite dans \zh{利} (via 刂). Désigne un droit légal, moral, civique. \textit{Ex. :} \zh{人权} (droits de l'homme), \zh{选举权利} (droit de vote), \zh{妇女权利} (droits des femmes). \textbf{C'est le mot du thème.}
\item \textbf{权力 quánlì} (pouvoir) : Radical \zh{力} (force) en bas à droite dans \zh{力}. Désigne le pouvoir politique, administratif, l'autorité de décision. \textit{Ex. :} \zh{行使权力} (exercer le pouvoir), \zh{滥用权力} (abuser de son pouvoir).
\end{itemize}
\textbf{Astuce mnémotechnique} : \zh{利} (lì) contient \zh{禾} (céréale/récolte) → un droit c'est un « gain » légitime. \zh{力} (lì) = force → le pouvoir c'est la « force » de commander.
\end{attention}

\courssub{Écriture et analyse des caractères clés (Radicaux, traits, ordre d'écriture)}

Avant la section dédiée à l'écriture pure, identifions ici les \cle{clés (radicaux)} qui portent le sens. Cela facilite la mémorisation, la recherche dictionnaire et l'écriture correcte.

\begin{center}
\begin{tabularx}{\linewidth}{|X|X|X|X|X|}
\hline
\rowcolor{poppurpleL} \textbf{Caractère} & \textbf{Radical (Clé)} & \textbf{Sens du radical} & \textbf{Nb traits (simplifié)} & \textbf{Caractères proches / Pièges} \\
\hline
女 nǚ & \zh{女} (nǚ) & femme & 3 & 奴 nú (esclave), 好 hǎo (bon = femme + enfant) \\
\hline
妇 fù & \zh{女} (nǚ) & femme & 8 & 妈 mā (maman), 姐 jiě (grande sœur) \\
\hline
平 píng & \zh{干} (gān) & sec, principal & 5 & 年 nián (année), 并 bìng (et/aussi) \\
\hline
权 quán & \zh{木} (mù) & bois, arbre & 6 & 模 mó (modèle), 榜 bǎng (liste/affiche) \\
\hline
地 dì & \zh{土} (tǔ) & terre, sol & 6 & 也 yě (aussi), 在 zài (à/être à) \\
\hline
应 yīng & \zh{心} (xīn) & cœur (bas) & 7 & 底 dǐ (fond), 盈 yíng (plein) \\
\hline
\end{tabularx}
\end{center}

\begin{methode}[Ordre des traits (règles de base + cas du thème)]
Respectez toujours : \textbf{1. Haut → Bas} | \textbf{2. Gauche → Droite} | \textbf{3. Horizontal → Vertical} | \textbf{4. Extérieur → Intérieur → Fermeture}.
\begin{itemize}
\item \zh{女} (3 traits) : 1. Horizontal court (haut) — 2. Vertical descendant courbe gauche — 3. Horizontal long (bas) traversant.
\item \zh{妇} (8 traits) : Écrire la clé \zh{女} (3 traits) puis le phonétique \zh{扫} (5 traits : 丶 一 丨 一 一).
\item \zh{平} (5 traits) : Clé \zh{干} : 1. Horizontal — 2. Vertical — 3. Horizontal — 4. Horizontal — 5. Horizontal (bas).
\item \zh{权} (6 traits) : Clé \zh{木} (4 traits : 一 丨 丿 丶) + \zh{又} (2 traits : 一 乙).
\item \zh{地} (6 traits) : Clé \zh{土} (3 traits : 一 丨 一) + \zh{也} (3 traits : 一 乙 一).
\item \zh{应} (7 traits) : Clé \zh{广} (3 traits : 丶 一 丿) + \zh{心} (4 traits : 丶 丶 丿 乀).
\end{itemize}
\end{methode}

\courssub{Texte d'entraînement : « Le parcours de Madame Ngo »}

Le texte ci-dessous (niveau HSK 3) réemploie \textbf{tout} le vocabulaire du glossaire. Lisez-le à haute voix, identifiez les structures, puis répondez aux questions.

\begin{textebox}[Le parcours de Madame Ngo — 恩戈女士的故事]
恩戈女士住在雅温得。她是一家大公司的经理。她工作很努力，每天早上八点去上班。\\
Ēngé Nǚshì zhù zài Yǎwēndé. Tā shì yī jiā dà gōngsī de jīnglǐ. Tā gōngzuò hěn nǔlì, měitiān zǎoshang bā diǎn qù shàngbān.\\
Madame Ngo habite à Yaoundé. Elle est directrice dans une grande entreprise. Elle travaille dur, chaque matin à 8h elle va travailler.

下班后，她回家做家务，看孩子。丈夫也帮忙做饭。\\
Xiàbān hòu, tā huí jiā zuò jiāwù, kàn háizi. Zhàngfu yě bāngmáng zuò fàn.\\
Après le travail, elle rentre faire le ménage, garder les enfants. Le mari aide aussi à cuisiner.

恩戈女士说：「男女平等很重要。妇女应该有同样的权利和地位。」\\
Ēngé Nǚshì shuō: « Nánnǚ píngděng hěn zhòngyào. Fùnǚ yīnggāi yǒu tóngyàng de quánlì hé dìwèi. »\\
Madame Ngo dit : « L'égalité homme-femme est très importante. Les femmes devraient avoir les mêmes droits et le même statut. »

现在，社会在进步，女性的地位越来越高。\\
Xiànzài, shèhuì zài jìnbù, nǚxìng de dìwèi yuè lái yuè gāo.\\
Maintenant, la société progresse, le statut des femmes est de plus en plus haut.
\end{textebox}

\begin{center}
\begin{tabularx}{\linewidth}{|X|X|X|}
\hline
\rowcolor{poppinkL} \textbf{Glossaire du texte (mots non vus)} & \textbf{Pinyin} & \textbf{Traduction} \\
\hline
恩戈 & Ēngé & Ngo (nom de famille) \\
\hline
雅温得 & Yǎwēndé & Yaoundé \\
\hline
公司 & gōngsī & entreprise, société \\
\hline
努力 & nǔlì & s'efforcer, travailler dur (adj./v.) \\
\hline
上班 & shàngbān & aller au travail, commencer le service \\
\hline
下班 & xiàbān & finir le travail, quitter le service \\
\hline
丈夫 & zhàngfu & mari \\
\hline
帮忙 & bāngmáng & aider \\
\hline
做饭 & zuò fàn & cuisiner, faire la cuisine \\
\hline
同样 & tóngyàng & même, identique \\
\hline
进步 & jìnbù & progresser, progrès \\
\hline
女性 & nǚxìng & femmes (terme formel, genre féminin) \\
\hline
越来越 & yuè lái yuè & de plus en plus \\
\hline
\end{tabularx}
\end{center}

\textbf{Questions de compréhension (réponses en français) :}
\begin{enumerate}
\item 恩戈女士在哪里工作？ (Elle travaille dans une grande entreprise / compagnie.)
\item 她下班后做什么？丈夫呢？ (Elle fait le ménage et garde les enfants. Son mari aide à cuisiner.)
\item 根据课文，妇女应该有什么？ (Les mêmes droits et le même statut que les hommes.)
\item 翻译最后一句：「现在，社会在进步，女性的地位越来越高。」 (Maintenant, la société progresse, le statut des femmes est de plus en plus élevé.)
\end{enumerate}

\courssub{Structure du texte et connecteurs utiles}

Pour rédiger un texte cohérent, il faut organiser ses idées et les relier. Analysez la structure du texte « Le parcours de Madame Ngo » :

\begin{propriete}[Analyse structurelle et connecteurs (连接词 liánjiēcí)]
\begin{center}
\begin{tabularx}{\linewidth}{|X|X|X|}
\hline
\rowcolor{popblueL} \textbf{Paragraphe / Fonction} & \textbf{Connecteurs / Marqueurs chinois} & \textbf{Exemple du texte} \\
\hline
\textbf{Introduction} (Présentation personne, lieu, métier) & \zh{是} (shì), \zh{住在} (zhù zài) & \zh{恩戈女士住在雅温得。她是一家大公司的经理。} \\
\hline
\textbf{Développement 1} (Vie professionnelle) & \zh{很} (hěn), \zh{每天} (měitiān), \zh{去} (qù) & \zh{她工作很努力，每天早上八点去上班。} \\
\hline
\textbf{Transition temporelle} & \zh{下班后} (xiàbān hòu) — « Après le travail » & \zh{下班后，她回家……} \\
\hline
\textbf{Développement 2} (Vie familiale, partage) & \zh{也} (yě) — « aussi » & \zh{丈夫也帮忙做饭。} \\
\hline
\textbf{Opinion / Argumentation} & \zh{说} (shuō) + guillemets \zh{« »}, \zh{应该} (yīnggāi) & \zh{恩戈女士说：「妇女应该有同样的权利……」。} \\
\hline
\textbf{Conclusion / Constat général} & \zh{现在} (xiànzài), \zh{在} (zài) + V (progression), \zh{越来越} (yuè lái yuè) & \zh{现在，社会在进步，女性的地位越来越高。} \\
\hline
\end{tabularx}
\end{center}

\textbf{Connecteurs à réinvestir pour votre rédaction :}
\begin{itemize}
\item \zh{首先} (shǒuxiān) — D'abord / \zh{其次} (qícì) — Ensuite / \zh{最后} (zuìhòu) — Enfin.
\item \zh{因为} (yīnwèi) … \zh{所以} (suǒyǐ) — Parce que … donc.
\item \zh{虽然} (suīrán) … \zh{但是} (dànshì) — Bien que … mais.
\item \zh{而且} (érqiě) — De plus / Et en plus (coordonne deux verbes/adjectifs).
\item \zh{跟……一样} (gēn… yīyàng) — Pareil que … (comparaison d'égalité).
\end{itemize}
\end{propriete}

\courssub{Carte mentale lexicale : le monde de la femme}

La figure ci-dessous organise le champ lexical autour du nœud central \zh{女人}. Mémorisez les quatre branches : elles correspondent aux quatre angles d'attaque d'une rédaction (vie familiale, vie professionnelle, droits, société).

\begin{popfigure}
\begin{tikzpicture}[
  mindmap,
  grow cyclic,
  every node/.style={concept, execute at begin node=\hskip0pt},
  concept color=popblue!80,
  level 1/.append style={level distance=4.5cm, sibling angle=90, concept color=popblue!60},
  level 2/.append style={level distance=3cm, sibling angle=45, concept color=poporange!70},
  text=white,
  font=\small\bfseries
]
\node {女人}
  child [concept color=popgreen!70] { node {Famille}
    child { node {家务 jiāwù} }
    child { node {孩子 háizi} }
    child { node {照顾 zhàogù} }
  }
  child [concept color=poporange!70] { node {Travail}
    child { node {工作 gōngzuò} }
    child { node {经理 jīnglǐ} }
    child { node {医生 yīshēng} }
  }
  child [concept color=poppurple!70] { node {Droits}
    child { node {平等 píngděng} }
    child { node {权利 quánlì} }
    child { node {应该 yīnggāi} }
  }
  child [concept color=poppink!70] { node {Société}
    child { node {地位 dìwèi} }
    child { node {社会 shèhuì} }
    child { node {妇女 fùnǚ} }
  };
\end{tikzpicture}
\legende{Carte mentale du vocabulaire « Situation de la femme » — Branches : Famille, Travail, Droits, Société.}
\end{popfigure}

\courssub{Thème guidé : du français vers le chinois (Méthode)}

Traduire (faire du thème) n'est pas remplacer mot à mot. Suivez cette méthode en 4 étapes pour la phrase : \textit{« Les femmes camerounaises devraient avoir le même statut social que les hommes. »}

\begin{methode}[Méthode du thème en 4 étapes]
\begin{enumerate}
\item \textbf{Analyser la phrase française} : Sujet = « Les femmes camerounaises » (groupe nominal). Verbe modal + verbe = « devraient avoir ». COD = « le même statut social ». Complément de comparaison = « que les hommes ».
\item \textbf{Identifier la structure chinoise cible} : Sujet + 应该 + 有 + COD + 跟/和 + Comparé + 一样 / 同样. Structure de comparaison d'égalité : \zh{A 跟 B 一样} ou \zh{A 和 B 同样 + Adj/Verbe}.
\item \textbf{Choisir le vocabulaire précis} : 
\begin{itemize}
\item Femmes camerounaises = \zh{喀麦隆的妇女} (terme officiel) ou \zh{喀麦隆的女性} (formel).
\item Statut social = \zh{社会地位}.
\item Même = \zh{同样的} ou \zh{一样}.
\item Hommes = \zh{男人} ou \zh{男性}.
\end{itemize}
\item \textbf{Assembler et vérifier l'ordre} : 
\zh{喀麦隆的妇女 应该 有 跟 男人 一样 的 社会地位。} \\
Kāmàilóng de fùnǚ yīnggāi yǒu gēn nánrén yīyàng de shèhuì dìwèi.
\end{enumerate}
\end{methode}

\begin{exoresolu}[Thème : « Madame Meng est médecin, elle s'occupe des patients et fait aussi le ménage. »]
\textbf{Énoncé :} Traduire en chinois.
\textbf{Solution détaillée :}
\begin{enumerate}
\item \textit{Analyse} : Deux propositions coordonnées par « et » (\zh{也} ou \zh{还}). Sujet unique : Madame Meng.
\item \textit{Vocabulaire} : Madame Meng = \zh{孟医生 Mèng Yīshēng} ou \zh{孟女士 Mèng Nǚshì}. Médecin = \zh{医生}. S'occuper de = \zh{照顾}. Patients = \zh{病人 bìngrén}. Faire le ménage = \zh{做家务}.
\item \textit{Assemblage} : 
\zh{孟医生照顾病人，也做家务。} (Mèng Yīshēng zhàogù bìngrén, yě zuò jiāwù.)
\item \textit{Vérification} : Ordre SOV respecté. \zh{也} placé avant le verbe \zh{做}. Ponctuation chinoise 。
\end{enumerate}
\end{exoresolu}

\courssub{Version guidée : du chinois vers le français (Méthode)}

La version exige de respecter le sens profond et le niveau de langue français. Appliquons la méthode au dernier paragraphe du texte modèle : \zh{现在，社会在进步，女性的地位越来越高。}

\begin{methode}[Méthode de la version en 4 étapes]
\begin{enumerate}
\item \textbf{Segmenter} : Couper en unités de sens (propositions). Ici : ① \zh{现在} (Maintenant) ② \zh{社会在进步} (la société progresse) ③ \zh{女性的地位越来越高} (le statut des femmes est de plus en plus haut).
\item \textbf{Analyser la structure de chaque segment} :
\begin{itemize}
\item ② Sujet \zh{社会} + \zh{在} (marque progressif) + Verbe \zh{进步}.
\item ③ Sujet \zh{女性的地位} (Groupe nominal : Femmes + 的 + Statut) + \zh{越来越} (de plus en plus) + Adjectif \zh{高}.
\end{itemize}
\item \textbf{Traduire le vocabulaire clé} : \zh{现在}=Maintenant ; \zh{社会}=société ; \zh{进步}=progresser ; \zh{女性}=femmes (formel) ; \zh{地位}=statut/position ; \zh{越来越}=de plus en plus ; \zh{élevé/élevé}.
\item \textbf{Recomposer en français naturel} : Ajuster la ponctuation, les temps, les articles. \rightarrow « Aujourd'hui, la société progresse et le statut des femmes ne cesse de s'élever. » (Ou : « est de plus en plus élevé. »)
\end{enumerate}
\end{methode}

\begin{exoresolu}[Version : \zh{在喀麦隆，很多女性既工作又照顾家庭。男女平等是基本权利。}]
\textbf{Énoncé :} Traduire en français.
\textbf{Solution détaillée :}
\begin{enumerate}
\item \textit{Segmentation} : ① \zh{在喀麦隆} (Au Cameroun) ② \zh{很多女性} (beaucoup de femmes) ③ \zh{既……又……} (à la fois … et …) ④ \zh{工作} (travaillent) ⑤ \zh{照顾家庭} (s'occupent de la famille) 。 ⑥ \zh{男女平等} (L'égalité homme-femme) ⑦ \zh{是} (est) ⑧ \zh{基本权利} (un droit fondamental).
\item \textit{Traduction} : « Au Cameroun, beaucoup de femmes travaillent tout en s'occupant de leur famille. L'égalité entre les hommes et les femmes est un droit fondamental. »
\item \textit{Note} : \zh{既……又……} structure corrélative = « à la fois … et … » / « aussi bien … que … ». \zh{基本} = fondamental / de base.
\end{enumerate}
\end{exoresolu}

\courssub{Le savais-tu ? : Le 8 mars, pont entre Yaoundé et Beijing}

\begin{savaistu}[La Journée internationale des droits des femmes]
Au Cameroun comme en Chine, le \textbf{8 mars} (\zh{三八妇女节 Sānbā Fùnǚ Jié}) est une fête officielle majeure.
\begin{itemize}
\item \textbf{Au Cameroun} : Défilés dans les grandes villes (Yaoundé, Douala, Bafoussam), port du pagne spécial « 8 mars » (\zh{妇女节布 fùnǚ jié bù}), discours officiels sur l'autonomisation des femmes (\zh{妇女自主 fùnǚ zìzhǔ}), remise de médailles aux femmes méritantes.
\item \textbf{En Chine} : Journée fériée de demi-journée pour les femmes (\zh{放半天假 fàng bàntiān jià}). Les entreprises offrent des cadeaux (\zh{礼物 lǐwù}), organisent des sorties. C'est aussi une journée de consommation massive (\zh{女神节 Nǚshén Jié} « Fête des déesses », terme marketing apparu récemment sur internet).
\item \textbf{Point commun} : Revendication de l'\textbf{égalité professionnelle} (\zh{职场平等 zhíchǎng píngděng}) et de la \textbf{répartition des tâches ménagères} (\zh{家务分担 jiāwù fēndān}).
\end{itemize}
\end{savaistu}

\courssub{Production écrite : sujet d'entraînement et grille d'évaluation}

\begin{exercice}[Production écrite — Sujet d'examen type]
\textbf{Consigne :} Rédigez un paragraphe d'environ \textbf{80 à 100 caractères chinois} sur le thème : \textit{« La situation de la femme dans la société camerounaise aujourd'hui »}.
\textbf{Contraintes obligatoires :}
\begin{enumerate}
\item Utilisez \textbf{au moins 7 mots} du glossaire principal (femme, travail, tâches ménagères, enfants, égalité, droits, statut, devrait).
\item Utilisez \textbf{au moins 3 connecteurs} vus en section 4 (ex: 首先, 但是, 现在, 越来越, 因为…所以…).
\item Respectez l'ordre des mots chinois (Sujet + Temps + Lieu + Verbe + Complément).
\item Soignez l'écriture des caractères (traits, proportions).
\end{enumerate}
\textbf{Amorce suggérée :} \zh{在喀麦隆，现在的妇女……} (Au Cameroun, les femmes d'aujourd'hui…)
\end{exercice}

\begin{aretenir}[Grille d'évaluation (20 points) — Auto-évaluation / Évaluation par les pairs]
\begin{center}
\begin{tabularx}{\linewidth}{|X|X|X|}
\hline
\rowcolor{popgoldL} \textbf{Critère} & \textbf{Indicateurs de réussite} & \textbf{Points} \\
\hline
\textbf{Richesse lexicale} & 7+ mots du glossaire employés correctement (nature, place). Variété (pas de répétition excessive). & /5 \\
\hline
\textbf{Exactitude grammaticale} & Structures justes (应该+V, 跟…一样, 在+V, 越来越+Adj). Pas d'erreur sur 的/得/地. Négation correcte. & /5 \\
\hline
\textbf{Cohérence \& Connecteurs} & 3+ connecteurs utilisés à bon escient. Progression logique (Intro → Développement → Conclusion). Ponctuation chinoise. & /5 \\
\hline
\textbf{Écriture des caractères} & Caractères lisibles, traits dans l'ordre, proportions respectées (clé + phonétique), pas de confusion visuelle (ex: 权/利, 女/母). & /5 \\
\hline
\end{tabularx}
\end{center}
\end{aretenir}

\begin{exercice}[Entraînement : Thème et Vocabulaire (Renforcement)]
\textbf{A. Complétez les phrases avec le mot juste du glossaire (caractères + pinyin) :}
1. \zh{我的母亲每天 \_\_\_\_\_\_ ，很辛苦。} (fait le ménage)
2. \zh{我们要争取男女 \_\_\_\_\_\_ 。} (égalité)
3. \zh{\_\_\_\_\_\_ 是一家医院的 \_\_\_\_\_\_ 。} (Elle / médecin)
4. \zh{妇女和男人有同样的 \_\_\_\_\_\_ 。} (droits)
5. \zh{社会 \_\_\_\_\_\_ 在提高。} (statut)

\textbf{B. Traduisez en chinois (Thème) :}
1. Les femmes font beaucoup de travail à la maison.
2. Mon père aide à garder les enfants.
3. Les femmes devraient avoir un bon travail.
4. L'égalité homme-femme est un droit.
5. Le statut social des femmes chinoises est haut.
\end{exercice}

\begin{corrige}[Exercice Lexique du thème]
\textbf{A. Complétion :}
1. \zh{做家务 zuò jiāwù}
2. \zh{平等 píngděng}
3. \zh{她 tā} / \zh{医生 yīshēng}
4. \zh{权利 quánlì}
5. \zh{地位 dìwèi}

\textbf{B. Thème (propositions de réponses standards) :}
1. \zh{女人做很多家务。} (Ou \zh{妇女做很多家务。})
2. \zh{我爸爸帮忙看孩子。}
3. \zh{妇女应该有好的工作。} (Ou \zh{女人应该找好工作。})
4. \zh{男女平等是权利。}
5. \zh{中国女性的社会地位很高。}
\end{corrige}

\coursec{Écriture des caractères : traits, radicaux, caractères proches}

Après avoir acquis le lexique essentiel pour décrire la situation de la femme dans la société camerounaise et chinoise (\emph{lexique du thème et glossaire}), nous devons maintenant maîtriser l'écriture correcte des caractères qui portent ce sens. En chinois, l'écriture n'est pas un simple exercice calligraphique : l'ordre des traits, la structure du caractère et le choix du radical conditionnent la lisibilité, la mémorisation et la distinction de sens. Cette section vous guide pas à pas, du radical fondamental \zh{女} aux caractères composés du thème, en insistant sur les pièges classiques pour les francophones.

\courssub{Le radical clé : 女 (nǚ) — « femme »}

\begin{textebox}[Observation : le radical 女 dans le vocabulaire du thème]
Regardons les mots-clés de notre liste de vocabulaire : \zh{她} (tā, elle), \zh{妇女} (fùnǚ, femme), \zh{好} (hǎo, bon/bien), \zh{妈妈} (māma, maman), \zh{姐姐} (jiějie, grande sœur), \zh{妹妹} (mèimei, petite sœur). Tous partagent un élément graphique commun à gauche : \zh{女}.
\end{textebox}

Le caractère \zh{女} (nǚ) est à la fois un caractère autonome (« femme », « féminin ») et un \cle{radical} (部首, bùshǒu) — clé de dictionnaire — qui signale une relation avec le monde féminin : personnes de sexe féminin, qualités traditionnellement associées aux femmes, relations familiales maternelles ou sororales.

\begin{propriete}[Ordre des traits de 女 (nǚ) — 3 traits]
L'écriture suit deux principes universels : \textbf{de haut en bas} (从上到下) et \textbf{de gauche à droite} (从左到右). Pour 女, l'ordre est le suivant :
\begin{enumerate}
    \item \textbf{撇点 (piědiǎn)} : trait brisé « pointe-pique » — part du haut-gauche, descend en courbe vers la droite, se termine par un petit crochet vers le bas. C'est le trait « maître » qui donne l'inclinaison.
    \item \textbf{撇 (piě)} : trait « pique » — part du milieu-haut, descend vers la gauche en s'évasant. Il équilibre le premier trait.
    \item \textbf{横 (héng)} : trait « horizontal » — part de la jonction des deux traits précédents, va vers la droite. Il « ferme » le caractère et lui donne sa base.
\end{enumerate}
\emph{Astuce mnémotechnique :} imaginez une femme qui marche gracieusement : le premier trait est la tête et l'épaule, le second la hanche, le troisième le sol sur lequel elle pose le pied.
\end{propriete}

\begin{popfigure}
\begin{tikzpicture}[scale=1.2, every node/.style={font=\small}]
    % Grille de fond
    \draw[popline, step=0.5cm] (0,0) grid (3,3);
    \draw[popdark, very thick] (0,0) rectangle (3,3);
    
    % Étape 1 : Trait 1 (piědiǎn)
    \begin{scope}[xshift=0cm]
        \draw[popblue, very thick, line cap=round] (1.5,2.7) .. controls (1.8,2.0) and (1.2,1.2) .. (1.5,0.8);
        \node[popblue, align=center] at (1.5,3.3) {Étape 1\\Trait 1 : 撇点};
        \node[popblue, circle, fill=popblue, text=white, font=\tiny\bfseries] at (1.5,2.7) {1};
    \end{scope}
    
    % Étape 2 : Trait 1 + 2
    \begin{scope}[xshift=4cm]
        \draw[popblue, very thick, line cap=round] (1.5,2.7) .. controls (1.8,2.0) and (1.2,1.2) .. (1.5,0.8);
        \draw[poporange, very thick, line cap=round] (1.5,2.2) .. controls (0.8,1.8) and (0.5,1.0) .. (0.8,0.3);
        \node[poporange, align=center] at (1.5,3.3) {Étape 2\\Trait 2 : 撇};
        \node[poporange, circle, fill=poporange, text=white, font=\tiny\bfseries] at (1.5,2.2) {2};
    \end{scope}
    
    % Étape 3 : Complet
    \begin{scope}[xshift=8cm]
        \draw[popblue, very thick, line cap=round] (1.5,2.7) .. controls (1.8,2.0) and (1.2,1.2) .. (1.5,0.8);
        \draw[poporange, very thick, line cap=round] (1.5,2.2) .. controls (0.8,1.8) and (0.5,1.0) .. (0.8,0.3);
        \draw[popgreen, very thick, line cap=round] (0.8,0.8) -- (2.7,0.8);
        \node[popgreen, align=center] at (1.5,3.3) {Étape 3\\Trait 3 : 横};
        \node[popgreen, circle, fill=popgreen, text=white, font=\tiny\bfseries] at (2.7,0.8) {3};
    \end{scope}
\end{tikzpicture}
\legende{Ordre des traits du radical 女 (nǚ) en 3 étapes. Couleurs : bleu = trait 1, orange = trait 2, vert = trait 3.}
\end{popfigure}

\courssub{Écriture des caractères composés avec 女 : 她, 妇, 好}

Les caractères du thème sont des \cle{caractères composés} (合字, hézì) de structure \textbf{gauche-droite} (左边结构) : le radical \zh{女} occupe la partie gauche (sémantique), une composante phonétique ou sémantique occupe la droite. Règle d'or : \textbf{on écrit d'abord la partie gauche (le radical), puis la partie droite}.

\begin{propriete}[Règle d'ordre des traits pour les caractères gauche-droite (左边结构)]
Pour tout caractère dont le radical est à gauche (如 \zh{她、妇、好、妈、姐、妹}) :
\begin{enumerate}
    \item Écrire complètement le radical de gauche (ici \zh{女}, 3 traits) en respectant son ordre propre.
    \item Écrire ensuite la partie de droite, trait par trait, selon ses propres règles (haut→bas, gauche→droite).
\end{enumerate}
\emph{Ne jamais « mélanger » les traits des deux côtés.} Cela garantit l'équilibre graphique et la vitesse d'écriture.
\end{propriete}

\begin{exemplebox}[Décomposition et ordre des traits : 她 (tā), 妇 (fù), 好 (hǎo)]
\begin{center}
\begin{tabularx}{\linewidth}{|X|X|X|X|}
\hline
\rowcolor{popblueL} \textbf{Caractère} & \textbf{Décomposition} & \textbf{Nb traits total} & \textbf{Ordre des traits (résumé)} \\
\hline
\zh{她} (tā) & 女 (3) + 也 (3) = 6 traits & 6 & 1-2-3 (女) → 4-5-6 (也: 横, 撇, 竖钩) \\
\hline
\zh{妇} (fù) & 女 (3) + 彐 (3) = 6 traits & 6 & 1-2-3 (女) → 4-5-6 (彐: 横折, 横, 竖提) \\
\hline
\zh{好} (hǎo) & 女 (3) + 子 (3) = 6 traits & 6 & 1-2-3 (女) → 4-5-6 (子: 横, 竖钩, 横) \\
\hline
\end{tabularx}
\end{center}
\end{exemplebox}

\begin{popfigure}
\begin{tikzpicture}[scale=1.1, every node/.style={font=\small}]
    % Étape A : 女 seul
    \begin{scope}[xshift=0cm, yshift=0cm]
        \draw[popline, step=0.5cm] (0,0) grid (3,3);
        \draw[popdark, very thick] (0,0) rectangle (3,3);
        \draw[popblue, very thick, line cap=round] (1.5,2.7) .. controls (1.8,2.0) and (1.2,1.2) .. (1.5,0.8);
        \draw[popblue, very thick, line cap=round] (1.5,2.2) .. controls (0.8,1.8) and (0.5,1.0) .. (0.8,0.3);
        \draw[popblue, very thick, line cap=round] (0.8,0.8) -- (2.7,0.8);
        \node[align=center] at (1.5,3.4) {\textbf{她} : Radical 女 complet};
        \node[popblue, circle, fill=popblue, text=white, font=\tiny\bfseries] at (2.7,0.8) {3};
    \end{scope}
    
    % Étape B : + 1er trait de 也 (横)
    \begin{scope}[xshift=4cm, yshift=0cm]
        \draw[popline, step=0.5cm] (0,0) grid (3,3);
        \draw[popdark, very thick] (0,0) rectangle (3,3);
        % 女
        \draw[popblue, very thick, line cap=round] (1.5,2.7) .. controls (1.8,2.0) and (1.2,1.2) .. (1.5,0.8);
        \draw[popblue, very thick, line cap=round] (1.5,2.2) .. controls (0.8,1.8) and (0.5,1.0) .. (0.8,0.3);
        \draw[popblue, very thick, line cap=round] (0.8,0.8) -- (2.7,0.8);
        % 也 - trait 1 (横)
        \draw[poporange, very thick, line cap=round] (0.5,2.0) -- (2.5,2.0);
        \node[align=center] at (1.5,3.4) {+ Trait 4 : 横 (de 也)};
        \node[poporange, circle, fill=poporange, text=white, font=\tiny\bfseries] at (2.5,2.0) {4};
    \end{scope}
    
    % Étape C : 也 complet (3 traits pour 也 = 6 traits total)
    \begin{scope}[xshift=8cm, yshift=0cm]
        \draw[popline, step=0.5cm] (0,0) grid (3,3);
        \draw[popdark, very thick] (0,0) rectangle (3,3);
        % 女
        \draw[popblue, very thick, line cap=round] (1.5,2.7) .. controls (1.8,2.0) and (1.2,1.2) .. (1.5,0.8);
        \draw[popblue, very thick, line cap=round] (1.5,2.2) .. controls (0.8,1.8) and (0.5,1.0) .. (0.8,0.3);
        \draw[popblue, very thick, line cap=round] (0.8,0.8) -- (2.7,0.8);
        % 也 complet (simplifié : 3 traits)
        \draw[poporange, very thick, line cap=round] (0.5,2.0) -- (2.5,2.0); % 横
        \draw[poporange, very thick, line cap=round] (1.5,2.0) .. controls (1.0,1.5) .. (0.8,1.0); % 撇
        \draw[poporange, very thick, line cap=round] (1.5,2.0) -- (1.5,0.5) .. controls (1.5,0.3) and (1.8,0.1) .. (2.0,0.1); % 竖钩
        \node[align=center] at (1.5,3.4) {\textbf{她} complet (6 traits)};
        \node[poporange, circle, fill=poporange, text=white, font=\tiny\bfseries] at (2.0,0.1) {6};
    \end{scope}
\end{tikzpicture}
\legende{Construction progressive de 她 (tā) : radical 女 (3 traits, bleu) puis partie droite 也 (3 traits, orange). Total : 6 traits (norme simplifiée).}
\end{popfigure}

\courssub{Caractères proches à ne pas confondre : l'importance du radical}

C'est le point critique pour l'expression écrite et la traduction. En français, « il » et « elle » s'écrivent différemment mais se prononcent différemment aussi. En chinois, \zh{他} (tā, il) et \zh{她} (tā, elle) sont \textbf{homophones} (même pinyin, même ton) mais \textbf{hétérographes} (écritures différentes) grâce au radical.

\begin{attention}[Piège majeur : 她 (tā) vs 他 (tā) — « elle » vs « il »]
\begin{itemize}
    \item \zh{他} (tā) : radical \zh{亻} (rén, « homme », forme compressée de \zh{人}) + \zh{也}. Sens : « il », « lui » (masculin ou générique).
    \item \zh{她} (tā) : radical \zh{女} (nǚ, « femme ») + \zh{也}. Sens : « elle » (féminin).
\end{itemize}
\emph{Erreur fréquente des francophones :} à l'oral, on entend \emph{tā} ; à l'écrit, on écrit machinalement \zh{他} (plus fréquent, forme « par défaut » historique). \textbf{Consigne :} dès que le contexte désigne une femme (mère, sœur, fille, femme, Madame X), le radical \zh{女} est \textbf{obligatoire}. Ne pas le mettre change le genre de la personne et fait perdre des points à l'examen (écrit et traduction).
\end{attention}

\begin{popfigure}
\begin{tikzpicture}[scale=1.3, every node/.style={font=\small}]
    % Tableau comparatif 她 vs 他
    \draw[popdark, thick] (0,0) rectangle (10,4);
    \draw[popdark, thick] (5,0) -- (5,4); % séparation milieu
    \draw[popdark, thick] (0,1) -- (10,1); % ligne titre
    \draw[popdark, thick] (0,2.5) -- (10,2.5); % ligne séparation radical/corps
    
    % En-têtes
    \node[popblue, align=center, font=\bfseries] at (2.5,3.3) {\zh{她} (tā) — « elle »};
    \node[poporange, align=center, font=\bfseries] at (7.5,3.3) {\zh{他} (tā) — « il »};
    
    % Ligne Radical
    \node[align=center] at (2.5,1.75) {Radical (gauche)};
    \node[align=center] at (7.5,1.75) {Radical (gauche)};
    
    % Radical 女 en bleu
    \begin{scope}[xshift=1.5cm, yshift=0.2cm, scale=0.6]
        \draw[popblue, very thick, line cap=round] (1.5,2.7) .. controls (1.8,2.0) and (1.2,1.2) .. (1.5,0.8);
        \draw[popblue, very thick, line cap=round] (1.5,2.2) .. controls (0.8,1.8) and (0.5,1.0) .. (0.8,0.3);
        \draw[popblue, very thick, line cap=round] (0.8,0.8) -- (2.7,0.8);
    \end{scope}
    \node[popblue, align=center, font=\bfseries] at (2.5,0.5) {\zh{女} (nǚ)\\femme};
    
    % Radical 亻 en orange
    \begin{scope}[xshift=6.5cm, yshift=0.5cm, scale=0.7]
        \draw[poporange, very thick, line cap=round] (0,2.5) -- (0,0.2); % verticale
        \draw[poporange, very thick, line cap=round] (0,2.0) -- (-0.5,1.0); % bras gauche
        \draw[poporange, very thick, line cap=round] (0,2.0) -- (0.5,1.0); % bras droit
    \end{scope}
    \node[poporange, align=center, font=\bfseries] at (7.5,0.5) {\zh{亻} (rén)\\homme};
    
    % Partie droite commune (也)
    \node[align=center] at (2.5,3.0) {Partie droite (phonétique)};
    \node[align=center] at (7.5,3.0) {Partie droite (phonétique)};
    \begin{scope}[xshift=1.5cm, yshift=1.5cm, scale=0.5]
        \draw[popdark, thick] (0.5,2.0) -- (2.5,2.0);
        \draw[popdark, thick] (1.5,2.0) .. controls (1.0,1.5) .. (0.8,1.0);
        \draw[popdark, thick] (1.5,2.0) -- (1.5,0.5) .. controls (1.5,0.3) and (1.8,0.1) .. (2.0,0.1);
    \end{scope}
    \begin{scope}[xshift=6.5cm, yshift=1.5cm, scale=0.5]
        \draw[popdark, thick] (0.5,2.0) -- (2.5,2.0);
        \draw[popdark, thick] (1.5,2.0) .. controls (1.0,1.5) .. (0.8,1.0);
        \draw[popdark, thick] (1.5,2.0) -- (1.5,0.5) .. controls (1.5,0.3) and (1.8,0.1) .. (2.0,0.1);
    \end{scope}
    \node[popdark, align=center, font=\bfseries] at (2.5,1.0) {\zh{也} (yě)};
    \node[popdark, align=center, font=\bfseries] at (7.5,1.0) {\zh{也} (yě)};
\end{tikzpicture}
\legende{Comparaison structurelle : 她 (radical femme 女) vs 他 (radical homme 亻). La partie droite 也 (yě) est identique (indice phonétique). La couleur du radical signale le genre.}
\end{popfigure}

\begin{exemplebox}[Autres paires de confusion à surveiller]
\begin{center}
\begin{tabularx}{\linewidth}{|X|X|X|X|}
\hline
\rowcolor{popblueL} \textbf{Caractère thème} & \textbf{Pinyin / Sens} & \textbf{Caractère proche} & \textbf{Pinyin / Sens / Différence} \\
\hline
\zh{妇} & fù — femme, épouse & \zh{扫} & sǎo — balayer. \textbf{Différence :} bas du caractère : \zh{彐} (jié, tête de porc stylisée) vs \zh{帚} (zhǒu, balai → main \zh{扌} + balai). \zh{扫} a le radical main \zh{扌} (action manuelle). \\
\hline
\zh{好} & hǎo — bon, bien, aimer & \zh{如} & rú — comme, si, ressembler. \textbf{Différence :} partie droite : \zh{子} (zǐ, enfant) vs \zh{口} (kǒu, bouche). \zh{如} = \zh{女} + \zh{口} (femme + bouche = « comme »/obéissance). \zh{好} = \zh{女} + \zh{子} (femme + enfant = « bon »). \\
\hline
\end{tabularx}
\end{center}
\end{exemplebox}

\begin{methode}[Comment mémoriser un caractère composé (radical + phonétique) — 5 étapes]
\begin{enumerate}
    \item \textbf{Décomposer :} identifier le radical (sens) et la partie droite (souvent indice de prononciation). Ex: \zh{妇} = \zh{女} (femme) + \zh{彐} (son « ü »/« u »).
    \item \textbf{Écrire en l'air (空写, kōngxiě) :} tracer le caractère avec l'index en disant l'ordre des traits à voix haute : « piědiǎn, piě, héng, zhé, héng, tí ».
    \item \textbf{Copier 5 fois sur quadrillage :} respecter la proportion (radical gauche ≈ 1/3, partie droite ≈ 2/3) et l'alignement horizontal (le trait horizontal de base du radical s'aligne souvent avec un trait de la partie droite).
    \item \textbf{Associer son + sens :} écrire le pinyin \emph{fù} à côté et dire une phrase : « \zh{妇女} (fùnǚ) — les femmes ».
    \item \textbf{Réviser espacé (J+1, J+3, J+7) :} réécrire le caractère sans modèle, vérifier l'ordre des traits.
\end{enumerate}
\end{methode}

\courssub{Exemples traduits : ancrage du radical 女 dans la phrase}

\begin{exemplebox}[Phrases modèles : distinction 她 / 他 et emploi de 妇, 好]
\begin{itemize}
    \item \zh{她是我的老师。} \emph{Tā shì wǒ de lǎoshī.} « Elle est mon professeur. » (Radical 女 : personne féminine)
    \item \zh{他是我的同学。} \emph{Tā shì wǒ de tóngxué.} « Il est mon camarade de classe. » (Radical 亻 : personne masculine)
    \item \zh{妇女能顶半边天。} \emph{Fùnǚ néng dǐng bànbiāntiān.} « Les femmes tiennent la moitié du ciel. » (Proverbe chinois, Mao Zedong. \zh{妇女} = femmes).
    \item \zh{我妈妈很好。} \emph{Wǒ māma hěn hǎo.} « Ma mère va très bien / est très gentille. » (\zh{好} = femme + enfant → bon).
    \item \zh{她姐姐喜欢跳舞。} \emph{Tā jiějie xǐhuan tiàowǔ.} « Sa grande sœur aime danser. » (\zh{姐} = femme + radical droite).
    \item \zh{我们要尊重妇女的权利。} \emph{Wǒmen yào zūnzhòng fùnǚ de quánlì.} « Nous devons respecter les droits des femmes. »
\end{itemize}
\end{exemplebox}

\courssub{Exercice de copie guidée et discrimination}

\begin{exoresolu}[Exercice guidé : copie, ordre des traits et choix du bon caractère]
\textbf{Consigne :} Recopiez chaque caractère 3 fois dans l'ordre des traits indiqué. Puis complétez les phrases en choisissant \zh{他} ou \zh{她}.

\textbf{1. Copie guidée (respectez l'ordre des traits numérotés) :}
\begin{center}
\begin{tabular}{c|c|c|c}
\zh{女} (1-2-3) & \zh{女} \zh{女} \zh{女} & \zh{她} (1-6) & \zh{她} \zh{她} \zh{她} \\
\zh{妇} (1-6) & \zh{妇} \zh{妇} \zh{妇} & \zh{好} (1-6) & \zh{好} \zh{好} \zh{好} \\
\end{tabular}
\end{center}

\textbf{2. Discrimination 她 / 他 :}
\begin{enumerate}
    \item Ma mère est médecin. \zh{我妈\_\_\_是医生。} → \zh{她} (contexte : mère = féminin).
    \item Mon père est professeur. \zh{我爸\_\_\_是老师。} → \zh{他} (contexte : père = masculin).
    \item Cette fille est camerounaise. \zh{\_\_\_是喀麦隆女孩。} → \zh{她} (sujet : fille = féminin).
    \item Ce garçon est chinois. \zh{\_\_\_是中国男孩。} → \zh{他} (sujet : garçon = masculin).
    \item Madame Li est très gentille. \zh{李女士\_\_\_很\_\_\_。} → \zh{她} ... \zh{好} (Madame = féminin ; \zh{好} = bon/gentil).
\end{enumerate}

\textbf{Correction détaillée :}
\begin{itemize}
    \item[1-4.] Le choix du radical (\zh{女} vs \zh{亻}) est dicté par le genre naturel de l'antécédent (mère/fille vs père/garçon). En chinois, pas de « genre neutre » pour les personnes : on \textbf{doit} choisir.
    \item[5.] \zh{李女士} (Lǐ Nǚshì) ou \zh{李太太} (Lǐ Tàitai) → pronom \zh{她}. L'adjectif « gentille » se traduit par \zh{好} (hǎo) ou \zh{善良} (shànliáng). Ici \zh{好} suffit (HSK 2-3).
\end{itemize}
\tcblower
\textbf{Solution détaillée :}
1. \zh{我妈\textbf{她}是医生。} (Wǒ mā \textbf{tā} shì yīshēng.)
2. \zh{我爸\textbf{他}是老师。} (Wǒ bà \textbf{tā} shì lǎoshī.)
3. \zh{\textbf{她}是喀麦隆女孩。} (\textbf{Tā} shì Kāmàilóng nǚhái.)
4. \zh{\textbf{他}是中国男孩。} (\textbf{Tā} shì Zhōngguó nánhái.)
5. \zh{李女士\textbf{她}很\textbf{好}。} (Lǐ Nǚshì \textbf{tā} hěn \textbf{hǎo}.)
\end{exoresolu}

\begin{exercice}[Entraînement personnel : écriture et traduction]
\begin{enumerate}
    \item Écrivez 5 fois chacun des caractères suivants sur votre cahier à quadrillage (format Tianzige 田字格) en respectant l'ordre des traits : \zh{女}, \zh{她}, \zh{妇}, \zh{好}, \zh{妈}, \zh{姐}, \zh{妹}. Cochez chaque trait à voix haute.
    \item Traduisez en chinois (caractères + pinyin) :
    \begin{enumerate}
        \item Ma grande sœur est une femme courageuse.
        \item Les femmes camerounaises aiment le ndolé.
        \item Elle n'est pas mon professeur, c'est ma mère.
        \item Ce n'est pas lui, c'est elle.
    \end{enumerate}
    \item Cerclez l'intrus dans chaque série et justifiez par le radical :
    \begin{enumerate}
        \item \zh{妈 姐 妹 他}
        \item \zh{好 妇 她 扫}
    \end{enumerate}
\end{enumerate}
\end{exercice}

\begin{corrige}[Exercice n°3 — Écriture et traduction]
\textbf{1. Écriture :} Vérification par le professeur (ordre traits, proportions, alignement horizontal du trait 3 de \zh{女} avec le trait horizontal médian de la partie droite).

\textbf{2. Traduction :}
\begin{enumerate}
    \item \zh{我姐姐是勇敢的妇女。} (Wǒ jiějie shì yǒnggǎn de fùnǚ.) — \emph{勇敢 yǒnggǎn = courageux.}
    \item \zh{喀麦隆妇女喜欢吃恩多莱。} (Kāmàilóng fùnǚ xǐhuan chī Ēnduōlái.) — \emph{Ndolé translittéré : 恩多莱 Ēnduōlái.}
    \item \zh{她不是我的老师，是我的妈妈。} (Tā bù shì wǒ de lǎoshī, shì wǒ de māma.)
    \item \zh{不是他，是她。} (Bù shì tā, shì tā.) — \emph{À l'oral indistinguable ; à l'écrit, le radical fait tout.}
\end{enumerate}

\textbf{3. Intrus :}
\begin{enumerate}
    \item \zh{他} (radical \zh{亻} homme) — les autres ont \zh{女} (femme).
    \item \zh{扫} (radical \zh{扌} main/action) — les autres ont \zh{女} (femme). \zh{好} a \zh{女} + \zh{子}.
\end{enumerate}
\end{corrige}

\begin{aretenir}[Synthèse : Le radical 女, votre repère « féminin »]
\begin{itemize}
    \item \textbf{Forme :} 3 traits (piědiǎn, piě, héng). S'écrit en premier (gauche) dans les composés.
    \item \textbf{Fonction :} Marque le genre féminin (pronoms \zh{她}, parenté \zh{妈、姐、妹}, collectif \zh{妇女}) ou la notion de « bien/bon » (\zh{好} = femme + enfant).
    \item \textbf{Piège n°1 :} \zh{他} (il) vs \zh{她} (elle) — même son \emph{tā}, radical différent. \textbf{Toujours vérifier le contexte de genre avant d'écrire.}
    \item \textbf{Piège n°2 :} \zh{妇} (femme) vs \zh{扫} (balayer) — même structure gauche-droite, radical différent (\zh{女} vs \zh{扌}).
    \item \textbf{Méthode :} Décomposer → Écrire en l'air (ordre traits) → Copier 5× (quadrillage) → Dire phrase → Réviser J+1/J+3.
\end{itemize}
\end{aretenir}

\coursec{Structure du texte et connecteurs}

Après avoir maîtrisé l'écriture des caractères clés du thème (\emph{女, 平, 等, 权, 家, 务, 领, 导, 担, 传, 统}…), nous abordons maintenant l'architecture du texte. En expression écrite, connaître le vocabulaire ne suffit pas : il faut savoir \textbf{organiser ses idées} avec des connecteurs logiques et temporels précis. Le chinois structure l'information différemment du français (place du complément de temps, coexistence \emph{因为…所以…}, absence de conjonction de subordination pour « que » après \emph{觉得}). Cette section vous donne le « squelette » d'un bon texte de niveau HSK 3-4 sur la situation de la femme.

\courssub{Plan type d'un texte simple : les quatre étapes}

Tout texte argumentatif ou descriptif simple sur l'évolution de la condition féminine suit naturellement une progression en quatre temps. Observez ce schéma directeur :

\begin{popfigure}
\begin{tikzpicture}[
    node distance=0.8cm and 1.2cm,
    box/.style={rectangle, rounded corners=4pt, minimum width=3.2cm, minimum height=1.6cm, align=center, font=\small, draw=popdark, thick, top color=white, bottom color=popblueL},
    arrow/.style={-Stealth, thick, popdark}
]
\node[box, bottom color=poporangeL, align=center] (etape1){\textbf{1. Situation passée}\\\zh{以前 yǐqián}\\« Autrefois »};
\node[box, bottom color=popblueL, right=of etape1, align=center] (etape2){\textbf{2. Situation actuelle}\\\zh{现在 xiànzài}\\« Maintenant »};
\node[box, bottom color=popgreenL, right=of etape2, align=center] (etape3){\textbf{3. Problème / Contraste}\\\zh{但是 dànshì}\\« Mais / Cependant »};
\node[box, bottom color=poppurpleL, right=of etape3, align=center] (etape4){\textbf{4. Opinion / Proposition}\\\zh{我觉得 wǒ juéde / 应该 yīnggāi}\\« Je pense que / Il faut »};

\draw[arrow] (etape1.east) -- (etape2.west);
\draw[arrow] (etape2.east) -- (etape3.west);
\draw[arrow] (etape3.east) -- (etape4.west);

\node[below=0.4cm of etape1, font=\tiny\itshape, text=popmuted] {Ex: \zh{以前女人不上学}};
\node[below=0.4cm of etape2, font=\tiny\itshape, text=popmuted] {Ex: \zh{现在女人上大学}};
\node[below=0.4cm of etape3, font=\tiny\itshape, text=popmuted] {Ex: \zh{但是工资不平等}};
\node[below=0.4cm of etape4, font=\tiny\itshape, text=popmuted] {Ex: \zh{我觉得应该改变}};
\end{tikzpicture}
\legende{Frise chronologique et logique du plan type : Passé $\rightarrow$ Présent $\rightarrow$ Contraste $\rightarrow$ Opinion}
\end{popfigure}

\begin{propriete}[Plan type en 4 séquences]
Un paragraphe cohérent sur ce thème suit l'ordre immuable :
\begin{enumerate}
    \item \textbf{Contexte passé} : \zh{以前 yǐqián} + faits (accès limité à l'école, au travail, mariages précoces).
    \item \textbf{Constat présent} : \zh{现在 xiànzài} + évolutions positives (études, carrière, lois).
    \item \textbf{Nuance / Obstacle restant} : \zh{但是 dànshì} (ou \zh{可是 kěshì}) + réalité persistante (écart salarial, double journée, stéréotypes).
    \item \textbf{Position personnelle} : \zh{我觉得 wǒ juéde} / \zh{我认为 wǒ rènwéi} + \zh{应该 yīnggāi} + solution.
\end{enumerate}
\emph{Ne sautez jamais d'étape.} En chinois, la progression temporelle (昔/今) précède toujours l'argumentation.
\end{propriete}

\courssub{Connecteurs de temps : situer les faits}

Le chinois place le repère temporel \textbf{en tout début de phrase} (Sujet + Temps + Verbe), contrairement au français qui le met souvent à la fin.

\begin{center}
\begin{tabularx}{\linewidth}{|X|X|X|X|}
\hline
\rowcolor{popblueL} \textbf{Caractères} & \textbf{Pinyin} & \textbf{Nature} & \textbf{Traduction} \\
\hline
以前 & yǐqián & nom / adv. & autrefois, avant, il y a longtemps \\
\hline
现在 & xiànzài & nom / adv. & maintenant, actuellement \\
\hline
以后 & yǐhòu & nom / adv. & après, plus tard, à l'avenir \\
\hline
过去 & guòqù & nom / adv. & le passé, autrefois (plus formel) \\
\hline
如今 & rújīn & adv. & de nos jours, aujourd'hui (écrit, soutenu) \\
\hline
\end{tabularx}
\end{center}

\begin{exemplebox}[Position du marqueur temporel]
\begin{itemize}
    \item \zh{以前，女人不能上大学。} (Yǐqián, nǚrén bù néng shàng dàxué.) « Autrefois, les femmes ne pouvaient pas aller à l'université. »
    \item \zh{现在，她们当医生、工程师。} (Xiànzài, tāmen dāng yīshēng, gōngchéngshī.) « Maintenant, elles sont médecins, ingénieures. »
    \item \zh{以后，男女应该平等分担家务。} (Yǐhòu, nánnǚ yīnggāi píngděng fēndān jiāwù.) « À l'avenir, hommes et femmes devraient partager les tâches ménagères équitablement. »
\end{itemize}
\emph{Règle :} Le marqueur temporel précède le sujet ou s'intercale Sujet + Temps. \zh{女人以前不上学} est correct ; \zh{女人不上学以前} est \textbf{faux}.
\end{exemplebox}

\courssub{Connecteurs logiques : cause, conséquence, concession, addition}

C'est le cœur de l'argumentation. La structure \zh{因为…所以…} est la plus fréquente et la plus piège pour les francophones.

\begin{propriete}[La structure \zh{因为 yīnwèi … 所以 suǒyǐ} — Coexistence obligatoire]
En français, « Parce que P, Q » ou « P, donc Q » sont exclusifs (redondance). En chinois standard, \zh{因为} (cause) et \zh{所以} (conséquence) \textbf{s'emploient ensemble dans la même phrase complexe} pour marquer explicitement le lien logique. C'est la norme à l'écrit.
\begin{center}
\begin{tabularx}{\linewidth}{|X|X|X|}
\hline
\rowcolor{popblueL} \textbf{Structure} & \textbf{Exemple (caractères + pinyin)} & \textbf{Traduction} \\
\hline
Suj + \zh{因为} Cause , \zh{所以} Conséquence . & \zh{因为女人也工作，所以家务应该一起做。} (Yīnwèi nǚrén yě gōngzuò, suǒyǐ jiāwù yīnggāi yīqǐ zuò.) & Parce que les femmes travaillent aussi, les tâches ménagères devraient être faites ensemble. \\
\hline
\zh{因为} Cause , \zh{所以} Suj + Conséquence . & \zh{因为社会进步了，所以女孩子都上学了。} (Yīnwèi shèhuì jìnbù le, suǒyǐ nǚháizi dōu shàngxué le.) & Comme la société a progressé, toutes les filles vont à l'école. \\
\hline
\end{tabularx}
\end{center}
\emph{Variante formelle (écrite) :} \zh{由于 yóuyú …, 因此 yīncǐ …} (plus soutenu que \zh{因为…所以…}).
\end{propriete}

\begin{attention}[Piège n°1 : Traduction mot à mot de « parce que »]
Ne dites jamais : \zh{因为她累了，她不做饭。} (sans \zh{所以}).
Ne traduisez pas « Elle ne cuisine pas \textbf{parce qu'}elle est fatiguée » par \zh{她不做饭因为她累了。} (ordre français).
\emph{Bonne forme :} \zh{因为她累了，所以她不做饭。} ou \zh{她因为累了，所以不做饭。} (Sujet + \zh{因为}… + \zh{所以}…).
\end{attention}

\begin{center}
\begin{tabularx}{\linewidth}{|X|X|X|X|}
\hline
\rowcolor{popblueL} \textbf{Connecteur} & \textbf{Pinyin} & \textbf{Fonction} & \textbf{Exemple court} \\
\hline
但是 & dànshì & Contraste, opposition (« mais ») & \zh{她工作忙，但是很开心。} \\
\hline
可是 & kěshì & Contraste (plus oral, plus faible) & \zh{他想帮忙，可是不会做饭。} \\
\hline
而且 & érqiě & Addition, progression (« de plus », « en plus ») & \zh{她会做饭，而且做得好吃。} \\
\hline
虽然…但是… & suīrán… dànshì… & Concession (« bien que… pourtant ») & \zh{虽然她很累，但是她坚持学习。} \\
\hline
\end{tabularx}
\end{center}

\begin{exemplebox}[Articulation \zh{因为…所以…} + \zh{但是} + \zh{而且}]
\zh{因为受教育机会多了，所以现在女人当领导也不奇怪。但是，工资还是不平等，而且家务全靠女人。我觉得这不对。}
(Yīnwèi shòu jiàoyù jīhuì duō le, suǒyǐ xiànzài nǚrén dāng lǐngdǎo yě bù qíguài. Dànshì, gōngzī hái shì bù píngděng, érqiě jiāwù quán kào nǚrén. Wǒ juéde zhè bù duì.)
« Parce que les opportunités d'éducation ont augmenté, ce n'est plus étrange que des femmes soient dirigeantes. Mais les salaires ne sont toujours pas égaux, et en plus les tâches ménagères reposent entièrement sur les femmes. Je trouve cela injuste. »
\end{exemplebox}

\courssub{Connecteurs d'énumération : structurer l'argumentaire}

Pour lister des arguments, des exemples ou des étapes, on utilise la triade classique. Ils se placent en \textbf{début de proposition}, suivis d'une virgule.

\begin{center}
\begin{tabularx}{\linewidth}{|X|X|X|X|}
\hline
\rowcolor{popblueL} \textbf{Caractères} & \textbf{Pinyin} & \textbf{Nature} & \textbf{Traduction} \\
\hline
首先 & shǒuxiān & adv. & d'abord, en premier lieu \\
\hline
然后 & ránhòu & adv. & ensuite, puis \\
\hline
最后 & zuìhòu & adv. & enfin, en dernier lieu \\
\hline
另外 & lìngwài & adv. & par ailleurs, de plus \\
\hline
\end{tabularx}
\end{center}

\begin{exemplebox}[Énumération d'arguments — Modèle réutilisable]
\zh{首先，女人可以上学；然后，她们可以工作；最后，她们也可以当领导。}
(Shǒuxiān, nǚrén kěyǐ shàngxué; ránhòu, tāmen kěyǐ gōngzuò; zuìhòu, tāmen yě kěyǐ dāng lǐngdǎo.)
« D'abord, les femmes peuvent étudier ; ensuite, elles peuvent travailler ; enfin, elles peuvent aussi devenir dirigeantes. »
\emph{Note :} Le point-virgule ; (chinois \zh{；}) sépare les items longs ; la virgule , (chinois \zh{,}) pour les items courts.
\end{exemplebox}

\courssub{Expression de l'opinion et formulation de propositions}

Pour conclure votre texte, vous devez prendre position. Trois structures indispensables :

\begin{definition}[Verbes d'opinion + Proposition]
\begin{itemize}
    \item \zh{我觉得 wǒ juéde} : « Je pense que / J'ai l'impression que » (subjectif, courant).
    \item \zh{我认为 wǒ rènwéi} : « Je considère que / J'estime que » (plus formel, argumenté).
    \item \zh{应该 yīnggāi} : « devoir / il faut que » (obligation morale, conseil). Place : Suj + \zh{应该} + V.
\end{itemize}
\end{definition}

\begin{attention}[Piège n°2 : « Je pense que oui / non »]
En français, on répond « Je pense que oui ». En chinois, \zh{我觉得} \textbf{ne s'utilise pas seul} avec « oui/non ».
\begin{itemize}
    \item \textbf{Faux :} \zh{你觉得对吗？ 我觉得是。} (Wǒ juéde shì.)
    \item \textbf{Vrai :} \zh{你觉得对吗？ 我觉得是这样。} (Wǒ juéde shì zhèyàng.) ou \zh{我觉得对。} (Wǒ juéde duì.)
    \item \textbf{Vrai (opinion complète) :} \zh{我觉得女人应该有平等权利。} (Wǒ juéde nǚrén yīnggāi yǒu píngděng quánlì.)
\end{itemize}
\end{attention}

\begin{exemplebox}[Opinion + Justification (intégration \zh{因为…所以…})]
\zh{我认为，因为男女平等是基本权利，所以社会应该给女人更多机会。}
(Wǒ rènwéi, yīnwèi nánnǚ píngděng shì jīběn quánlì, suǒyǐ shèhuì yīnggāi gěi nǚrén gèng duō jīhuì.)
« Je considère que, parce que l'égalité hommes-femmes est un droit fondamental, la société devrait donner plus d'opportunités aux femmes. »
\end{exemplebox}

\courssub{Texte modèle commenté : « La femme d'aujourd'hui »}

Voici un texte original (environ 260 caractères) respectant le plan en 4 étapes et utilisant tous les connecteurs vus. Lisez-le à voix haute, repérez les connecteurs (surlignés mentalement), puis consultez le glossaire.

\begin{textebox}[Texte modèle : 今天的女人 Jīntiān de nǚrén]
以前，在中国和喀麦隆，女人地位很低。因为传统思想重男轻女，所以女孩子不上学，早早结婚，只做家务、带孩子。\\
现在，情况变了。首先，法律保护女人的权利；然后，学校向女学生敞开大门；最后，女人可以当医生、老师、部长，甚至当总统。而且，越来越多的男人开始做家务。\\
但是，问题还没有完全解决。虽然女人也工作，但是工资往往比男人少。而且，家务和照顾孩子，主要还是女人做。这不公平。\\
我觉得，男女平等不只是一句口号。因为家庭是社会的细胞，所以男女应该平等分担家务，共同决定大事。我认为，教育从小告诉孩子：男女一样重要，这才是根本。
\end{textebox}

\begin{center}
\begin{tabularx}{\linewidth}{|X|X|X|X|}
\hline
\rowcolor{popblueL} \textbf{Caractères} & \textbf{Pinyin} & \textbf{Nature} & \textbf{Traduction} \\
\hline
地位 & dìwèi & nom & statut, position sociale \\
\hline
传统思想 & chuántǒng sīxiǎng & nom & pensée traditionnelle \\
\hline
重男轻女 & zhòng nán qīng nǚ & idiome & privilégier les garçons, dévaloriser les filles \\
\hline
法律 & fǎlǜ & nom & loi \\
\hline
权利 & quánlì & nom & droit \\
\hline
敞开大门 & chǎngkāi dàmén & verbe + obj & ouvrir grand les portes (fig.) \\
\hline
甚至 & shènzhì & adv. & même, jusqu'à \\
\hline
完全 & wánquán & adv. & complètement, entièrement \\
\hline
解决 & jiějué & verbe & résoudre \\
\hline
往往 & wǎngwǎng & adv. & souvent, d'ordinaire \\
\hline
公平 & gōngpíng & adj. & juste, équitable \\
\hline
口号 & kǒuhào & nom & slogan \\
\hline
细胞 & xìbāo & nom & cellule (biol. / fig. unité de base) \\
\hline
分担 & fēndān & verbe & partager (une charge) \\
\hline
共同 & gòngtóng & adj. & commun, ensemble \\
\hline
根本 & gēnběn & nom/adj & fondement, fondamental \\
\hline
\end{tabularx}
\end{center}

\begin{savaistu}[Femmes au Cameroun et en Chine : regards croisés]
\begin{itemize}
    \item \textbf{Cameroun :} La loi (Constitution 1996, Code de la famille 2023 en révision) garantit l'égalité. Progrès : scolarisation des filles (taux brut primaire \textgreater 100\%), femmes ministres, députés, cheffes d'entreprise. Défis persistants : mariages précoces (régions Nord/Extrême-Nord), accès à la terre, violences basées sur le genre, double journée de travail.
    \item \textbf{Chine :} Slogan maoïste « Les femmes tiennent la moitié du ciel » (\zh{妇女能顶半边天}). Réforme agraire 1950 (Loi sur le mariage) : abolition mariages forcés, liberté divorce. Aujourd'hui : taux d'activité féminin parmi les plus élevés OCDE (\textasciitilde 61\%), forte présence dans STEM. Défis : écart salarial (\textasciitilde 20\%), « plafond de verre » (\zh{玻璃天花板}), pression sociale \zh{剩女 shèngnǔ} (« femmes restantes » non mariées à 27+ ans), charge domestique inégale.
    \item \textbf{Point commun :} La loi a changé plus vite que les mentalités. L'éducation (\zh{教育}) reste le levier principal cité dans les deux pays.
\end{itemize}

\textbf{Questions de compréhension / analyse (à faire oralement ou par écrit) :}
\begin{enumerate}
    \item Repérez les 4 étapes du plan (marqueurs \zh{以前}, \zh{现在}, \zh{但是}, \zh{我觉得/我认为}).
    \item Combien de fois la structure \zh{因为…所以…} apparaît-elle ? Soulignez-les.
    \item Quel connecteur d'énumération introduit les 3 progrès majeurs (école, travail, pouvoir) ?
    \item Traduisez la dernière phrase : \zh{我认为，教育从小告诉孩子：男女一样重要，这才是根本。}
\end{enumerate}

\end{savaistu}

\courssub{Méthode : Construire un paragraphe cohérent (étape par étape)}

\begin{methode}[Rédiger un paragraphe « Situation de la femme » en 5 minutes]
\begin{enumerate}
    \item \textbf{Poser le repère temporel :} Choisissez \zh{以前} ou \zh{过去} pour le passé ; \zh{现在} ou \zh{如今} pour le présent. Placez-le \textbf{en tout début de phrase 1}.
    \item \textbf{Énoncer 2 faits simples (SVO) :} Un fait passé négatif, un fait présent positif. Utilisez \zh{因为…所以…} pour lier cause (tradition/loi) et conséquence (études/travail).
    \item \textbf{Introduire la nuance :} Commencez une nouvelle phrase par \zh{但是} (ou \zh{可是}). Ajoutez un fait chiffré ou observable (salaire, tâches ménagères, plafond de verre).
    \item \textbf{Enchaîner avec \zh{而且} :} Ajoutez un second obstacle pour montrer l'ampleur.
    \item \textbf{Conclure par votre avis :} \zh{我觉得 / 我认为} + \zh{应该} + solution concrète (éducation, loi, partage tâches). Terminez par \zh{这才是根本} (c'est là le fondement) ou \zh{这样社会才会更好} (ainsi la société ira mieux).
\end{enumerate}
\textbf{Astuce :} Écrivez d'abord le squelette en pinyin/connecteurs, puis habillez de vocabulaire précis.
\end{methode}

\courssub{Thème et version guidés (Exercice résolu)}

\begin{exoresolu}[Thème : Du français au chinois — « Situation de la femme au Cameroun »]
\textbf{Énoncé :} Traduisez le paragraphe suivant en chinois (caractères + pinyin). Utilisez le plan 4 étapes et les connecteurs vus.
« Autrefois, au Cameroun, les filles ne allaient pas à l'école parce que les parents pensaient que le mariage était plus important. Maintenant, la situation a changé. D'abord, la loi protège les droits des filles. Ensuite, elles vont à l'université. Mais, le salaire des femmes est encore plus bas que celui des hommes. De plus, elles font presque toutes les tâches ménagères. Je pense que les hommes devraient partager le travail à la maison. »

\tcblower
\textbf{Solution détaillée :}

\textbf{Étape 1 (Passé + Cause) :}
\zh{以前，在喀麦隆，因为父母觉得结婚更重要，所以女孩子不上学。}
(Yǐqián, zài Kāmàilóng, yīnwèi fùmǔ juéde jiéhūn gèng zhòngyào, suǒyǐ nǚháizi bù shàngxué.)
\emph{Note :} \zh{在喀麦隆} (au Cameroun) placé après \zh{以前}. \zh{觉得} (penser que) introduit la cause. \zh{更重要} (plus important).

\textbf{Étape 2 (Présent + Énumération) :}
\zh{现在，情况变了。首先，法律保护女孩子的权利；然后，她们上大学。}
(Xiànzài, qíngkuàng biàn le. Shǒuxiān, fǎlǜ bǎohù nǚháizi de quánlì; ránhòu, tāmen shàng dàxué.)
\emph{Note :} \zh{情况变了} (la situation a changé) = transition standard. \zh{上大学} (aller à l'université).

\textbf{Étape 3 (Contraste + Addition) :}
\zh{但是，女人的工资还是比男人少。而且，家务几乎全是女人做。}
(Dànshì, nǚrén de gōngzī hái shì bǐ nánrén shǎo. Érqiě, jiāwù jīhū quán shì nǚrén zuò.)
\emph{Note :} Comparatif \zh{A 比 B Adj} (voir leçon grammaire). \zh{几乎 jīhū} (presque). \zh{全是} (entièrement fait par).

\textbf{Étape 4 (Opinion + Proposition) :}
\zh{我觉得男人应该分担家务。}
(Wǒ juéde nánrén yīnggāi fēndān jiāwù.)
\emph{Variante plus formelle :} \zh{我认为，男女应该平等分担家务，这才是根本。}

\textbf{Texte complet assemblé :}
\zh{以前，在喀麦隆，因为父母觉得结婚更重要，所以女孩子不上学。现在，情况变了。首先，法律保护女孩子的权利；然后，她们上大学。但是，女人的工资还是比男人少。而且，家务几乎全是女人做。我觉得男人应该分担家务。}
\end{exoresolu}

\courssub{Exercices d'application}

\begin{exercice}[Mise en ordre et complétion]
\textbf{A. Remettez les phrases dans l'ordre logique (1 à 6) pour former un paragraphe cohérent.}
\begin{enumerate}
    \item \zh{最后，我觉得社会应该改变这种观念。}
    \item \zh{现在，很多女人上大学，当领导。}
    \item \zh{因为传统思想，所以以前女人不工作。}
    \item \zh{但是，她们回家还要做所有家务。}
    \item \zh{而且，工资往往不平等。}
    \item \zh{以前，女人地位很低。}
\end{enumerate}

\textbf{B. Complétez avec le connecteur approprié : \zh{因为 / 所以 / 但是 / 而且 / 首先 / 然后 / 最后 / 我觉得 / 应该}.}
1. \zh{\_\_\_\_\_\_ 女人也上班，\_\_\_\_\_\_ 家务男人也要做。}
2. \zh{\_\_\_\_\_\_，法律给了权利；\_\_\_\_\_\_，女孩上了学；\_\_\_\_\_\_，女人当了部长。}
3. \zh{\_\_\_\_\_\_ 她很忙，\_\_\_\_\_\_ 她坚持学中文。}
4. \zh{\_\_\_\_\_\_ 男女平等是基本人权。}

\textbf{C. Version (Chinois $\rightarrow$ Français) : Traduisez fidèlement.}
\zh{虽然现在女大学生比男大学生多，但是高层领导里女人很少。我觉得这是因为社会对女人的要求太高了。男人应该多承担家庭责任，这样女人才有机会发展事业。}
\end{exercice}

\begin{corrige}[Exercice : Structure du texte et connecteurs]
\textbf{A. Ordre logique :} 6 $\rightarrow$ 3 $\rightarrow$ 2 $\rightarrow$ 4 $\rightarrow$ 5 $\rightarrow$ 1.
\textit{Parcours :} Passé (6) + Cause passé (3) $\rightarrow$ Présent (2) $\rightarrow$ Contraste 1 (4) $\rightarrow$ Contraste 2 / Addition (5) $\rightarrow$ Opinion finale (1).

\textbf{B. Complétion :}
1. \zh{\textbf{因为} 女人也上班，\textbf{所以} 家务男人也要做。} (Cause $\rightarrow$ Conséquence)
2. \zh{\textbf{首先}，法律给了权利；\textbf{然后}，女孩上了学；\textbf{最后}，女人当了部长。} (Énumération chronologique/argumentative)
3. \zh{\textbf{虽然} 她很忙，\textbf{但是} 她坚持学中文。} (Concession — \emph{attention : bien que = 虽然, mais = 但是})
4. \zh{\textbf{我觉得} 男女平等是基本人权。} (Opinion) \quad \textit{ou} \zh{男女平等\textbf{应该是}基本人权。} (Devoir moral — structure préférée à \zh{应该 男女平等...}).

\textbf{C. Version :}
« Bien qu'actuellement il y ait plus d'étudiantes que d'étudiants, il y a très peu de femmes dans les hautes sphères du pouvoir. Je pense que c'est parce que la société a des exigences trop élevées envers les femmes. Les hommes devraient assumer davantage de responsabilités familiales, ainsi seulement les femmes auront l'opportunité de développer leur carrière. »
\emph{Points de vocabulaire :} \zh{高层领导} (hauts dirigeants), \zh{要求太高} (exigences trop fortes), \zh{承担责任} (assumer des responsabilités), \zh{发展事业} (développer sa carrière), \zh{才 cái} (marque la condition nécessaire : « seulement ainsi / alors seulement »).
\end{corrige}

\begin{aretenir}[Check-list « Expression écrite : Situation de la femme »]
Avant de rendre votre copie, vérifiez :
\begin{itemize}
    \item [ ] Plan respecté : \zh{以前} $\rightarrow$ \zh{现在} $\rightarrow$ \zh{但是} $\rightarrow$ \zh{我觉得/我认为}.
    \item [ ] Au moins une structure \zh{因为…所以…} correcte (les deux présents).
    \item [ ] Connecteurs d'énumération (\zh{首先/然后/最后}) pour lister arguments.
    \item [ ] Opinion marquée par \zh{应该} + verbe d'action (partager, changer, éduquer).
    \item [ ] Caractères clés du thème écrits correctement (\zh{女, 平, 权, 务, 领, 导, 担, 传, 统}…).
    \item [ ] Ponctuation chinoise (， 。 ；) et pas d'espaces entre caractères.
\end{itemize}
\end{aretenir}

\coursec{Thème et version guidés}

Après avoir analysé la structure du texte modèle et les connecteurs logiques dans la section précédente, nous abordons maintenant le cœur de l'épreuve écrite du Baccalauréat : la traduction. Passer du français au chinois (thème) et du chinois au français (version) ne s'improvise pas ; cela requiert une méthode rigoureuse, une connaissance précise du vocabulaire du module et une maîtrise de l'ordre des mots, si différent de notre langue. C'est ici que se joue la justesse de votre expression écrite.

\courssub{Méthode du thème : du français vers le chinois}

Le thème (français $\rightarrow$ chinois) est souvent perçu comme plus difficile car il exige une production active de caractères et de structures. Suivez ces quatre étapes systématiquement.

\begin{methode}[Méthode du thème (français $\rightarrow$ chinois)]
\begin{enumerate}
    \item \textbf{Analyser la phrase française} : identifiez le Sujet (S), le Verbe (V), les Compléments (COI, COD, Circonstanciels de temps, lieu, manière). Repérez la négation, le temps (passé, présent, futur), la modalité (devoir, vouloir, pouvoir).
    \item \textbf{Choisir le vocabulaire} : piochez dans le glossaire de la leçon (section 2) et vos fiches de révision. Vérifiez la nature du mot (verbe, nom, adjectif predicatif) et son classificateur si c'est un nom concret.
    \item \textbf{Construire la phrase chinoise} : appliquez l'ordre canonique \cle{Sujet + Temps + Lieu + Manière + Verbe + Complément d'objet}. \textbf{Règle d'or} : le complément de temps (以前, 现在, 每天, 昨天) se place \emph{avant} le verbe, jamais à la fin comme en français.
    \item \textbf{Vérifier et écrire} : contrôlez l'ordre des traits pour chaque caractère, les tons (pinyin), la ponctuation chinoise (， 。) et l'accord des adjectifs predicatifs (avec \zh{很}, \zh{是}... \zh{的}).
\end{enumerate}
\end{methode}

\begin{attention}[Position du complément de temps : piège n°1 des francophones]
En français, on dit : « \emph{Les femmes travaillent maintenant.} » (Temps à la fin).
En chinois, c'est \textbf{interdit} : *\zh{女人工作现在} est faux.
La structure correcte est : \zh{女人现在工作} (Sujet + Temps + Verbe).
De même : « \emph{Autrefois, elles ne travaillaient pas.} » $\rightarrow$ \zh{以前，女人不工作。} (Temps en tête de phrase ou après le sujet).
\end{attention}

Appliquons la méthode aux trois phrases types du programme.

\begin{exemplebox}[Thème 1 : « Autrefois, les femmes ne travaillaient pas. »]
\textbf{1. Analyse FR :} Sujet = « les femmes » ; Temps = « Autrefois » (passé) ; Négation = « ne ... pas » ; Verbe = « travailler ».
\textbf{2. Vocabulaire :} \zh{女人} (nǚrén) ; \zh{以前} (yǐqián) ; \zh{不} (bù) ; \zh{工作} (gōngzuò).
\textbf{3. Construction CN :} Sujet (\zh{女人}) + Temps (\zh{以前}) + Négation (\zh{不}) + Verbe (\zh{工作}).
\textbf{4. Résultat :} \zh{以前，女人不工作。} (Yǐqián, nǚrén bù gōngzuò.)
\emph{Note :} La virgule après \zh{以前} marque la topicalisation du temps (thème rhématique).
\end{exemplebox}

\begin{exemplebox}[Thème 2 : « Elle doit s'occuper des enfants et faire le ménage. »]
\textbf{1. Analyse FR :} Sujet = « Elle » ; Modalité = « doit » (obligation) ; V1 = « s'occuper de » + COD « les enfants » ; Liaison = « et » (addition) ; V2 = « faire le ménage ».
\textbf{2. Vocabulaire :} \zh{她} (tā) ; \zh{要} (yào) ; \zh{照顾} (zhàogù) ; \zh{孩子} (háizi) ; \zh{还要} (hái yào) / \zh{又} (yòu) ; \zh{做家务} (zuò jiāwù).
\textbf{3. Construction CN :} Sujet (\zh{她}) + Verbe modal (\zh{要}) + V1 (\zh{照顾}) + Obj1 (\zh{孩子}) + Conjonction addition (\zh{还要}) + V2 (\zh{做家务}).
\textbf{4. Résultat :} \zh{她要照顾孩子，还要做家务。} (Tā yào zhàogù háizi, hái yào zuò jiāwù.)
\end{exemplebox}

\begin{attention}[Le verbe modal \cle{要} : plusieurs sens selon le contexte]
\zh{要} est un « caméléon » grammatical. Dans nos thèmes :
\begin{itemize}
    \item \zh{我要去北京。} (Wǒ yào qù Běijīng.) = « Je \textbf{veux} aller à Pékin. » (Volonté, désir).
    \item \zh{明天要下雨。} (Míngtiān yào xiàyǔ.) = « Demain, il \textbf{va} pleuvoir. » (Futur proche, prédiction).
    \item \zh{她要照顾孩子。} (Tā yào zhàogù háizi.) = « Elle \textbf{doit} s'occuper des enfants. » (Obligation, devoir moral/social).
\end{itemize}
Au Bac, si le contexte français impose « doit / avoir l'obligation de », \zh{要} est le bon choix (registre courant). \zh{必须} (bìxū) est plus fort (« impératif absolu »).
\end{attention}

\begin{exemplebox}[Thème 3 : « Je pense que l'homme et la femme sont égaux. »]
\textbf{1. Analyse FR :} Sujet = « Je » ; Verbe opinion = « pense que » ; Proposition subordonnée : Sujet = « l'homme et la femme » (sujet complexe) ; Verbe d'état = « sont » ; Attribut = « égaux ».
\textbf{2. Vocabulaire :} \zh{我} (wǒ) ; \zh{觉得} (juéde) ; \zh{男人} (nánrén) ; \zh{和} (hé) ; \zh{女人} (nǚrén) ; \zh{是} (shì) ; \zh{平等} (píngděng) ; \zh{的} (de).
\textbf{3. Construction CN :} Proposition principale : \zh{我觉得} (Je pense que). Proposition subordonnée (devient objet de \zh{觉得}) : Sujet \zh{男人和女人} + Verbe \zh{是} + Adj predicatif \zh{平等的}.
\textbf{4. Résultat :} \zh{我觉得男人和女人是平等的。} (Wǒ juéde nánrén hé nǚrén shì píngděng de.)
\emph{Point de grammaire :} L'adjectif \zh{平等} devient attribut du sujet via \zh{是} + \zh{的} (structure nominale). On ne dit pas *\zh{他们平等}.
\end{exemplebox}

\begin{exoresolu}[Thème guidé n°1 : phrases à traduire]
\textbf{Énoncé :} Traduisez en chinois (caractères + pinyin).
\begin{enumerate}
    \item Autrefois, les hommes ne faisaient pas la cuisine.
    \item Aujourd'hui, ma mère a son propre travail.
    \item Je crois que les filles et les garçons sont les mêmes.
\end{enumerate}
\tcblower
\textbf{Solution détaillée :}
\begin{enumerate}
    \item \textbf{Analyse :} Temps=\zh{以前}, Sujet=\zh{男人}, Négation=\zh{不}, Verbe=\zh{做饭}.
    \zh{以前，男人不做饭。} (Yǐqián, nánrén bù zuòfàn.)
    \item \textbf{Analyse :} Temps=\zh{现在}, Sujet=\zh{我妈妈}, Verbe=\zh{有}, Objet=\zh{自己的工作}.
    \zh{现在，我妈妈有自己的工作。} (Xiànzài, wǒ māma yǒu zìjǐ de gōngzuò.)
    \item \textbf{Analyse :} Opinion=\zh{我觉得}, Sujet complexe=\zh{女孩子和男孩子}, Verbe=\zh{是}, Attribut=\zh{一样的}.
    \zh{我觉得女孩子和男孩子是一样的。} (Wǒ juéde nǚháizi hé nánháizi shì yíyàng de.)
\end{enumerate}
\end{exoresolu}

\courssub{Méthode de la version : du chinois vers le français}

La version (chinois $\rightarrow$ français) teste votre compréhension fine et votre capacité à produire un français correct, naturel, sans calquer la syntaxe chinoise.

\begin{methode}[Méthode de la version (chinois $\rightarrow$ français)]
\begin{enumerate}
    \item \textbf{Lecture globale} : lisez la phrase ou le texte entier sans dictionnaire pour saisir le sens général (qui ? quoi ? quand ?).
    \item \textbf{Identification lexicale et structurelle} : soulignez les mots connus. Repérez la structure S-V-C. Identifiez les \cle{mots-outils} (了, 的, 得, 地, 过, 着, 比, 把, 被) et les \cle{connecteurs} (因为...所以, 虽然...但是, 现在, 以前, 也, 还).
    \item \textbf{Traduction mot à mot (brouillon)} : écrivez le sens littéral de chaque mot dans l'ordre chinois.
    \item \textbf{Reformulation en français correct} : réordonnez selon la syntaxe française (SVO, temps à la fin, prépositions). Choisissez le bon temps verbal (passé composé, imparfait, présent). Vérifiez l'accord (genre, nombre).
\end{enumerate}
\end{methode}

Appliquons-la à la phrase type du programme.

\begin{exemplebox}[Version guidée : \zh{现在的女人跟以前不一样了。她们有自己的工作，也有自己的权利。}]
\textbf{Texte source :}
\zh{现在的女人跟以前不一样了。她们有自己的工作，也有自己的权利。}
(Xiànzài de nǚrén gēn yǐqián bù yíyàng le. Tāmen yǒu zìjǐ de gōngzuò, yě yǒu zìjǐ de quánlì.)

\textbf{Étape 1 -- Lecture globale :} Comparaison femmes maintenant / avant. Changement positif (travail, droits).
\textbf{Étape 2 -- Analyse détaillée :}
\begin{itemize}
    \item \zh{现在的} (xiànzài de) : « d'aujourd'hui / actuelles » (particule \zh{的} marquant l'adjectif).
    \item \zh{女人} (nǚrén) : « femmes » (sujet).
    \item \zh{跟...不一样} (gēn... bù yíyàng) : structure « ne pas être comme... / être différent de... ».
    \item \zh{以前} (yǐqián) : « avant / autrefois » (complément de comparaison).
    \item \zh{了} (le) : particule de changement d'état / nouvelle situation.
    \item \zh{她们} (tāmen) : « elles » (féminin pluriel, reprise du sujet).
    \item \zh{有} (yǒu) : « avoir / il y a ».
    \item \zh{自己的} (zìjǐ de) : « son propre / leur propre » (possessif réfléchi).
    \item \zh{工作} (gōngzuò) : « travail / emploi ».
    \item \zh{也} (yě) : « aussi / également ».
    \item \zh{权利} (quánlì) : « droits ».
\end{itemize}
\textbf{Étape 3 -- Brouillon littéral :} « Maintenant-de femmes avec avant pas pareil *le*. Elles avoir leur-propre travail, aussi avoir leur-propre droits. »
\textbf{Étape 4 -- Français naturel :}
\emph{« Les femmes d'aujourd'hui ne sont plus comme avant. Elles ont leur propre travail et aussi leurs propres droits. »}
$\rightarrow$ Version soignée : \textbf{« Les femmes d'aujourd'hui ne sont plus comme avant. Elles ont leur propre travail et leurs propres droits. »}
\end{exemplebox}

\begin{exoresolu}[Version guidée n°1]
\textbf{Énoncé :} Traduisez en français les phrases suivantes.
\begin{enumerate}
    \item \zh{我妈妈以前不工作，现在在医院工作。}
    \item \zh{男人和女人应该平等，因为他们有同样的权利。}
\end{enumerate}
\tcblower
\textbf{Solution détaillée :}
\begin{enumerate}
    \item \textbf{Analyse :} \zh{我妈妈} (Ma mère) / \zh{以前} (autrefois) / \zh{不工作} (ne travaillait pas) / \zh{现在} (maintenant) / \zh{在医院} (à l'hôpital) / \zh{工作} (travaille).
    \textbf{Structure :} Deux propositions opposées (autrefois vs maintenant). Connecteur implicite « mais ».
    \textbf{Traduction :} \emph{« Ma mère ne travaillait pas autrefois, mais maintenant elle travaille à l'hôpital. »} (Imparfait pour l'habitude passée, présent pour la situation actuelle).
    \item \textbf{Analyse :} \zh{男人和女人} (Les hommes et les femmes) / \zh{应该} (yīnggāi, doivent / devraient) / \zh{平等} (píngděng, égaux) / \zh{因为} (yīnwèi, parce que) / \zh{他们} (tāmen, ils) / \zh{有} (yǒu, ont) / \zh{同样的} (tóngyàng de, les mêmes) / \zh{权利} (quánlì, droits).
    \textbf{Traduction :} \emph{« Les hommes et les femmes devraient être égaux, car ils ont les mêmes droits. »} (Conditionnel de recommandation pour \zh{应该}, présent pour le fait général).
\end{enumerate}
\end{exoresolu}

\courssub{Pièges de traduction et points de vigilance}

Voici les trois écueils majeurs qui font perdre des points au Bac, même quand le vocabulaire est connu.

\begin{attention}[Piège 1 : \cle{跟...不一样} vs \cle{不像}]
\begin{itemize}
    \item \zh{跟...不一样} (gēn... bù yíyàng) : « être différent de / ne pas être comme... ». Structure : A \zh{跟} B \zh{不一样}. Ex: \zh{这件衣服跟那件不一样。} (Zhè jiàn yīfu gēn nà jiàn bù yíyàng.) — « Ce vêtement n'est pas comme celui-là. »
    \item \zh{不像} (bù xiàng) : « ne pas ressembler à » (souvent physique ou apparence) ou « contrairement à » (connecteur de phrase). Ex: \zh{他不像医生。} (Tā bù xiàng yīshēng.) — « Il ne ressemble pas à un médecin / Il n'a pas l'air médecin. »
    \item \textbf{Au Bac :} Pour exprimer « la situation a changé, ce n'est plus comme avant », utilisez \zh{跟以前不一样了}. N'écrivez pas *\zh{不像以前}.
\end{itemize}
\end{attention}

\begin{attention}[Piège 2 : \cle{自己的} -- Le possessif réfléchi universel]
En français, le possessif s'accorde avec \textbf{l'objet possédé} : « \emph{son} travail » (masc. sg), « \emph{sa} maison » (fém. sg), « \emph{leurs} droits » (pluriel).
En chinois, \zh{自己的} (zìjǐ de) est \textbf{invariant}. Il renvoie toujours au \textbf{sujet de la phrase}.
\begin{center}
\begin{tabularx}{\linewidth}{|X|X|X|}\hline
\rowcolor{popblueL} Sujet & Chinois & Français \\ \hline
\zh{我} (wǒ, Je) & \zh{我的书 / 我自己的书} & mon livre / \textbf{mon propre} livre \\ \hline
\zh{她} (tā, Elle) & \zh{她自己的工作} & \textbf{son propre} travail \\ \hline
\zh{他们} (tāmen, Ils) & \zh{他们自己的权利} & \textbf{leurs propres} droits \\ \hline
\zh{女人} (nǚrén, Les femmes) & \zh{她们自己的权利} & \textbf{leurs propres} droits \\ \hline
\end{tabularx}
\end{center}
\textbf{Erreur fréquente :} Traduire \zh{她们有自己的权利} par « Elles ont \emph{son} propre droit » (calque sur le singulier chinois \zh{权利} non marqué au pluriel). \textbf{Pensez à l'accord français !}
\end{attention}

\begin{attention}[Piège 3 : \cle{要} -- Devoir vs Vouloir (Rappel)]
Ne confondez pas l'obligation (\zh{要} = devoir) et le désir (\zh{想} = vouloir).
\begin{itemize}
    \item \zh{我要去上班。} (Wǒ yào qù shàngbān. Contexte : obligation horaire) $\rightarrow$ « Je \textbf{dois} aller au travail. »
    \item \zh{我想去旅游。} (Wǒ xiǎng qù lǚyóu. Contexte : envie) $\rightarrow$ « Je \textbf{veux} / \textbf{aimerais} voyager. »
\end{itemize}
Dans le thème « Elle doit s'occuper des enfants », c'est une contrainte sociale $\rightarrow$ \zh{要}. Si c'était un choix personnel « Elle veut s'occuper... », on écrirait \zh{她想照顾孩子}.
\end{attention}

\begin{exemplebox}[Exemples traduits supplémentaires pour l'entraînement]
\begin{enumerate}
    \item \zh{女人跟男人一样，都有自己的权利。} (Nǚrén gēn nánrén yíyàng, dōu yǒu zìjǐ de quánlì.) \\
    \emph{« Les femmes, comme les hommes, ont toutes leurs propres droits. »} (Structure \zh{跟...一样} = « comme... » ; \zh{都} = « toutes / tous » distribuatif).
    \item \zh{以前女人没有权利，现在法律保护她们。} (Yǐqián nǚrén méiyǒu quánlì, xiànzài fǎlǜ bǎohù tāmen.) \\
    \emph{« Autrefois les femmes n'avaient pas de droits, maintenant la loi les protège. »} (\zh{没有} = n'avoir pas ; \zh{法律} = loi ; \zh{保护} = protéger).
    \item \zh{做家务不只是女人的事，男人也应该做。} (Zuò jiāwù bù zhǐ shì nǚrén de shì, nánrén yě yīnggāi zuò.) \\
    \emph{« Faire le ménage n'est pas seulement l'affaire des femmes, les hommes devraient le faire aussi. »} (\zh{不只是} = n'est pas seulement ; \zh{应该} = devoir moral).
\end{enumerate}
\end{exemplebox}

\courssub{Entraînement : texte à traduire (version / thème)}

Le texte ci-dessous réutilise le vocabulaire et les structures de la leçon. Il sert de support d'entraînement complet (version chinois $\rightarrow$ français, puis thème français $\rightarrow$ chinois en masquant l'original).

\begin{textebox}[La place des femmes dans la société moderne]
\zh{现在的社会跟以前很不一样。以前，大多数女人没有机会上学，也没有自己的工作。她们每天在家里做家务，照顾孩子，照顾老人，很辛苦。} \\
\zh{现在，女孩子也可以上大学，找到好工作。我的姐姐是医生，她在大医院工作，工资很高。她结婚了，丈夫也做家务，带孩子。他们平分家务，日子过得很幸福。} \\
\zh{我觉得男人和女人应该平等。女人有权利选择自己的生活，男人也有责任分担家务。这样，家庭才和谐，社会才进步。}

\textbf{Glossaire (nouveaux mots / rappel) :}
\begin{center}\begin{tabularx}{\linewidth}{|X|X|X|X|}\hline
\rowcolor{popblueL} Caractères & Pinyin & Nature & Traduction \\ \hline
社会 & shèhuì & n. & société \\ \hline
大多数 & dàduōshù & pron. & la grande majorité \\ \hline
机会 & jīhuì & n. & occasion, opportunité \\ \hline
上学 & shàngxué & v. & aller à l'école, étudier \\ \hline
辛苦 & xīnkǔ & adj. & dur, pénible, fatigant \\ \hline
大学 & dàxué & n. & université \\ \hline
医生 & yīshēng & n. & médecin \\ \hline
工资 & gōngzī & n. & salaire \\ \hline
丈夫 & zhàngfu & n. & mari \\ \hline
带孩子 & dài háizi & v. & s'occuper des enfants (porter/élever) \\ \hline
平分 & píngfēn & v. & partager équitablement \\ \hline
日子 & rìzi & n. & vie, quotidien, jours \\ \hline
幸福 & xìngfú & adj. & heureux (vie, mariage) \\ \hline
权利 & quánlì & n. & droit \\ \hline
选择 & xuǎnzé & v. & choisir \\ \hline
生活 & shēnghuó & n. & vie, vie quotidienne \\ \hline
责任 & zérèn & n. & responsabilité, devoir \\ \hline
分担 & fēndān & v. & partager (une charge) \\ \hline
和谐 & héxié & adj. & harmonieux \\ \hline
进步 & jìnbù & v./n. & progresser / progrès \\ \hline
\end{tabularx}\end{center}

\textbf{Questions de compréhension / expression (à l'oral ou par écrit) :}
\begin{enumerate}
    \item Selon le texte, quelles étaient les occupations principales des femmes « autrefois » ? (Citez 3 verbes).
    \item Quelle est la profession de la sœur de l'auteur ? Où travaille-t-elle ?
    \item Comment le couple gère-t-il les tâches ménagères ? (Verbe clé : \zh{平分}).
    \item Traduisez la dernière phrase en français : \zh{这样，家庭才和谐，社会才进步。}
    \item \textbf{Thème inverse :} Re-traduisez le paragraphe 2 en chinois à partir de votre version française, sans regarder l'original.
\end{enumerate}
\end{textebox}

\begin{savaistu}[Le thème et la version au Baccalauréat]
À l'épreuve écrite du Bac (LV2 Chinois), l'exercice de traduction (thème + version) représente souvent \textbf{8 à 10 points sur 20}.
\begin{itemize}
    \item \textbf{Le thème (FR $\rightarrow$ CN)} porte sur 3 à 5 phrases isolées reprenant \textbf{exactement} les structures du programme (比, 把, 被, 了, 过, 要/想/应该, compléments de temps/lieu). La justesse des \textbf{caractères} (traits, radicaux) et de l'\textbf{ordre des mots} compte pour 50\% de la note de l'exercice. Un mot juste mal écrit = 0 point.
    \item \textbf{La version (CN $\rightarrow$ FR)} propose un court paragraphe (40-60 caractères) sur un thème du programme (famille, école, société, environnement). La note sanctionne la \textbf{qualité du français} (pas de « chinois-français » style « Les femmes d'avant pas pareil maintenant »). Évitez le mot-à-mot.
    \item \textbf{Astuce de pro :} Relisez vos thèmes en vous demandant : « Ai-je mis le temps avant le verbe ? Ai-je mis le 的 / 得 / 地 correct ? Le sujet est-il clair ? »
\end{itemize}
\end{savaistu}

\begin{popfigure}
\begin{tikzpicture}[
    node distance=1.1cm and 0.5cm,
    >=Stealth,
    box/.style={draw, rounded corners=6pt, thick, align=center, minimum width=2.8cm, minimum height=1.6cm, font=\small},
    arrow/.style={->, thick, shorten >=2pt, shorten <=2pt}
]
% Nœuds
\node[box, fill=popblueL, align=center] (step1){\textbf{1. LIRE}\\\emph{Globalement}\\(Sens général)};
\node[box, fill=poporangeL, right=of step1, align=center] (step2){\textbf{2. ANALYSER}\\\emph{S-V-C + Mots-outils}\\(Temps, Négation, Connecteurs)};
\node[box, fill=popgreenL, right=of step2, align=center] (step3){\textbf{3. TRADUIRE}\\\emph{Mot-à-mot $\rightarrow$ Réordonner}\\(Syntaxes FR / CN)};
\node[box, fill=poppinkL, right=of step3, align=center] (step4){\textbf{4. VÉRIFIER}\\\emph{Caractères, Tons, Accords}\\(Ponctuation, Français naturel)};

% Flèches principales
\draw[arrow, popblue] (step1.east) -- (step2.west) node[midway, above, font=\tiny] {Identifier};
\draw[arrow, poporange] (step2.east) -- (step3.west) node[midway, above, font=\tiny] {Structurer};
\draw[arrow, popgreen] (step3.east) -- (step4.west) node[midway, above, font=\tiny] {Contrôler};

% Boucle de retour (optionnel, style subtil)
\draw[arrow, popmuted, dashed, bend left=40] (step4.north) to node[midway, above, font=\tiny] {Si erreur} (step2.north);

% Titre au-dessus
\node[above=0.3cm of step2, font=\bfseries\large, text=popdark] {Organigramme : Les 4 étapes universelles de la traduction};
\end{tikzpicture}
\legende{Méthode unique pour le Thème et la Version : Lire $\rightarrow$ Analyser $\rightarrow$ Traduire $\rightarrow$ Vérifier.}
\end{popfigure}

\coursec{Production écrite et grille d'évaluation}

\courssub{Transition et objectif de la séance}

Après avoir travaillé le thème et la version guidés (section 5), vous disposez maintenant du vocabulaire, des structures grammaticales et des connecteurs nécessaires pour produire un texte autonome. L'épreuve écrite du Baccalauréat (série A) vous demandera souvent de rédiger un court paragraphe (5 à 8 phrases) sur un sujet de société. Le thème « la situation de la femme » est un grand classique : il permet de mobiliser le lexique de la famille, du travail, de l'éducation et des comparaisons, tout en exigeant une progression logique (passé $\rightarrow$ présent $\rightarrow$ problème $\rightarrow$ opinion). Cette section vous donne la méthode, la grille d'évaluation officielle et un entraînement complet pour réussir cet exercice le jour de l'examen.

\begin{methode}[Réussir la production écrite au Bac — écrire simple et juste]
\emph{Objectif :} produire un texte de 5 à 8 phrases correctes, cohérentes et notées sur 10 points.
\begin{enumerate}
\item \textbf{Analysez le sujet} (30 s) : identifiez le thème (femme), le pays (votre pays = Cameroun), la structure attendue (passé / présent / problème / opinion).
\item \textbf{Préparez un plan mental} (1 min) : 2 phrases sur le passé, 2 sur le présent, 1 sur le problème persistant, 1--2 sur votre opinion / espoir.
\item \textbf{Choisissez des phrases « clés en main »} que vous maîtrisez parfaitement (voir la \textbf{banque de phrases de démarrage} ci-dessous). N'inventez pas de structures complexes le jour J.
\item \textbf{Rédigez au brouillon} en caractères + pinyin (5 min). Vérifiez : ordre des mots (SVO, compléments avant le verbe), radicaux \cle{女} pour les mots féminins, connecteurs variés, \cle{了} de changement d'état.
\item \textbf{Recopiez proprement} sur la copie d'examen (2 min). Comptez les phrases (5--8). Relisez une dernière fois.
\end{enumerate}
\textbf{Règle d'or :} 6 phrases correctes valent mieux que 10 phrases fautives. La note porte sur la qualité, pas la quantité.
\end{methode}

\courssub{Consigne type Bac et plan de rédaction}

\begin{textebox}[Sujet type Baccalauréat]
Rédigez un texte de 5 à 8 phrases en chinois sur la situation de la femme dans votre pays.
\\ \emph{Consigne :} « 请用中文写一段 5 到 8 句话的短文，谈谈你国家妇女的情况。 » (Qǐng yòng Zhōngwén xiě yī duǎn 5 dào 8 jù huà de duǎnwén, tántán nǐ guójiā fùnǚ de qíngkuàng.)
\end{textebox}

\begin{propriete}[Plan type en 4 étapes (passé $\rightarrow$ présent $\rightarrow$ problème $\rightarrow$ opinion)]
Chaque étape = 1 à 2 phrases. Total visé : 6--7 phrases.
\begin{center}
\begin{tabularx}{\linewidth}{|X|X|X|}\hline
\rowcolor{popblueL} \textbf{Étape} & \textbf{Idée directrice} & \textbf{Connecteurs / structures utiles} \\ \hline
1. Passé & Situation ancienne : peu d'éducation, rôle domestique, dépendance & 以前 (yǐqián), 过去 (guòqù), 没有…机会 (méiyǒu… jīhuì), 只能 (zhǐnéng) \\ \hline
2. Présent & Progrès : scolarisation, travail, lois, statut qui monte & 现在 (xiànzài), 如今 (rújīn), 已经 (yǐjīng), 越来越…了 (yuè lái yuè… le) \\ \hline
3. Problème & Inégalités persistantes : salaire, charge mentale, violences & 但是 (dànshì), 不过 (bùguò), 还存在 (hái cúnzài), 依然 (yīrán) \\ \hline
4. Opinion & Égalité nécessaire, espoir, engagement personnel & 我觉得 (wǒ juéde), 我认为 (wǒ rènwéi), 应该 (yīnggāi), 希望 (xīwàng) \\ \hline
\end{tabularx}
\end{center}
\end{propriete}

\courssub{Banque de phrases de démarrage (à mémoriser et adapter)}

Ces 8 phrases « passe-partout » couvrent les 4 étapes. Apprenez-les par cœur avec leurs caractères, pinyin et traduction. Vous pourrez les combiner ou les modifier légèrement le jour de l'examen.

\begin{center}
\begin{tabularx}{\linewidth}{|X|X|X|X|}\hline
\rowcolor{popblueL} \textbf{Caractères} & \textbf{Pinyin} & \textbf{Étape} & \textbf{Traduction} \\ \hline
以前，妇女没有受教育的机会。 & Yǐqián, fùnǚ méiyǒu shòu jiàoyù de jīhuì. & Passé & Autrefois, les femmes n'avaient pas la possibilité de recevoir une éducation. \\ \hline
过去，女人只能在家做家务。 & Guòqù, nǚrén zhǐnéng zài jiā zuò jiāwù. & Passé & Autrefois, les femmes ne pouvaient que faire le ménage à la maison. \\ \hline
现在，女人的地位越来越高了。 & Xiànzài, nǚrén de dìwèi yuè lái yuè gāo le. & Présent & Maintenant, le statut de la femme est de plus en plus élevé. \\ \hline
如今，越来越多的女性上大学、工作。 & Rújīn, yuè lái yuè duō de nǚxìng shàng dàxué, gōngzuò. & Présent & De nos jours, de plus en plus de femmes vont à l'université et travaillent. \\ \hline
但是，男女工资有时还不一样。 & Dànshì, nánnǚ gōngzī yǒushí hái bù yīyàng. & Problème & Mais les salaires des hommes et des femmes sont parfois encore différents. \\ \hline
不过，家庭责任依然主要落在女人身上。 & Bùguò, jiātíng zérèn yīrán zhǔyào luò zài nǚrén shēnshàng. & Problème & Cependant, la responsabilité familiale retombe encore principalement sur les femmes. \\ \hline
我觉得男女应该平等，互相尊重。 & Wǒ juéde nánnǚ yīnggāi píngděng, hùxiāng zūnzhòng. & Opinion & Je pense que les hommes et les femmes devraient être égaux et se respecter mutuellement. \\ \hline
希望未来没有性别歧视。 & Xīwàng wèilái méiyǒu xìngbié qíshì. & Opinion & J'espère qu'à l'avenir il n'y aura plus de discrimination sexiste. \\ \hline
\end{tabularx}
\end{center}

\begin{exemplebox}[Structure 越来越 + adjectif + 了 — très rentable à l'écrit]
\zh{女人的地位越来越高了。} \\
Nǚrén de dìwèi yuè lái yuè gāo le. \\
« Le statut de la femme est de plus en plus élevé. » \\
\emph{Schéma :} Sujet + 越来越 + adjectif qualitatif + 了 (marque le changement d'état progressif). \\
Autres exemples utiles : \zh{生活越来越好} (la vie est de mieux en mieux), \zh{女性受教育程度越来越高} (le niveau d'éducation des femmes est de plus en plus élevé).
\end{exemplebox}

\courssub{Relecture guidée : checklist avant de rendre la copie}

\begin{aretenir}[Check-list de relecture (5 points de contrôle)]
\begin{enumerate}
\item \textbf{Nombre de phrases} : entre 5 et 8 (comptez les 。).
\item \textbf{Ordre des mots} : Sujet + Temps + Lieu + Complément d'objet + Verbe + Compléments / 了. Pas d'inversion à la française.
\item \textbf{Caractères à radical \cle{女}} : 妇 (fù), 妈 (mā), 姐 (jiě), 妹 (mèi), 妻 (qī), 她 (tā), 好 (hǎo), 如 (rú), 奶 (nǎi), 婚 (hūn), 嫁 (jià), 娘 (niáng). Vérifiez l'ordre des traits du radical \cle{女} : 1. 撇 (piě, trait diagonal haut-gauche vers bas-droite), 2. 横 (héng, trait horizontal), 3. 横 (héng, trait horizontal). Puis le reste du caractère (haut $\rightarrow$ bas, gauche $\rightarrow$ droite).
\item \textbf{Connecteurs variés} : au moins 3 différents parmi 以前 / 现在 / 但是 / 不过 / 因为…所以 / 我觉得 / 希望.
\item \textbf{Vocabulaire du thème} : au moins 5 mots du lexique ciblé (地位, 平等, 歧视, 受教育, 工作, 家务, 责任, 独立, 法律, 机会).
\end{enumerate}
\end{aretenir}

\begin{attention}[Erreur fréquente — oublier le 了 de changement d'état]
Beaucoup d'élèves écrivent : \zh{女人的地位越来越高} (sans 了). \\
\emph{Incorrect.} La structure \cle{越来越 + adjectif} exprime une évolution en cours ; elle \textbf{exige} le 了 final pour marquer le changement d'état réalisé / constaté. \\
\emph{Correct :} \zh{女人的地位越来越高了。} \\
Même piège avec : \zh{生活越来越好了} (pas \zh{生活越来越好}).
\end{attention}

\courssub{Grille d'évaluation sur 10 points (officielle, à recopier dans le cahier)}

\begin{textebox}[Grille d'évaluation — Production écrite « Situation de la femme » (10 pts)]
Critères (2 pts chacun) — Cochez la case correspondante : 0 / 1 / 2 pts
\begin{tabularx}{\linewidth}{|X|c|}
\hline
1. Respect du sujet (thème, pays, 5–8 phrases) & $\square$ 0 $\square$ 1 $\square$ 2 \\ \hline
2. Exactitude des caractères (traits, radicaux, pas de caractères inventés) & $\square$ 0 $\square$ 1 $\square$ 2 \\ \hline
3. Ordre des mots et grammaire (SVO, compléments, 了, 比, 越来越) & $\square$ 0 $\square$ 1 $\square$ 2 \\ \hline
4. Connecteurs variés et cohérence (au moins 3, progression logique) & $\square$ 0 $\square$ 1 $\square$ 2 \\ \hline
5. Richesse du vocabulaire du thème (min. 5 mots ciblés bien employés) & $\square$ 0 $\square$ 1 $\square$ 2 \\ \hline
\end{tabularx}\par\medskip
\textbf{Total : /10}
\end{textebox}

\begin{popfigure}
\begin{tikzpicture}[
  ampersand replacement=\&,
  cell/.style={draw=popline, minimum height=1.1cm, align=center, font=\small},
  header/.style={cell, fill=popblueL, font=\small\bfseries, text=popdark},
  crit/.style={cell, fill=popcream, font=\small\bfseries, minimum width=5.5cm},
  score/.style={cell, fill=popgreenL, minimum width=1.8cm},
  total/.style={cell, fill=popgoldL, font=\bfseries, minimum width=3cm}
]
\matrix (grid) [matrix of nodes, nodes in empty cells, column sep=-\pgflinewidth, row sep=-\pgflinewidth,
  column 1/.style={nodes={crit}},
  column 2/.style={nodes={score}},
  column 3/.style={nodes={score}},
  column 4/.style={nodes={score}},
  column 5/.style={nodes={total}},
  row 1/.style={nodes={header}}
]{
Critère (2 pts) \& 0 pt \& 1 pt \& 2 pts \& Total \\
1. Respect du sujet \& \& \& \& \\
2. Exactitude des caractères \& \& \& \& \\
3. Ordre des mots \& grammaire \& \& \& \& \\
4. Connecteurs \& cohérence \& \& \& \& \\
5. Vocabulaire du thème \& \& \& \& \\
\textbf{TOTAL} \& \& \& \& \textbf{/10} \\
};
\draw[popline, thick] (grid-1-1.north west) rectangle (grid-7-5.south east);
\end{tikzpicture}
\legende{Grille d'évaluation à recopier pour l'auto-correction et l'évaluation croisée.}
\end{popfigure}

\courssub{Auto-évaluation puis évaluation croisée entre pairs}

\begin{methode}[Auto-évaluation $\rightarrow$ Évaluation croisée (15 min en classe)]
\begin{enumerate}
\item \textbf{Auto-évaluation (5 min)} : relisez votre texte avec la grille. Attribuez 0, 1 ou 2 pts par critère. Notez le total. Entourez en rouge les caractères douteux, en bleu les connecteurs utilisés.
\item \textbf{Échange de copies (5 min)} : donnez votre texte à un camarade. Il/elle applique la grille sur \emph{votre} texte (pas sur le sien) et écrit son total en bas.
\item \textbf{Comparaison (5 min)} : comparez les deux grilles. Discutez les écarts : « Pourquoi as-tu mis 1 pt en grammaire ? » « J'ai vu l'ordre « 我去北京了 » au lieu de « 我去了北京 » ». Corrigez ensemble.
\item \textbf{Rédaction finale} : si le total pair est $<$ 7, réécrivez le texte en corrigeant les points faibles.
\end{enumerate}
\end{methode}

\courssub{Deux phrases « passe-partout » d'opinion à connaître par cœur (complément examen)}

\begin{exemplebox}[Phrases réutilisables — Opinion / Conclusion]
\begin{enumerate}
\item \zh{我觉得男女应该平等。} \\
Wǒ juéde nánnǚ yīnggāi píngděng. \\
« Je pense que les hommes et les femmes devraient être égaux. » \\
\emph{Structure :} Sujet + 觉得 + Sujet 2 + 应该 + adjectif/verbe. Utilisable sur \emph{tout} sujet d'opinion (environnement, éducation, santé…).
\item \zh{女人的地位越来越高了。} \\
Nǚrén de dìwèi yuè lái yuè gāo le. \\
« Le statut de la femme est de plus en plus élevé. » \\
\emph{Structure :} Sujet + 越来越 + adjectif + 了. Utilisable pour décrire toute amélioration (économie, éducation, santé, relations sino-camerounaises…).
\end{enumerate}
\emph{Astuce :} écrivez ces deux phrases au dos de votre convocation le jour du Bac. Elles vous garantissent 2 phrases correctes, riches en vocabulaire et structures, pour la partie « opinion ».
\end{exemplebox}

\courssub{Exercice résolu : rédaction complète corrigée}

\begin{exoresolu}[Rédiger un texte de 7 phrases sur la situation de la femme au Cameroun]
\textbf{Énoncé :} En vous appuyant sur la banque de phrases et le plan en 4 étapes, écrivez un texte de 7 phrases. Utilisez au moins 6 mots du lexique du thème et 4 connecteurs différents.

\tcblower
\textbf{Proposition de rédaction (niveau HSK 3–4, naturel et sans faute) :}

\zh{以前，喀麦隆的妇女没有受教育的机会，只能在家做家务。} \\
Yǐqián, Kāmàilóng de fùnǚ méiyǒu shòu jiàoyù de jīhuì, zhǐnéng zài jiā zuò jiāwù. \\
« Autrefois, les femmes du Cameroun n'avaient pas la possibilité de recevoir une éducation, elles ne pouvaient que faire le ménage à la maison. »

\zh{现在，情况不一样了，女人的地位越来越高了。} \\
Xiànzài, qíngkuàng bù yīyàng le, nǚrén de dìwèi yuè lái yuè gāo le. \\
« Maintenant, la situation n'est plus la même, le statut de la femme est de plus en plus élevé. »

\zh{越来越多的女性上大学，找到工作，变得独立。} \\
Yuè lái yuè duō de nǚxìng shàng dàxué, zhǎodào gōngzuò, biànde dúlì. \\
« De plus en plus de femmes vont à l'université, trouvent du travail, deviennent indépendantes. »

\zh{但是，男人的工资比女人的高，家庭责任依然主要落在女人身上。} \\
Dànshì, nánrén de gōngzī bǐ nǚrén de gāo, jiātíng zérèn yīrán zhǔyào luò zài nǚrén shēnshàng. \\
« Mais le salaire des hommes est plus élevé que celui des femmes, la responsabilité familiale retombe encore principalement sur les femmes. »

\zh{我觉得男女应该平等，互相尊重，共同分担家务。} \\
Wǒ juéde nánnǚ yīnggāi píngděng, hùxiāng zūnzhòng, gòngtóng fēndān jiāwù. \\
« Je pense que les hommes et les femmes devraient être égaux, se respecter mutuellement, et partager le ménage. »

\zh{希望未来没有性别歧视，每个女性都能实现梦想。} \\
Xīwàng wèilái méiyǒu xìngbié qíshì, měi gè nǚxìng dōu néng shíxiàn mèngxiǎng. \\
« J'espère qu'à l'avenir il n'y aura plus de discrimination sexiste, et que chaque femme pourra réaliser son rêve. »

\textbf{Analyse selon la grille :}
\begin{itemize}
\item \textbf{Respect du sujet} (2/2) : 7 phrases, Cameroun mentionné, progression passé/présent/problème/opinion.
\item \textbf{Caractères} (2/2) : radicaux \cle{女} corrects (妇, 女, 如), pas de caractères inventés. Attention : \zh{性} (xìng, dans \zh{女性}) a pour radical \cle{忄} (cœur), pas \cle{女}.
\item \textbf{Grammaire / ordre} (2/2) : SVO respecté, 了 de changement (不一样了, 越来越高了), 越来越+adj+了, 比较句 (男人的工资比女人的高), 觉得/认为 + proposition.
\item \textbf{Connecteurs} (2/2) : 以前 / 现在 / 但是 / 越来越多… / 我觉得 / 希望 = 6 connecteurs/marqueurs distincts.
\item \textbf{Vocabulaire thème} (2/2) : 妇女, 受教育, 机会, 家务, 地位, 独立, 工资, 责任, 平等, 歧视, 梦想 = 11 mots ciblés.
\end{itemize}
\textbf{Total : 10/10.}
\end{exoresolu}

\courssub{Exercices d'entraînement (à faire au brouillon, puis évaluation croisée)}

\begin{exercice}[Entraînement 1 — Variante « Mon village / Ma ville »]
Rédigez 6 phrases sur la situation de la femme dans votre village ou votre ville (Yaoundé, Douala, Buea, Ngaoundéré…). Adaptez le vocabulaire : \zh{农村} (nóngcūn, campagne), \zh{城市} (chéngshì, ville), \zh{市场} (shìchǎng, marché), \zh{卖菜} (màicài, vendre des légumes), \zh{创业} (chuàngyè, entreprendre). Utilisez au moins 3 connecteurs.
\end{exercice}

\begin{exercice}[Entraînement 2 — Version inversée (chinois $\rightarrow$ français)]
Traduisez le texte suivant en français, puis réécrivez-le en chinois en changeant l'opinion (ex : « Je pense qu'il faut plus de lois » au lieu de l'égalité).
\zh{过去，中国的女性也很少上学。现在，法律保护妇女权利，女大学生比男大学生还多。但是，高层管理职位里女性还是少。我觉得社会应该给女性更多机会。}
\end{exercice}

\begin{corrige}[Exercice 2]
\textbf{Traduction :} « Autrefois, les femmes chinoises allaient aussi rarement à l'école. Maintenant, la loi protège les droits des femmes, les étudiantes sont même plus nombreuses que les étudiants. Mais dans les postes de haute direction, les femmes sont encore peu nombreuses. Je pense que la société devrait donner plus d'opportunités aux femmes. »
\\ \textbf{Réécriture chinoise (opinion modifiée) :}
\zh{过去，中国的女性也很少上学。现在，法律保护妇女权利，女大学生比男大学生还多。但是，高层管理职位里女性还是少。我觉得社会应该制定更多法律，保障女性晋升机会。} \\
Guòqù, Zhōngguó de nǚxìng yě hěn shǎo shàngxué. Xiànzài, fǎlǜ bǎohù fùnǚ quánlì, nǚ dàxuéshēng bǐ nán dàxuéshēng hái duō. Dànshì, gāocéng guǎnlǐ zhíwèi lǐ nǚxìng hái shì shǎo. Wǒ juéde shèhuì yīnggāi zhìdìng gèng duō fǎlǜ, bǎozhàng nǚxìng jìnshēng jīhuì.
\end{corrige}

\courssub{Bilan de la leçon et conseils pour la révision}

\begin{aretenir}[Ce qu'il faut retenir pour le Bac — Module 1, Leçon 3]
\begin{itemize}
\item \textbf{Lexique clé (20 mots)} : 妇女, 女性, 地位, 平等, 歧视, 受教育, 机会, 独立, 法律, 权利, 家务, 责任, 工资, 管理, 职位, 晋升, 创业, 梦想, 互相尊重, 共同分担.
\item \textbf{Structures grammaticales maîtrisées} : 越来越+adj+了 ; 比较句 (A 比 B + adj) ; 觉得/认为 + proposition ; 应该/必须 + V ; 以前/现在/但是/不过/因为…所以/希望.
\item \textbf{Caractères à radical \cle{女}} : 妇, 妈, 姐, 妹, 妻, 她, 好, 如, 奶, 婚, 嫁, 娘. Ordre des traits du radical \cle{女} : 1. 撇 (piě), 2. 横 (héng), 3. 横 (héng) ; puis le reste du caractère (haut$\rightarrow$bas, gauche$\rightarrow$droite).
\item \textbf{Méthode de rédaction} : Plan 4 étapes $\rightarrow$ Phrases clés en main $\rightarrow$ Brouillon caractères+pinyin $\rightarrow$ Relecture check-list $\rightarrow$ Copie propre.
\item \textbf{Grille 10 pts} : Sujet (2) + Caractères (2) + Grammaire (2) + Connecteurs (2) + Vocabulaire (2).
\item \textbf{2 phrases « passe-partout »} : \zh{我觉得男女应该平等。} / \zh{女人的地位越来越高了。}
\end{itemize}
\end{aretenir}

\begin{savaistu}[Regard croisé Cameroun — Chine : la femme dans les deux sociétés]
\begin{itemize}
\item \textbf{Cameroun :} Loi n° 2016/007 du 12 juil. 2016 portant Code pénal (répression des violences conjugales) ; quota 30\% femmes aux élections locales (loi 2012) ; femmes très actives au marché (Bayam-Sellam, Marché Central) et dans l'agriculture vivrière ; scolarisation des filles en hausse mais décrochage au secondaire dans le Grand-Nord.
\item \textbf{Chine :} Loi sur la protection des droits des femmes (1992, révisée 2022) ; taux de scolarisation féminin $>$ 99\% au primaire ; 45\% de la main-d'œuvre ; mais « plafond de verre » dans les postes dirigeants ($\approx$ 13\% au comité central du PCC) ; pression sociale « femme restante » (剩女, shèngnǚ) après 27 ans.
\item \textbf{Point commun :} Dans les deux pays, la loi a progressé, l'éducation s'est féminisée, mais la charge domestique et les stéréotypes persistent. Le débat public évolue vite (réseaux sociaux, ONG, médias).
\end{itemize}
\end{savaistu}

\courssub{Devoir maison (noté sur 10 avec la grille)}

\begin{exercice}[Devoir — Production écrite notée]
\textbf{Consigne :} « Rédigez un texte de 6 à 8 phrases en chinois sur le thème : \emph{La femme camerounaise d'hier à aujourd'hui}. » \\
Utilisez obligatoirement : (1) au moins 6 mots du lexique de la leçon, (2) la structure 越来越+adj+了 une fois, (3) un comparatif avec 比, (4) au moins 4 connecteurs différents, (5) les deux phrases « passe-partout » d'opinion. \\
Remettez : (a) le brouillon annoté (check-list cochée), (b) la copie finale en caractères, (c) votre auto-évaluation sur la grille (total /10).
\end{exercice}

\begin{methode}[Conseils pour le devoir]
\begin{enumerate}
\item Commencez par lister au brouillon les 6+ mots de vocabulaire que vous allez caser.
\item Construisez vos 6–8 phrases une par une en vous référant au plan 4 étapes.
\item Vérifiez chaque phrase avec la check-list (radical 女, 了, ordre SVO, connecteur).
\item Faites l'auto-évaluation \textbf{avant} de rendre : soyez honnête, c'est ainsi qu'on progresse.
\end{enumerate}
\end{methode}

\newpage

\coursec{Activité d'intégration}

\coursec{Leçon 3 — Expression écrite : situation de la femme}

\courssub{Objectifs de l'activité}
\noindent À l'issue de cette activité d'intégration, l'élève sera capable de :
\begin{itemize}
    \item Mobiliser le lexique thématique (famille, travail, société, évolution) et les connecteurs logiques (\cle{以前/现在}, \cle{虽然/但是}, \cle{不但/而且}, \cle{因为/所以}, \cle{首先/然后}) pour structurer un texte cohérent.
    \item Produire un article court (6 à 8 phrases) en chinois simplifié, respectant le plan en quatre parties imposé (\textit{Autrefois $\rightarrow$ Aujourd'hui $\rightarrow$ Nuance $\rightarrow$ Opinion personnelle}).
    \item Écrire correctement les caractères clés du thème (\zh{女}, \zh{好}, \zh{妈}, \zh{平}, \zh{等}, \zh{独}, \zh{立}) en respectant l'ordre des traits et les proportions.
    \item S'auto-évaluer et évaluer un pair à l'aide d'une grille critériée (contenu, langue, caractères, présentation).
    \item Apprécier les évolutions sociétales comparées (Cameroun / Chine) sur la condition féminine.
\end{itemize}

\courssub{Situation}
\noindent Le club de presse du lycée bilingue de \textbf{Yaoundé} prépare l'édition spéciale « \textit{Femmes et Sociétés} » du journal mural \textit{L'Écho des Lycées}, publié en français, en anglais et en chinois. Votre binôme est chargé de rédiger la contribution en chinois. Le titre imposé est \zh{现在的女人} (\textit{La femme d'aujourd'hui}). L'article doit refléter la réalité camerounaise (évolution du rôle maternel, accès aux études et à l'emploi, persistance de certains stéréotypes, fierté de l'identité féminine) tout en utilisant les structures grammaticales et le vocabulaire travaillés dans la leçon (Module 1 : Vie familiale et intégration sociale). Un élève chinois en échange, \textbf{Lǐ Wén} (李文), relit les productions avant affichage.

\begin{textebox}[Consigne de rédaction remise aux élèves]
\textbf{Thème :} \zh{现在的喀麦隆女人} (La femme camerounaise d'aujourd'hui) \\
\textbf{Titre de l'article :} \zh{现在的女人} \\
\textbf{Plan obligatoire (4 paragraphes / parties) :}
\begin{enumerate}
    \item \zh{以前} (Autrefois) : situation traditionnelle (famille, foyer, dépendance).
    \item \zh{现在} (Aujourd'hui) : études, travail, indépendance, rôle social.
    \item \zh{但是} (Mais / Cependant) : défis restants, inégalités, charge mentale.
    \item \zh{我觉得} (Je pense que) : opinion personnelle, espoir, fierté.
\end{enumerate}
\textbf{Contraintes linguistiques :}
\begin{itemize}
    \item 6 à 8 phrases complètes en caractères simplifiés.
    \item Minimum \textbf{3 connecteurs} différents parmi la liste étudiée.
    \item Minimum \textbf{5 mots} du glossaire thématique (voir Ressources).
    \item Calligraphier soigneusement \textbf{un caractère} au choix (\zh{女} ou \zh{好}) en grand format pour illustrer l'article.
\end{itemize}
\end{textebox}

\courssub{Consignes}
\noindent \textbf{Travail en binôme (45 min + 15 min correction + 10 min affichage).}

\begin{enumerate}
    \item \textbf{Préparation et planification (10 min).} 
    \begin{itemize}
        \item Lire la situation et la consigne ci-dessus.
        \item Au brouillon, lister en français les idées principales pour chaque partie du plan (2--3 idées par partie).
        \item Sélectionner les 5+ mots du glossaire et les 3+ connecteurs à utiliser. Les noter avec leur pinyin.
        \item Vérifier l'ordre des traits des caractères \zh{女} et \zh{好} (voir Ressources) pour choisir celui à calligraphier.
    \end{itemize}

    \item \textbf{Rédaction de la première version (20 min).}
    \begin{itemize}
        \item Rédiger l'article en chinois (caractères + pinyin au-dessus si besoin) sur la fiche fournie (format A5).
        \item Respecter scrupuleusement l'enchaînement \zh{以前} $\rightarrow$ \zh{现在} $\rightarrow$ \zh{但是} $\rightarrow$ \zh{我觉得}.
        \item Soigner l'écriture des caractères (taille, proportion, ordre des traits). Calligraphier le caractère choisi (\zh{女} ou \zh{好}) en haut à droite de la fiche, dans un cadre prévu à cet effet.
        \item Relire à deux : vérifier l'ordre des mots (Sujet + Temps + Lieu + Verbe + Complément), l'accord des classificateurs, la place de \zh{了} (changement d'état) et \zh{过} (expérience).
    \end{itemize}

    \item \textbf{Échange et correction par les pairs (15 min).}
    \begin{itemize}
        \item Échanger votre fiche avec un autre binôme.
        \item Utiliser la \textbf{Grille d'évaluation} (distribuée par le professeur, voir ci-dessous) pour annoter la production : \cle{✓} (correct), \cle{∼} (maladroit), \cle{✗} (erreur). Écrire un commentaire constructif en bas de page.
        \item Récupérer votre fiche, lire les annotations, corriger les erreurs identifiées au stylo vert.
    \end{itemize}

    \item \textbf{Version finale et affichage (10 min).}
    \begin{itemize}
        \item Réécrire proprement la version corrigée sur une fiche neuve (ou au verso) si nécessaire.
        \item Afficher l'article sur le « Mur de chinois » (panneau d'affichage de la salle).
        \item L'enseignant sélectionne 2--3 productions pour une lecture à haute voix et une analyse collective en fin de séance.
    \end{itemize}
\end{enumerate}

\courssub{Ressources à mobiliser}
\noindent \textbf{1. Glossaire thématique (révision Leçon 3 -- Vocabulaire de base HSK 3/4)}

\begin{center}
\begin{tabularx}{\linewidth}{|X|X|X|X|}
\hline
\rowcolor{popblueL} \textbf{Caractères} & \textbf{Pinyin} & \textbf{Nature} & \textbf{Traduction} \\ \hline
女人 & nǚrén & n. & femme (adulte) \\ \hline
以前 & yǐqián & adv. & autrefois, avant \\ \hline
现在 & xiànzài & adv. & maintenant, aujourd'hui \\ \hline
家庭主妇 & jiātíng zhǔfù & n. & femme au foyer \\ \hline
工作 & gōngzuò & v./n. & travailler / travail, emploi \\ \hline
独立 & dúlì & adj./v. & indépendant / devenir indépendant \\ \hline
平等 & píngděng & adj./n. & égal / égalité \\ \hline
地位 & dìwèi & n. & statut, position sociale \\ \hline
进步 & jìnbù & v./n. & progresser / progrès \\ \hline
社会 & shèhuì & n. & société \\ \hline
照顾 & zhàogù & v. & s'occuper de, prendre soin de \\ \hline
丈夫 & zhàngfu & n. & mari, époux \\ \hline
孩子 & háizi & n. & enfant(s) \\ \hline
努力 & nǔlì & adv./v. & avec effort / faire des efforts \\ \hline
成功 & chénggōng & v. & réussir \\ \hline
挑战 & tiǎozhàn & n./v. & défi / défier \\ \hline
骄傲 & jiāo'ào & adj. & fier / fierté \\ \hline
\end{tabularx}
\end{center}

\noindent \textbf{2. Connecteurs logiques (outils de cohérence textuelle)}

\begin{center}
\begin{tabularx}{\linewidth}{|X|X|X|}
\hline
\rowcolor{poporangeL} \textbf{Structure / Connecteur} & \textbf{Exemple (Caractères + Pinyin)} & \textbf{Traduction} \\ \hline
\zh{以前……, 现在……} & \zh{以前女人在家，现在女人工作。} & Autrefois les femmes restaient à la maison, maintenant elles travaillent. \\ \hline
\zh{不但……，而且……} & \zh{她不但工作，而且照顾孩子。} & Elle non seulement travaille, mais en plus s'occupe des enfants. \\ \hline
\zh{虽然……，但是/可是……} & \zh{虽然她累，但是她快乐。} & Bien qu'elle soit fatiguée, elle est heureuse. \\ \hline
\zh{因为……，所以……} & \zh{因为她努力，所以她成功。} & Parce qu'elle fait des efforts, elle réussit. \\ \hline
\zh{首先……，然后……} & \zh{首先读书，然后找工作。} & D'abord étudier, ensuite trouver un travail. \\ \hline
\zh{此外} & \zh{此外，社会也变了。} & De plus, la société a aussi changé. \\ \hline
\zh{因此 / 所以} & \zh{地位提高了，因此她们自信。} & Le statut a monté, donc elles sont confiantes. \\ \hline
\end{tabularx}
\end{center}

\noindent \textbf{3. Points de grammaire clés (rappel synthétique)}

\begin{propriete}[Structures attendues dans l'article]
\begin{itemize}
    \item \textbf{Comparaison \zh{比} (bǐ) :} \zh{Sujet A + 比 + Sujet B + Adj.} \\
    \zh{现在的女人比以前自由。} (Les femmes d'aujourd'hui sont plus libres qu'avant.)
    \item \textbf{Changement d'état \zh{了} (le) en fin de phrase :} \zh{Adj/Verbe + 了.} \\
    \zh{她的地位变高了。} (Son statut est devenu plus élevé \textit{maintenant}.)
    \item \textbf{Expérience vécue \zh{过} (guò) :} \zh{Verbe + 过 + (Objet).} \\
    \zh{我妈妈当过医生。} (Ma mère a été médecin \textit{une fois dans sa vie}.)
    \item \textbf{Complément de degré \zh{得} (de) :} \zh{Verbe + 得 + Adj.} \\
    \zh{她工作得很努力。} (Elle travaille très dur.)
    \item \textbf{\zh{越来越} (yuè lái yuè) + Adj :} « De plus en plus... » \\
    \zh{社会越来越平等。} (La société est de plus en plus égalitaire.)
\end{itemize}
\end{propriete}

\noindent \textbf{4. Écriture des caractères : ordre des traits et radicaux (焦点)}

\begin{definition}[Caractères à calligraphier pour l'illustration]
\begin{itemize}
    \item \textbf{\zh{女} (nǚ) -- Femme.} \textbf{Radical :} \zh{女} (nǚzìpáng). \textbf{Traits (3) :} 1. Trait oblique descendant gauche \zh{ノ} ; 2. Trait oblique descendant droit \zh{丶} (partant du milieu du 1) ; 3. Trait horizontal long \zh{一} (traversant le bas). \textit{Astuce mnémotechnique :} La « femme » marche, le trait horizontal final représente le sol.
    \item \textbf{\zh{好} (hǎo) -- Bien, bon.} \textbf{Composé :} \zh{女} (gauche, radical) + \zh{子} (droite, enfant). \textbf{Traits (6) :} Écrire \zh{女} (3 traits) puis \zh{子} (3 traits : horizontal, vertical crochet, horizontal). \textit{Sens étymologique :} Une femme et un enfant ensemble = c'est bien / bon. \textbf{Proportion :} La partie gauche (\zh{女}) est plus étroite que la droite (\zh{子}) ; aligner les horizontales du bas.
\end{itemize}
\end{definition}

\begin{popfigure}
\begin{tikzpicture}[
    node distance=0.8cm and 1.2cm,
    box/.style={draw=popblue, thick, rounded corners, fill=popblueL, text width=2.8cm, align=center, minimum height=1.2cm},
    arrow/.style={-Stealth, thick, popdark}
]
\node[box, align=center] (p1){\textbf{以前}\\(Autrefois)\\\zh{家庭主妇, 照顾丈夫孩子}};
\node[box, right=of p1, align=center] (p2){\textbf{现在}\\(Aujourd'hui)\\\zh{上大学, 工作, 独立, 地位高}};
\node[box, right=of p2, align=center] (p3){\textbf{但是}\\(Cependant)\\\zh{挑战, 不平等, 压力}};
\node[box, right=of p3, align=center] (p4){\textbf{我觉得}\\(Mon avis)\\\zh{骄傲, 努力, 成功, 未来}};
\draw[arrow] (p1) -- (p2) node[midway, above, font=\small] {\zh{现在}};
\draw[arrow] (p2) -- (p3) node[midway, above, font=\small] {\zh{虽然...但是}};
\draw[arrow] (p3) -- (p4) node[midway, above, font=\small] {\zh{因为...所以}};
\end{tikzpicture}
\legende{Schéma de structure du texte : les 4 parties obligatoires et leurs connecteurs pivots.}
\end{popfigure}

\begin{attention}[Erreurs fréquentes à surveiller (Grille d'auto-correction)]
\begin{itemize}
    \item \textbf{Ordre des mots (Temps) :} Placer le complément de temps (\zh{以前}, \zh{现在}, \zh{今天}) \textbf{avant} le verbe ou en tête de phrase (Thème), jamais à la fin. \zh{✗ 我工作现在} $\rightarrow$ \zh{✓ 我现在工作} / \zh{✓ 现在我工作}.
    \item \textbf{\zh{的 / 得 / 地} (les trois \textit{de}) :} 
    \begin{itemize}
        \item \zh{的} (de) : après adjectif/nom $\rightarrow$ Nom. \textit{Ex:} \zh{努力的女人} (femme laborieuse).
        \item \zh{得} (de) : après Verbe $\rightarrow$ Adj/Degré. \textit{Ex:} \zh{工作得很努力} (travailler très dur).
        \item \zh{地} (de) : après Adj/Adv $\rightarrow$ Verbe (manière). \textit{Ex:} \zh{努力地工作} (travailler laborieusement).
    \end{itemize}
    \item \textbf{Classificateurs :} \zh{一个女人} (une femme, neutre), \zh{一位老师/女士} (un prof/une dame -- respect), \zh{个} vs \zh{位}. Ne pas oublier le classifieur après \zh{这/那/数量}.
    \item \textbf{\zh{和} vs \zh{而且} / \zh{也} :} \zh{和} lie des \textbf{noms} ; \zh{而且}/\zh{也} lient des \textbf{verbes/phrases}. \zh{✗ 她工作和照顾孩子} $\rightarrow$ \zh{✓ 她不但工作，而且照顾孩子} / \zh{✓ 她工作，也照顾孩子}.
    \item \textbf{Tons et Sandhi :} \zh{女人} (nǚrén $\rightarrow$ sandhi \textbf{núrén} 2+2 à l'oral), \zh{好} (hǎo -- 3ème ton), \zh{平等} (píngděng -- 2+3), \zh{独立} (dúlì -- 2+4). Attention au \zh{ü} (\zh{女} nǚ, \zh{绿} lǜ) : pas de tréma après j, q, x, y (\zh{居} jū, \zh{去} qù, \zh{需} xū, \zh{鱼} yú).
\end{itemize}
\end{attention}

\noindent \textbf{5. Grille d'évaluation par les pairs (à distribuer)}

\begin{center}
\begin{tabularx}{\linewidth}{|X|c|c|c|}
\hline
\rowcolor{popgreenL} \textbf{Critères} & \textbf{✓ Acquis} & \textbf{∼ En cours} & \textbf{✗ Non acquis} \\ \hline
\textbf{Contenu / Plan} : 4 parties respectées (以前/现在/但是/我觉得) & & & \\ \hline
\textbf{Lexique} : au moins 5 mots du glossaire utilisés correctement & & & \\ \hline
\textbf{Grammaire / Connecteurs} : au moins 3 connecteurs, structures \zh{比}/\zh{了}/\zh{过} correctes & & & \\ \hline
\textbf{Cohérence} : Sens logique, transitions fluides & & & \\ \hline
\textbf{Caractères} : Écriture lisible, ordre des traits respecté, proportions (calligraphie \zh{女} ou \zh{好}) & & & \\ \hline
\textbf{Pinyin / Tons} : Pinyin noté si hésitation, tons corrects à l'oral & & & \\ \hline
\textbf{Présentation} : Propreté, mise en page (titre, nom, illustration) & & & \\ \hline
\end{tabularx}
\end{center}

\courssub{Méthode : Rédiger un article d'opinion structuré}
\begin{methode}[Étapes de la rédaction]
\begin{enumerate}
    \item \textbf{Analyser le sujet et le plan imposé :} Identifier les 4 séquences temporelles/logiques (Passé $\rightarrow$ Présent $\rightarrow$ Concession $\rightarrow$ Opinion).
    \item \textbf{Brainstorming lexical (brouillon) :} Lister 8--10 mots du glossaire + 4--5 connecteurs. Vérifier la nature grammaticale (n, v, adj, adv) pour les placer correctement.
    \item \textbf{Rédiger phrase par phrase (Chinois) :} 
    \begin{itemize}
        \item P1 : \zh{以前} + Sujet + Verbe/Adj. (Contexte passé).
        \item P2 : \zh{现在} + Sujet + \zh{不但} V1, \zh{而且} V2 + \zh{越来越} Adj.
        \item P3 : \zh{虽然} Fait positif, \zh{但是} Problème (\zh{比如} Exemple).
        \item P4 : \zh{我觉得} Sujet + Adj (+ \zh{因为} Cause + \zh{了} Résultat).
    \end{itemize}
    \item \textbf{Relire : Check-list « S-T-L-V-C »} (Sujet - Temps - Lieu - Verbe - Complément), \zh{的/得/地}, \zh{了/过}, Classifieurs, Tons.
    \item \textbf{Calligraphier l'illustration :} Choisir \zh{女} (simple, radical) ou \zh{好} (composé, sens positif). Respecter l'ordre des traits (Haut$\rightarrow$Bas, Gauche$\rightarrow$Droite, Horizontal$\rightarrow$Vertical).
\end{enumerate}
\end{methode}

\courssub{Proposition de production et corrigé commenté}
\noindent Voici une production modèle respectant toutes les contraintes, suivie d'un commentaire linguistique détaillé pour l'enseignant.

\begin{exoresolu}[Article modèle : \zh{现在的女人} -- Niveau HSK 3/4 visé]
\textbf{Production de l'élève (Caractères + Pinyin + Traduction)} \\
\vspace{0.3cm}
\noindent \zh{以前，喀麦隆的女人大多是家庭主妇，在家照顾丈夫和孩子。} \\
\emph{Yǐqián, Kāmàilóng de nǚrén duōshì jiātíng zhǔfù, zài jiā zhàogù zhàngfu hé háizi.} \\
« Autrefois, la plupart des femmes camerounaises étaient femmes au foyer, s'occupant du mari et des enfants à la maison. » \\

\noindent \zh{现在，她们不但上大学，而且找到了好工作，地位越来越高。} \\
\emph{Xiànzài, tāmen bùdàn shàng dàxué, érqiě zhǎodào le hǎo gōngzuò, dìwèi yuè lái yuè gāo.} \\
« Aujourd'hui, elles non seulement vont à l'université, mais en plus trouvent de bons emplois, leur statut est de plus en plus élevé. » \\

\noindent \zh{虽然社会进步了，但是挑战还是有，比如工资不平等。} \\
\emph{Suīrán shèhuì jìnbù le, dànshì tiǎozhàn háishì yǒu, bǐrú gōngzī bù píngděng.} \\
« Bien que la société ait progressé, il reste des défis, par exemple les salaires inégaux. » \\

\noindent \zh{我觉得现在的女人很独立，也很骄傲，因为她们努力成功了。} \\
\emph{Wǒ juéde xiànzài de nǚrén hěn dúlì, yě hěn jiāo'ào, yīnwèi tāmen nǔlì chénggōng le.} \\
« Je pense que les femmes d'aujourd'hui sont très indépendantes et très fières, parce qu'elles ont réussi grâce à leurs efforts. » \\
\tcblower
\textbf{Commentaire linguistique et pédagogique (pour l'enseignant)} \\
\vspace{0.2cm}
\noindent \textbf{1. Respect du plan et connecteurs (Objectif atteint +) :} Les quatre parties sont clairement délimitées par les marqueurs discursifs \zh{以前}, \zh{现在}, \zh{虽然……但是……}, \zh{我觉得}. On note l'utilisation de \zh{不但……而且……} (Partie 2) et \zh{因为……} (Partie 4), soit 4 connecteurs complexes utilisés correctement (supérieur à 3 requis).

\noindent \textbf{2. Lexique thématique (Objectif atteint +) :} 10 mots du glossaire intégrés naturellement : \zh{女人}, \zh{家庭主妇}, \zh{照顾}, \zh{丈夫}, \zh{孩子}, \zh{工作}, \zh{地位}, \zh{进步}, \zh{平等}, \zh{独立}, \zh{骄傲}, \zh{努力}, \zh{成功} (occurrences multiples > 5 requis).

\noindent \textbf{3. Structures grammaticales ciblées :}
\begin{itemize}
    \item \zh{不但 A，而且 B} : Structure HSK 4 parfaite pour l'énumération d'actions positives (liaison verbe-verbe).
    \item \zh{越来越 + Adj} (\zh{越来越高}) : Expression de l'évolution progressive, très idiomatique.
    \item \zh{了} (changement d'état / aspect accompli) : \zh{进步了} (a progressé -- nouvel état), \zh{找到了} (action aboutie avec résultat \zh{到} + \zh{了}), \zh{成功了} (action accomplie avec résultat actuel).
    \item \zh{虽然……但是……} : Concession correcte (pas de « mais » redondant en chinois, structure standard).
    \item \zh{比如} (bǐrú) : Introduire un exemple concret (niveau HSK 4, bonus bienvenu).
    \item \zh{大多是} (duōshì) : Quantificateur « pour la plupart », évite l'écueil du classifieur manquant sur pluriel générique.
\end{itemize}

\noindent \textbf{4. Points d'attention / Difficultés pour les francophones (corrigés ici) :}
\begin{itemize}
    \item \textbf{Ordre des mots Temps/Sujet :} \zh{现在，她们……} (Correct : Thème \zh{现在} en tête). Pas de \zh{她们现在……} en tête de phrase pour la mise en relief thématique.
    \item \textbf{Négation \zh{不} vs \zh{没} :} \zh{不但} (structure figée, pas de négation réelle), \zh{不平等} (adjectif négatif = \zh{不} + Adj). Pas de \zh{没} car pas d'action passée non réalisée.
    \item \textbf{\zh{和} vs \zh{而且} :} Correctement évité \zh{和} entre verbes (\zh{上大学和找工作} $\times$) ; \zh{而且} utilisé pour lier les propositions verbales.
    \item \textbf{Complément de résultat \zh{找到} :} \zh{找} (chercher) + \zh{到} (trouver/résultat) + \zh{了} (accompli). Structure V+Résultat+了 = action aboutie.
\end{itemize}

\noindent \textbf{5. Calligraphie suggérée :} Le caractère \zh{好} (hǎo) est un excellent choix pédagogique : composé de \zh{女} (thème central) et \zh{子} (génération future/enfant), il symbolise la continuité et la valeur positive de la femme. L'élève doit veiller à l'alignement horizontal des deux composants (\zh{女} plus étroit, \zh{子} plus large) et à la longueur du trait final de \zh{子} (horizontal long fermant le caractère).
\end{exoresolu}

\courssub{Bilan de l'activité}
\noindent Cette activité d'intégration clôture la séquence « Expression écrite : situation de la femme » en mobilisant l'ensemble des compétences visées au programme officiel MINESEC pour la Terminale A (LV2 Chinois).

\begin{aretenir}[Compétences validées]
\begin{itemize}
    \item \textbf{Compréhension de l'écrit / Expression écrite :} L'élève a produit un texte authentique (article de journal), contextualisé (Cameroun), structuré (plan 4 parties), cohésif (connecteurs) et lexicalement riche (glossaire HSK 3/4).
    \item \textbf{Grammaire :} Application en contexte de production des structures \zh{比}, \zh{越来越}, \zh{了/过}, \zh{不但……而且……}, \zh{虽然……但是……}, \zh{因为……所以……}, \zh{大多是}.
    \item \textbf{Écriture des caractères :} Renforcement de la mémoire motrice pour les caractères à haute fréquence (\zh{女}, \zh{好}, \zh{妈}, \zh{平}, \zh{等}, \zh{独}, \zh{立}) via la calligraphie illustrative et la rédaction complète.
    \item \textbf{Traduction (Thème/Version) :} L'étape de rédaction du brouillon en français puis passage au chinois (thème guidé) et la relecture par les pairs (version implicite) consolident la médiation linguistique.
    \item \textbf{Expression orale / Interaction :} La lecture à haute voix des meilleures productions et le débat final (\textit{« Selon vous, quel est le plus grand défi pour la femme camerounaise aujourd'hui ? »}) ancrent la prononciation (tons, \zh{ü}, rétroflexes \zh{zh/ch/sh/r}) et l'argumentation.
    \item \textbf{Socioculturel :} Comparaison factuelle Cameroun/Chine (accès à l'éducation, taux d'activité féminine, congé maternité, place au politique) sans stéréotypes, ancrée dans le vécu des élèves (mères, tantes, enseignantes).
\end{itemize}

\noindent \textbf{Prolongements possibles (Différenciation / Approfondissement) :}
\begin{itemize}
    \item \textbf{Niveau HSK 4+ :} Rédiger une version longue (120--150 caractères) avec paragraphe supplémentaire sur « \zh{未来} (L'avenir) » en utilisant \zh{将来} / \zh{会} / \zh{希望} / \zh{应该}.
    \item \textbf{Projet interdisciplinaire (Histoire/Géo/EMC) :} Étude de cas comparée : \textit{La loi sur la parité au Cameroun vs La loi sur la protection des droits des femmes en Chine (révisée 2022)}. Vocabulaire juridique simple (\zh{法律} fǎlǜ, \zh{权利} quánlì, \zh{保护} bǎohù, \zh{歧视} qíshì).
    \item \textbf{Numérique :} Saisir l'article final sur l'application OnBuch+ (module « Portfolio élève »), enregistrer l'oral (reconnaissance vocale pinyin/tones), partager sur le mur de classe virtuel.
\end{itemize}

\noindent \textbf{Trace écrite obligatoire dans le cahier de l'élève (à dicter ou photocopier) :}
\begin{enumerate}
    \item Le plan type \zh{以前 $\rightarrow$ 现在 $\rightarrow$ 但是 $\rightarrow$ 我觉得} (schéma fléché ci-dessus).
    \item Le tableau des 5 connecteurs complexes (\zh{不但……而且……}, \zh{虽然……但是……}, \zh{因为……所以……}, \zh{首先……然后……}, \zh{越来越}).
    \item L'ordre des traits détaillés de \zh{女} (3 traits) et \zh{好} (6 traits = 3+3) avec flèches directionnelles.
    \item La grille d'auto-évaluation simplifiée (3 items : « J'ai utilisé 5 mots du thème », « Mes phrases suivent l'ordre S-T-L-V-C », « Mes caractères sont lisibles »).
\end{enumerate}
\end{aretenir}


\newpage

\coursec{Méthodes et exercices résolus}

\coursec{Leçon 3 — Expression écrite : situation de la femme}

\courssub{Méthodes et exercices résolus}

\begin{methode}[Analyser un texte modèle : compréhension globale, structure et connecteurs]
\textbf{Objectif :} Savoir extraire l'information principale, identifier la structure argumentative (thèse, arguments, exemples, conclusion) et repérer les connecteurs logiques pour les réutiliser.
\begin{enumerate}
    \item \textbf{Lecture rapide (1\textsuperscript{re} lecture) :} Identifier le \cle{sujet} (ici : la femme moderne), le \cle{point de vue} de l'auteur (évolution positive, défis restants) et le \cle{ton} (objectif, bienveillant).
    \item \textbf{Lecture détaillée (2\textsuperscript{de} lecture) :} Souligner les \textbf{mots-clés} du thème (vocabulaire cible) et les \textbf{connecteurs} (d'abord, cependant, de plus, en conclusion). Noter la structure : \textbf{Introduction (constat)} $\rightarrow$ \textbf{Développement (arguments contrastés : progrès / difficultés)} $\rightarrow$ \textbf{Conclusion (souhait / perspective)}.
    \item \textbf{Analyse grammaticale ciblée :} Repérer les structures modèles : \zh{随着……的发展} (avec le développement de...), \zh{不仅……而且……} (non seulement... mais aussi...), \zh{尽管……但是……} (bien que... cependant...), \zh{越来越} (de plus en plus).
    \item \textbf{Synthèse :} Rédiger un résumé en 3 phrases en français, puis tenter de le reformuler en chinois en réutilisant le vocabulaire et les connecteurs du texte.
\end{enumerate}
\textbf{Réflexe Bac :} Au brouillon, faites un \emph{mind-map} : Sujet au centre, branches = arguments principaux, sous-branches = exemples/chiffres/connecteurs.
\end{methode}

\begin{exoresolu}[Analyse du texte modèle « La femme d'aujourd'hui »]
\textbf{Énoncé :} Lisez le texte ci-dessous. Répondez aux questions en chinois (si possible) ou en français.
\begin{textebox}[Texte modèle : 现代女性的新面貌]
现代女性的新面貌 (Xiàndài nǚxìng de xīn miànmào)

随着社会的快速发展，女性的地位发生了巨大的变化。过去，妇女主要负责家务和照顾孩子，很少有机会接受高等教育或参加工作。现在，情况完全不同了。女性不仅受教育程度越来越高，而且在职场上也展现出了强大的竞争力。\\

例如，在喀麦隆和中国，越来越多的女性成为医生、工程师、教师甚至企业家。她们在工作中表现优秀，同时也能很好地平衡家庭和事业。然而，挑战依然存在。尽管法律保障了男女平等，但在现实生活中，部分女性仍面临就业歧视或家庭暴力的问题。\\

因此，全社会都应该努力，消除性别偏见，创造更公平的环境。只有这样，女性才能真正实现自我价值，为社会贡献更大的力量。
\end{textebox}
\textbf{Questions :}
\begin{enumerate}
    \item Donnez le titre du texte en français.
    \item Selon le texte, quelles étaient les occupations traditionnelles des femmes ? (Citez le passage en chinois + pinyin).
    \item Identifiez 4 connecteurs logiques utilisés dans le texte et donnez leur fonction (addition, opposition, cause, conséquence, exemple).
    \item Relevez deux structures grammaticales HSK 3/4 servant à exprimer le contraste et la progression.
    \item Traduisez la dernière phrase en français : \zh{只有这样，女性才能真正实现自我价值，为社会贡献更大的力量。}
\end{enumerate}
\tcblower
\textbf{Solution détaillée :}
\begin{enumerate}
    \item \textbf{Titre :} « Le nouveau visage de la femme moderne » / « La femme d'aujourd'hui ».
    \item \textbf{Occupations traditionnelles :} \zh{过去，妇女主要负责家务和照顾孩子} (Guòqù, fùnǚ zhǔyào fùzé jiāwù hé zhàogù háizi.) — « Autrefois, les femmes étaient principalement chargées des tâches ménagères et de la garde des enfants. »
    \item \textbf{Connecteurs :}
        \begin{itemize}
            \item \zh{现在} (xiànzài) — \textbf{Marque temporelle / Opposition} (contraste avec \zh{过去}).
            \item \zh{不仅……而且……} (bùjǐn... érqiě...) — \textbf{Addition / Progression} (accumulation d'arguments positifs).
            \item \zh{然而} (rán'ér) — \textbf{Opposition / Contraste} (introduction des défis).
            \item \zh{尽管……但是……} (jǐnguǎn... dànshì...) — \textbf{Concession / Opposition} (fait admis vs réalité).
            \item \zh{因此} (yīncǐ) — \textbf{Conséquence / Conclusion} (appel à l'action).
            \item \zh{例如} (lìrú) — \textbf{Illustration / Exemple}.
        \end{itemize}
    \item \textbf{Structures ciblées :}
        \begin{itemize}
            \item \zh{尽管……但是……} (Jǐnguǎn... dànshì...) : « Bien que... cependant... » (Concession). Niveau HSK 4.
            \item \zh{越来越 + Adj/Verbe} (Yuè lái yuè...) : « De plus en plus... » (Changement progressif). Niveau HSK 3.
            \item \zh{只有……才……} (Zhǐyǒu... cái...) : « Ce n'est que si... que... » (Condition nécessaire). Niveau HSK 4. \textbf{Point clé :} \zh{才} placé devant le verbe principal de la proposition principale.
        \end{itemize}
    \item \textbf{Traduction :} « Ce n'est qu'ainsi / Seulement de cette manière que les femmes pourront véritablement réaliser leur valeur personnelle, et contribuer d'une force plus grande à la société. »
    \begin{itemize}
        \item \zh{只有这样} (Zhǐyǒu zhèyàng) = Ce n'est qu'ainsi / Seulement ainsi.
        \item \zh{才能} (cái néng) = pouvoir (alors) / seulement alors pouvoir.
        \item \zh{实现自我价值} (shíxiàn zìwǒ jiàzhí) = réaliser sa valeur personnelle / s'accomplir.
        \item \zh{为社会贡献更大的力量} (wèi shèhuì gòngxiàn gèng dà de lìliàng) = contribuer d'une force plus grande à la société.
    \end{itemize}
\end{enumerate}
\textbf{Conclusion :} Ce texte suit une structure classique \textbf{Constat $\rightarrow$ Exemples $\rightarrow$ Nuance/Problème $\rightarrow$ Appel à l'action}. Les connecteurs \zh{然而} et \zh{donc} sont indispensables pour l'expression écrite au Bac.
\end{exoresolu}

\begin{methode}[Maîtriser le lexique thématique et l'écriture des caractères clés (radicaux, traits, confusions)]
\textbf{Objectif :} Mémoriser le vocabulaire « Situation de la femme » (HSK 3-4) et écrire correctement les caractères en respectant l'ordre des traits et en évitant les confusions visuelles.
\begin{enumerate}
    \item \textbf{Classement par radicaux (clés) :} Regrouper le vocabulaire par clé sémantique pour créer des liens logiques.
        \begin{itemize}
            \item Clé \zh{女} (nǚ, femme) : \zh{妈} (mère), \zh{好} (bon), \zh{妇} (femme mariée/épouse), \zh{姊} (sœur aînée), \zh{妹} (sœur cadette), \zh{娃} (bébé/enfant).
            \item Clé \zh{力} (lì, force) : \zh{努} (effort), \zh{功} (mérite/résultat), \zh{加} (ajouter), \zh{勉} (s'efforcer).
            \item Clé \zh{心} (xīn, cœur) : \zh{想} (penser), \zh{意} (idée/sens), \zh{平} (égal/paix $\rightarrow$ \zh{平等}), \zh{衡} (peser/mesurer $\rightarrow$ \zh{衡量}).
        \end{itemize}
    \item \textbf{Ordre des traits (règles d'or) :} \textbf{Haut $\rightarrow$ Bas}, \textbf{Gauche $\rightarrow$ Droite}, \textbf{Horizontal avant Vertical} (si croisant), \textbf{Extérieur avant Intérieur avant Fermeture}, \textbf{Trait central avant ailes}.
    \item \textbf{Chasse aux « faux jumeaux » (caractères proches) :} Créer des fiches de discrimination visuelle.
        \begin{itemize}
            \item \zh{妇} (fù, femme) vs \zh{母} (mǔ, mère) : \zh{妇} a \zh{女} à gauche + \zh{扌} (main) + \zh{帚} (balai) $\rightarrow$ femme qui tient le balai. \zh{母} = \zh{女} + deux points (seins) = mère qui allaite.
            \item \zh{平} (píng, égal/plat) vs \zh{年} (nián, an) : \zh{平} a un trait horizontal bas long (le sol plat). \zh{年} a \zh{禾} (riz) au-dessus de \zh{千}.
            \item \zh{争} (zhēng, lutter/concurrencer) vs \zh{事} (shì, affaire) : \zh{争} = \zh{爪} (griffe) au-dessus de \zh{亻} (homme) $\rightarrow$ se battre. \zh{事} = \zh{亻} + \zh{史} / \zh{吏}.
        \end{itemize}
    \item \textbf{Entraînement :} Écrire chaque caractère 5 fois en dictant l'ordre des traits à voix haute. Faire des phrases contextuelles.
\end{enumerate}
\end{methode}

\begin{exoresolu}[Vocabulaire, radicaux et discrimination de caractères]
\textbf{Énoncé :}
\begin{enumerate}
    \item Complétez le tableau de vocabulaire ci-dessous (Caractères, Pinyin, Nature, Traduction).
    \item Pour les caractères \zh{平}, \zh{妇}, \zh{争}, \zh{努}, indiquez : la clé (radical), le nombre de traits total, et un caractère « piège » visuellement proche.
    \item Choisissez le bon caractère pour compléter les phrases (attention aux faux jumeaux) :
        \begin{enumerate}
            \item 男女\_\_\_\_\_\_ (píngděng) 是基本人权。
            \item 她是一位伟大的\_\_\_\_\_\_ (fùnǚ) 代表。
            \item 在职场上，女性不应\_\_\_\_\_\_ (zhēngqǔ) 权利，而应平等获得。
            \item 我们要\_\_\_\_\_\_ (nǔlì) 学习，实现梦想。
        \end{enumerate}
\end{enumerate}

\begin{center}
\begin{tabularx}{\linewidth}{|X|X|X|X|}
\hline
\rowcolor{popblueL} \textbf{Caractères} & \textbf{Pinyin} & \textbf{Nature} & \textbf{Traduction} \\ \hline
平等 &  &  &  \\ \hline
妇女 &  &  &  \\ \hline
争取 &  &  &  \\ \hline
努力 &  &  &  \\ \hline
歧视 &  &  &  \\ \hline
暴力 &  &  &  \\ \hline
企业家 &  &  &  \\ \hline
平衡 &  &  &  \\ \hline
偏见 &  &  &  \\ \hline
实现 &  &  &  \\ \hline
\end{tabularx}
\end{center}
\tcblower
\textbf{Solution détaillée :}
\begin{enumerate}
    \item \textbf{Tableau complété :}
\begin{center}
\begin{tabularx}{\linewidth}{|X|X|X|X|}
\hline
\rowcolor{popblueL} \textbf{Caractères} & \textbf{Pinyin} & \textbf{Nature} & \textbf{Traduction} \\ \hline
平等 & píngděng & Nom / Adj & Égalité / Égal(e) \\ \hline
妇女 & fùnǚ & Nom & Femme(s) (terme formel) \\ \hline
争取 & zhēngqǔ & Verbe & Lutter pour / Revendiquer / Conquérir \\ \hline
努力 & nǔlì & Verbe / Adv & Faire des efforts / Avec effort \\ \hline
歧视 & qíshì & Verbe / Nom & Discriminer / Discrimination \\ \hline
暴力 & bàolì & Nom / Adj & Violence / Violent(e) \\ \hline
企业家 & qǐyèjiā & Nom & Entrepreneur / Chef d'entreprise \\ \hline
平衡 & pínghéng & Verbe / Nom & Équilibrer / Équilibre \\ \hline
偏见 & piānjiàn & Nom & Préjugé \\ \hline
实现 & shíxiàn & Verbe & Réaliser / Concrétiser \\ \hline
\end{tabularx}
\end{center}
    \item \textbf{Analyse des caractères :}
    \begin{center}
    \begin{tabularx}{\linewidth}{|X|X|X|X|}
    \hline
    \rowcolor{poporangeL} \textbf{Caractère} & \textbf{Clé (Radical)} & \textbf{Nb traits} & \textbf{Piège visuel (Différence)} \\ \hline
    \zh{平} & \zh{干} (gān, sec/perche) & 5 & \zh{年} (nián) : \zh{平} a un grand trait horizontal final (le sol). \\ \hline
    \zh{妇} & \zh{女} (nǚ, femme) & 8 & \zh{母} (mǔ) : \zh{妇} a \zh{女} à g. + \zh{扌}+\zh{帚} ; \zh{母} = \zh{女} + 2 points. \\ \hline
    \zh{争} & \zh{爪} (zhǎo, griffe) & 6 & \zh{事} (shì) : \zh{争} = griffe + homme ; \zh{事} = homme + histoire/officiel. \\ \hline
    \zh{努} & \zh{力} (lì, force) & 7 & \zh{奴} (nú, esclave) : \zh{努} a \zh{力} à dr. (effort) ; \zh{奴} a \zh{女} à g. + \zh{又}. \\ \hline
    \end{tabularx}
    \end{center}
    \item \textbf{Complétion (Choix du bon caractère) :}
    \begin{enumerate}
        \item 男女\zh{平等} (píngděng) — \zh{平} (plat/égal) et non \zh{年}.
        \item 她是一位伟大的\zh{妇女} (fùnǚ) — \zh{妇} (femme adulte/épouse) avec clé \zh{女}.
        \item 在职场上，女性不应\zh{争取} (zhēngqǔ) — \zh{争} (lutter/revendiquer) avec clé \zh{爪}. \textit{Note : « Revendiquer » est positif ici, \zh{争} n'est pas péjoratif.}
        \item 我们要\zh{努力} (nǔlì) — \zh{努} (effort) avec clé \zh{力} (force). Pas \zh{奴} (esclave).
    \end{enumerate}
\end{enumerate}
\textbf{Conclusion :} La connaissance des radicaux (\zh{女}, \zh{力}, \zh{爪}, \zh{干}) est le levier principal pour mémoriser l'écriture et le sens, et éviter les confusions sanctionnées à l'écrit (0 pt au mot si caractère faux).
\end{exoresolu}

\begin{methode}[Structurer un paragraphe argumentatif : connecteurs, ordre des mots et progression logique]
\textbf{Objectif :} Produire un paragraphe cohérent (80-120 caractères) sur « Les défis de la femme moderne » en respectant la structure chinoise : \textbf{Sujet + Contexte (时间/地点/条件) + Argument 1 + Connecteur + Argument 2 + Conclusion}.
\begin{enumerate}
    \item \textbf{Planifier au brouillon (5 min) :} Idée directrice $\rightarrow$ 2 arguments (ex: 1. Progrès éducatif/professionnel 2. Double fardeau / discriminations) $\rightarrow$ Connecteurs choisis $\rightarrow$ Conclusion (solution / espoir).
    \item \textbf{Maîtriser la « Boîte à outils » des connecteurs (HSK 3-4) :}
        \begin{center}
        \begin{tabularx}{\linewidth}{|X|X|X|}
        \hline
        \rowcolor{popgreenL} \textbf{Fonction} & \textbf{Connecteur (Car. + Pin.)} & \textbf{Position / Structure} \\ \hline
        \textbf{Suite / Addition} & \zh{而且} (érqiě), \zh{此外} (cǐwài), \zh{另外} (lìngwài) & Après le sujet (souvent) : S + 而且 + V... \\ \hline
        \textbf{Contraste / Opposition} & \zh{但是} (dànshì), \zh{然而} (rán'ér), \zh{可是} (kěshì) & Début 2\textsuperscript{e} proposition : S1..., 但是 S2... \\ \hline
        \textbf{Concession} & \zh{虽然} (suīrán) S V, \zh{但/但是} (dàn/dànshì) S V & \zh{虽然} en tête de prop. sub., \zh{但} en tête de prop. princ. \\ \hline
        \textbf{Cause} & \zh{因为} (yīnwèi) S V, \zh{所以} (suǒyǐ) S V & Paire corrélative obligatoire. \\ \hline
        \textbf{Conséquence} & \zh{因此} (yīncǐ), \zh{所以} (suǒyǐ) & Début de phrase ou après sujet. \\ \hline
        \textbf{Condition} & \zh{只有} (zhǐyǒu) S V, \zh{才} (cái) S V & \zh{只有}... \zh{才}... (Seule condition). \\ \hline
        \textbf{Exemple} & \zh{例如} (lìrú), \zh{比如} (bǐrú) & En début de phrase ou après deux-points. \\ \hline
        \end{tabularx}
        \end{center}
    \item \textbf{Ordre des mots (Règle d'or) :} \textbf{Sujet + Temps + Lieu + Manière + Verbe + Complément}. Ne jamais mettre le Lieu après le Verbe (sauf \zh{在} + lieu comme compl. de lieu antéposé). Ex: \zh{她在公司工作} (Elle travaille \textbf{dans l'entreprise}) et non \zh{她工作在公司}.
    \item \textbf{Vérification finale :} Accords \zh{的/得/地}, particules d'aspect (\zh{了}, \zh{过}, \zh{着}), classificateurs.
\end{enumerate}
\end{methode}

\begin{exoresolu}[Rédaction guidée : connecteurs et ordre des mots]
\textbf{Énoncé :} Rédigez un paragraphe d'environ 80 caractères chinois sur le thème : \zh{现代女性面临的挑战与机遇} (Défis et opportunités des femmes modernes). Utilisez \textbf{au moins 5 connecteurs} de la boîte à outils et les mots : \zh{平等}, \zh{压力}, \zh{平衡}, \zh{实现}.
\tcblower
\textbf{Proposition de rédaction (Modèle corrigé) :}
\begin{textebox}[Production écrite élève - Modèle]
随着社会发展，女性获得了更多受教育和就业的机会，而且在科技、教育等领域表现优秀。然而，现实中女性依然面临巨大的压力。虽然法律规定男女平等，但职场歧视和家庭重担依然存在，女性很难平衡事业与家庭。因此，全社会应该消除偏见，提供更多支持。只有这样，女性才能真正实现自我价值。
\end{textebox}
\textbf{Pinyin :} Suízhe shèhuì fāzhǎn, nǚxìng huòdéle gèngduō shòu jiàoyù hé jiùyè de jīhuì, érqiě zài kējì, jiàoyù děng lǐngyù biǎoxiàn yōuxiù. Rán'ér, xiànshí zhōng nǚxìng yīrán miànlín jùdà de yālì. Suīrán fǜlǜ guīdìng nánnǚ píngděng, dàn zhíchǎng qíshì hé jiātíng zhòngdān yīrán cúnzài, nǚxìng hěn nán pínghéng shìyè yǔ jiātíng. Yīncǐ, quán shèhuì yīnggāi xiāochú piānjiàn, tígōng gèngduō zhīchí. Zhǐyǒu zhèyàng, nǚxìng cái néng zhēnzhèng shíxiàn zìwǒ jiàzhí.

\textbf{Analyse structurelle (Grille d'évaluation) :}
\begin{itemize}
    \item \textbf{Introduction contextuelle :} \zh{随着社会发展} (Préposition + Sujet + Verbe) — Correct.
    \item \textbf{Argument 1 (Progrès) :} \zh{而且} (Addition) lie « éducation » et « performance au travail ».
    \item \textbf{Pivot (Contraste) :} \zh{然而} (Cependant) — Marque le tournant vers les problèmes.
    \item \textbf{Argument 2 (Problèmes) :} \zh{虽然}... \zh{但}... (Concession) — Structure HSK 4 maîtrisée. \zh{平衡} utilisé comme verbe « équilibrer ».
    \item \textbf{Conclusion / Solution :} \zh{因此} (Conséquence) + \zh{只有}... \zh{才能}... (Condition nécessaire) — Structure complexe HSK 4 parfaite.
    \item \textbf{Vocabulaire imposé :} \zh{平等}, \zh{压力}, \zh{平衡}, \zh{实现} — Tous employés correctement.
    \item \textbf{Ordre des mots :} \zh{女性很难平衡事业与家庭} (S + Adv + V + O) — Correct. \zh{在科技、教育等领域} (Locuteur) placé AVANT le verbe \zh{表现}.
\end{itemize}
\textbf{Note attendue :} 18-20/20 (Cohérence, structures variées, vocabulaire précis, zéro faute de caractère).
\end{exoresolu}

\begin{methode}[Thème (Français $\rightarrow$ Chinois) : stratégies de transposition grammaticale]
\textbf{Objectif :} Traduire des phrases françaises complexes en chinois correct en évitant le calque syntaxique.
\begin{enumerate}
    \item \textbf{Analyse de la phrase source :} Identifier le \textbf{sujet}, le \textbf{verbe}, les \textbf{compléments} (temps, lieu, manière, objet, but). Repérer les \textbf{pièges francophones} : « C'est... qui... », passif, subjonctif, noms abstraits, genres/nombres.
    \item \textbf{Transposition mot à mot (Brouillon) :} Écrire les mots chinois clés dans l'ordre français, puis \textbf{réordonner} selon la syntaxe chinoise (SVO strict, Compléments \textbf{avant} le verbe).
    \item \textbf{Application des règles critiques :}
        \begin{itemize}
            \item \textbf{Passif « être + PP » :} $\rightarrow$ \zh{被} (bèi, formel/négatif), \zh{让} (ràng, oral/neutre), \zh{叫} (jiào, oral) ou \textbf{sujet inanimé + Verbe} (ex: \zh{这个问题解决了}).
            \item \textbf{« C'est... que... » (Insistance) :} $\rightarrow$ \zh{是}... \zh{的} (shì... de) \textbf{uniquement pour le passé accompli}. Présent/Futur : pas de \zh{是...的}.
            \item \textbf{Comparatif :} « A est plus... que B » $\rightarrow$ \zh{A 比 B + Adj} (Pas de \zh{更} après \zh{比} sauf si « beaucoup plus » : \zh{比... 更/还/得多}).
            \item \textbf{Durée / Fréquence :} Complément de durée \textbf{après} le verbe (éventuellement redoublé) : \zh{学习了两年}. Fréquence \textbf{avant} le verbe : \zh{每天去一次}.
            \item \textbf{Subjonctif / Obligation :} « Il faut que... » $\rightarrow$ \zh{应该} (yīnggāi), \zh{必须} (bìxū), \zh{得} (děi) + Verbe.
        \end{itemize}
    \item \textbf{Vérification :} Tons, particules (\zh{了} aspect accompli / \zh{过} expérience / \zh{着} durée), \zh{的/得/地}, classificateurs.
\end{enumerate}
\end{methode}

\begin{exoresolu}[Thème guidé : phrases types Bac]
\textbf{Énoncé :} Traduisez en chinois (caractères + pinyin).
\begin{enumerate}
    \item La situation des femmes s'est beaucoup améliorée par rapport au passé.
    \item Bien qu'elle soit très occupée, elle arrive à s'occuper de sa famille.
    \item Il faut que les hommes fassent plus de travaux ménagers.
    \item C'est l'année dernière qu'elle a commencé à travailler comme ingénieure.
    \item Plus elle travaille dur, plus elle réussit.
\end{enumerate}
\tcblower
\textbf{Solution détaillée :}
\begin{enumerate}
    \item \textbf{Analyse :} Sujet = \zh{女性的地位} (Situation des femmes). Verbe = \zh{改善} (s'améliorer) + \zh{了} (changement d'état). Complément comparaison = \zh{比过去} (par rapport au passé). Degré = \zh{很大} / \zh{明显}.
    \textbf{Réponse :} \zh{女性的地位比过去改善了很多。} (Nǚxìng de dìwèi bǐ guòqù gǎishànle hěnduō.) — \textit{Structure : Sujet + 比 + Référence + Verbe + Complément de degré.}
    \item \textbf{Analyse :} Concession (\zh{虽然}... \zh{但是}...). Sujet = \zh{她}. Adj = \zh{忙} (occupée). Verbe principal = \zh{照顾} (s'occuper) + \zh{好} (résultat) ou \zh{兼顾} (concilier).
    \textbf{Réponse :} \zh{虽然她很忙，但(是)她还是能照顾好家庭。} (Suīrán tā hěn máng, dàn(shì) tā háishì néng zhàogù hǎo jiātíng.) — \textit{Note : \zh{还是} (encore/toujours) renforce la concession. \zh{照顾好} = V + Complément de résultat.}
    \item \textbf{Analyse :} Obligation (\zh{应该} / \zh{必须}). Sujet = \zh{男人} / \zh{男性}. Verbe = \zh{做} (faire) + Objet = \zh{家务} (tâches ménagères). « Plus » = \zh{多做} / \zh{多承担}.
    \textbf{Réponse :} \zh{男性应该多承担（做）一些家务。} (Nánxìng yīnggāi duō chéngdàn (zuò) yīxiē jiāwù.) — \textit{Éviter le passif « il faut que » $\rightarrow$ Sujet actif + Modal.}
    \item \textbf{Analyse :} Insistance sur le temps passé (\zh{是}... \zh{的}). Sujet = \zh{她}. Temps = \zh{去年}. Verbe = \zh{开始} (commencer) + \zh{当} / \zh{做} (travailler comme) + Métier = \zh{工程师}.
    \textbf{Réponse :} \zh{是去年她开始当工程师的。} (Shì qùnián tā kāishǐ dāng gōngchéngshī de.) — \textit{Structure \zh{是}... \zh{的} obligatoire car action passée datée. Ordre : \zh{是} + Temps + Sujet + Verbe + \zh{的}.}
    \item \textbf{Analyse :} Corrélation « Plus... plus... » $\rightarrow$ \zh{越……越……} (Yuè... yuè...). Verbe 1 = \zh{努力} / \zh{努力工作}. Verbe 2 = \zh{成功} / \zh{顺利}.
    \textbf{Réponse :} \zh{她工作越努力，越成功。} (Tā gōngzuò yuè nǔlì, yuè chénggōng.) — \textit{Structure : Sujet + Verbe/Adj 1 + \zh{越} + Verbe/Adj 2. Pas de « plus » (更) dans la structure \zh{越...越...}.}
\end{enumerate}
\textbf{Point de vigilance :} Phrase 4 : Ne pas écrire \zh{她是去年开始当工程师的} (Bien que compréhensible, la place de \zh{是} avant le sujet met l'insistance sur \zh{去年} de façon plus standard pour l'examen). Phrase 5 : Ne pas dire \zh{更努力，更成功}.
\end{exoresolu}

\begin{methode}[Version (Chinois $\rightarrow$ Français) : décodage des particules, aspects et structures figées]
\textbf{Objectif :} Traduire un texte chinois vers un français naturel, en rendant correctement les particules aspectuelles (\zh{了}, \zh{过}, \zh{着}), les structures \zh{被/把/是...的}, et les modaux.
\begin{enumerate}
    \item \textbf{Segmentation (Découpage) :} Couper la phrase en \textbf{Groupes Nominaux (GN)} et \textbf{Groupes Verbaux (GV)}. Identifier le \textbf{Sujet} (souvent en tête, mais peut être un thème/Topique).
    \item \textbf{Décodage des particules clés (Tableau de correspondance) :}
    \begin{center}
    \begin{tabularx}{\linewidth}{|X|X|X|}
    \hline
    \rowcolor{poppurpleL} \textbf{Particule / Structure} & \textbf{Sens grammatical} & \textbf{Rendu français type} \\ \hline
    \zh{V + 了} (phrase finale) & Aspect accompli / Changement d'état / Nouveau fait & Passé composé / « Maintenant... » / « Déjà... » \\ \hline
    \zh{V + 了 + Durée} & Action terminée, durée écoulée & « A fait qqch \textbf{pendant} [Durée] » \\ \hline
    \zh{V + 过} & Expérience vécue (au moins une fois) & « A \textbf{déjà} fait qqch » / « A \textbf{eu l'occasion de} » \\ \hline
    \zh{V + 着} & État durable / Action en cours (statique) & « Est en train de... » / Participe présent / Adjectif \\ \hline
    \zh{被 / 让 / 叫 + Agent + V} & Passif (subi / neutre / oral) & « Être / Se faire / Se laisser + Infinitif » \\ \hline
    \zh{把 + Objet + V + Résultat/Dir} & Transposition objet (action intentionnelle, aboutie) & « Prendre/Prendre l'objet et le... » $\rightarrow$ souvent \textbf{Voix active} en FR : « Elle \textbf{a rangé} la chambre. » \\ \hline
    \zh{是 ... 的} (Passé) & Focalisation sur circonstance (temps, lieu, manière) passé & Cleft sentence : « C'est [circonstance] \textbf{que}... » \\ \hline
    \zh{越 ... 越 ...} & Corrélation proportionnelle & « Plus... plus... » / « De plus en plus... » \\ \hline
    \end{tabularx}
    \end{center}
    \item \textbf{Rédaction française :} Ne pas traduire mot à mot. Chercher l'équivalent idiomatique. Attention aux \textbf{faux amis} : \zh{方便} (pratique/commode $\neq$ « commode » seul), \zh{麻烦} (ennui/problème $\neq$ « embêter » toujours), \zh{辛苦} (dur/pénible $\rightarrow$ « Merci pour votre peine / Bon courage »).
    \item \textbf{Relecture :} Temps des verbes (Passé composé vs Imparfait), concordance des temps, register de langue (soutenu pour le Bac).
\end{enumerate}
\end{methode}

\begin{exoresolu}[Version guidée : texte court « Témoignage »]
\textbf{Énoncé :} Traduisez en français le texte suivant.
\begin{textebox}[Texte à traduire]
王芳是我的同事。去年，她被公司派到北京开会，住在了一家很方便的酒店。会议结束后，她去过故宫，觉得非常壮观。现在，她正在学习法语，因为她想将来去非洲工作。她常说：「只有努力，才能成功。」
\end{textebox}
\tcblower
\textbf{Solution détaillée :}
\textbf{Segmentation \& Analyse :}
\begin{enumerate}
    \item \zh{王芳是我的同事。} (Wáng Fāng shì wǒ de tóngshì.) $\rightarrow$ « Wang Fang est mon collègue. » (Présent de description).
    \item \zh{去年，她被公司派到北京开会，} (Qùnián, tā bèi gōngsī pài dào Běijīng kāihuì,) $\rightarrow$ \zh{被} = passif formel. \zh{派到} = envoyer à. $\rightarrow$ « L'année dernière, elle a été envoyée par son entreprise à Pékin pour une réunion, »
    \item \zh{住在了一家很方便的酒店。} (Zhù zài le yī jiā hěn fāngbiàn de jiǔdiàn.) $\rightarrow$ \zh{住在} + \zh{了} (accompli) + Lieu. \zh{方便} = pratique/commode. $\rightarrow$ « elle a logé dans un hôtel très pratique (bien situé). »
    \item \zh{会议结束后，} (Huìyì jiéshù hòu,) $\rightarrow$ « Après la fin de la réunion / À l'issue de la réunion, »
    \item \zh{她去过故宫，} (Tā qùguò Gùgōng,) $\rightarrow$ \zh{去过} = \zh{去} + \zh{过} (Expérience). $\rightarrow$ « elle a été au Palais Impérial (la Cité Interdite), » / « elle a visité... »
    \item \zh{觉得非常壮观。} (Juéde fēicháng zhuàngguān.) $\rightarrow$ \zh{觉得} = trouver / avoir l'impression. $\rightarrow$ « l'a trouvé très impressionnant / grandiose. »
    \item \zh{现在，她正在学习法语，} (Xiànzài, tā zhèngzài xuéxí Fǎyǔ,) $\rightarrow$ \zh{正在} + V = En train de (Progressif). $\rightarrow$ « Actuellement, elle apprend le français / est en train d'apprendre le français, »
    \item \zh{因为她想将来去非洲工作。} (Yīnwèi tā xiǎng jiānglái qù Fēizhōu gōngzuò.) $\rightarrow$ \zh{将来} = à l'avenir / plus tard. $\rightarrow$ « car elle veut aller travailler en Afrique plus tard / à l'avenir. »
    \item \zh{她常说：「只有努力，才能成功。」} (Tā cháng shuō: « Zhǐyǒu nǔlì, cáinéng chénggōng. ») $\rightarrow$ Structure \zh{只有}... \zh{才}... $\rightarrow$ « Elle dit souvent : "Ce n'est qu'en faisant des efforts / Seuls les efforts permettent de réussir." »
\end{enumerate}
\textbf{Traduction finale soignée :}
\begin{quote}
« Wang Fang est ma collègue. L'an dernier, elle a été envoyée par son entreprise à Pékin pour une conférence et a logé dans un hôtel très bien situé. À l'issue de la conférence, elle a visité la Cité Interdite, qu'elle a trouvée absolument grandiose. Actuellement, elle apprend le français car elle souhaite aller travailler en Afrique par la suite. Elle répète souvent : "Seuls les efforts mènent à la réussite." »
\end{quote}
\textbf{Points grammaticaux notés :} \zh{被} (passif), \zh{了} (accompli \zh{住在了}), \zh{过} (expérience \zh{去过}), \zh{正在} (progressif), \zh{只有...才...} (condition nécessaire), \zh{方便} (pratique $\neq$ commode seul), \zh{壮观} (grandiose/impressionnant).
\end{exoresolu}

\begin{methode}[Production écrite autonome : de la planification à l'auto-évaluation (Grille Bac)]
\textbf{Objectif :} Rédiger un texte argumentatif structuré (120-150 caractères) en 30 min, puis s'auto-évaluer selon les critères officiels.
\begin{enumerate}
    \item \textbf{Analyse du sujet (3 min) :} Identifier le \textbf{type de texte} (article d'opinion, lettre, discours, message forum), le \textbf{public visé}, le \textbf{registre} (formel \zh{您/您们} vs courant \zh{你/你们}), les \textbf{contraintes} (mots obligatoires, nombre de caractères).
    \item \textbf{Remue-méninges \& Plan (7 min) :} \textit{Mind-map} au crayon. Structure type 3 parties :
        \begin{itemize}
            \item \textbf{Introduction (2-3 phrases) :} Accroche (fait social, citation, proverbe \zh{妇女能顶半边天}) $\rightarrow$ Annonce du thèse $\rightarrow$ Plan implicite.
            \item \textbf{Développement (2 paragraphes, 4-5 phrases chacun) :}
                \begin{itemize}
                    \item Para 1 : Progrès / Droits acquis / Exemples (Chine/Cameroun). Connecteurs : \zh{首先}, \zh{而且}, \zh{例如}.
                    \item Para 2 : Défis restants / Double journée / Plafond de verre. Connecteurs : \zh{然而}, \zh{虽然}...\zh{但}, \zh{依然}.
                \end{itemize}
            \item \textbf{Conclusion (2 phrases) :} Synthèse + Appel à l'action / Perspective. Connecteurs : \zh{因此}, \zh{只有}...\zh{才}.
        \end{itemize}
    \item \textbf{Rédaction (15 min) :} Écrire au brouillon en caractères (pas en pinyin !). Vérifier l'ordre des mots \textbf{à chaque phrase}. Utiliser des structures « sûres » maîtrisées (HSK 3) plutôt que des structures risquées (HSK 5+ mal maîtrisées).
    \item \textbf{Relecture \& Auto-correction (5 min) :} Appliquer la \textbf{Grille d'auto-évaluation} ci-dessous.
\end{enumerate}
\end{methode}

\begin{exoresolu}[Production écrite complète + Auto-évaluation]
\textbf{Sujet d'examen type :} \zh{请以「女性与社会发展」为题，写一篇短文（120-150字），谈谈女性在现代社会中的作用、面临的挑战以及你的建议。要求：使用至少5个连接词，正确书写汉字。}
\textbf{Proposition d'élève (Brouillon corrigé) :}
\begin{textebox}[Production écrite - Version finale]
题目：女性与社会发展

妇女能顶半边天。随着社会的快速发展，女性在各个领域都发挥着重要作用。首先，受教育程度的提高让女性拥有了更多就业机会，而且在科学、政治、商业等方面取得了杰出成就。例如，喀麦隆和中国都有许多优秀的女性部长、企业家和科学家。\\

然而，挑战依然存在。虽然法律保障了男女平等，但职场歧视、薪资差距和家庭暴力等问题仍然普遍。许多女性不得不在事业和家庭之间做艰难的选择，承受巨大的心理压力。\\

因此，全社会必须采取行动。只有消除性别偏见，提供平等的托儿服务和晋升机会，女性才能真正实现自我价值，为构建和谐社会贡献力量。
\end{textebox}
\textbf{Auto-évaluation (Grille Bac - 20 pts) :}
\begin{center}
\begin{tabularx}{\linewidth}{|X|X|X|}
\hline
\rowcolor{popgoldL} \textbf{Critère} & \textbf{Indicateurs de réussite} & \textbf{Auto-note /5} \\ \hline
\textbf{Contenu / Respect sujet} (5 pts) & Thèse claire, 2 arguments distincts, exemples Cameroun/Chine, conclusion avec proposition. & 5/5 \\ \hline
\textbf{Organisation / Cohérence} (5 pts) & Structure 3 parties visible, \textbf{5+ connecteurs} (\zh{随着, 首先, 而且, 例如, 然而, 虽然...但, 因此, 只有...才}), paragraphes distincts. & 5/5 \\ \hline
\textbf{Grammaire / Structures} (5 pts) & \zh{随着...}, \zh{受教育程度提高} (Sujet verbal), \zh{不得不} (contrainte), \zh{只有...才...} (HSK4), \zh{被} (absent mais passif non forcé), ordre mots correct. & 4,5/5 (mineur: \zh{普遍} place) \\ \hline
\textbf{Vocabulaire / Caractères} (5 pts) & Lexique thème précis (\zh{薪资差距, 晋升机会, 和谐社会, 托儿服务}), 0 faute de caractère (vérif. \zh{妇, 平, 争, 努}), tons corrects implicites. & 5/5 \\ \hline
\textbf{TOTAL} &  & \textbf{19,5/20} \\ \hline
\end{tabularx}
\end{center}
\textbf{Commentaire correcteur :} Excellent respect du format. L'introduction par le proverbe \zh{妇女能顶半边天} ancre culturellement le propos. L'usage de \zh{不得不} (devoir/être obligé de) et \zh{薪资差距} (écart salarial) montre un niveau HSK 4 solide. La structure \zh{只有}...\zh{才能} en conclusion est le marqueur syntaxique attendu pour la « proposition de solution ».
\end{exoresolu}

\begin{attention}[Les 5 erreurs qui coûtent des points au Bac (Spécial « Situation de la femme »)]
\begin{enumerate}
    \item \textbf{Confusion \zh{妇} (fù) / \zh{母} (mǔ) / \zh{女} (nǚ) :} Écrire \zh{妇女} (femmes) avec \zh{母} ou \zh{女} seul. \textbf{Règle :} \zh{妇女} = terme formel « femmes » (clé \zh{女} + balai). \zh{母亲} = mère (clé \zh{女} + points). \zh{女性} = genre féminin (clé \zh{生} + \zh{女}). \textit{Sanction : 0,5 pt / erreur de caractère.}
    \item \textbf{Ordre des mots : Complément de lieu APRÈS le verbe.} \textit{Faux :} \zh{她工作在公司。} \textbf{Juste :} \zh{她在公司工作。} (Préposition \zh{在} + Lieu \textbf{avant} Verbe). Idem pour \zh{从} (depuis), \zh{跟} (avec), \zh{给} (pour/à).
    \item \textbf{Mauvaise place de \zh{了} (Aspect) :} \textit{Faux :} \zh{她了去年结婚。} \textbf{Juste :} \zh{她去年结婚了。} (Changement d'état / Nouveau fait en fin de phrase). \textit{Faux :} \zh{我学了中文两年。} (Ambigu). \textbf{Juste :} \zh{我学中文了两年了。} (Durée écoulée + action continue ou non).
    \item \textbf{Calque du comparatif français : « Plus que » $\rightarrow$ \zh{比...更...} sans \zh{得多} pour « beaucoup plus ».} \textit{Faux :} \zh{女性比男性更努力。} (Correct pour « plus... »). \textit{Mais si « beaucoup plus » :} \zh{女性比男性努力得多。} (Pas de \zh{更} devant \zh{得多}). \textit{Erreur fréquente :} \zh{比更努力得多}.
    \item \textbf{Oubli de \zh{才} dans \zh{只有}... \zh{才}... :} \textit{Faux :} \zh{只有努力，能成功。} \textbf{Juste :} \zh{只有努力，才能成功。} \zh{才} est \textbf{obligatoire} en proposition principale pour marquer la condition nécessaire exclusive. C'est une structure « à points » systématique au Bac.
\end{enumerate}
\textbf{Astuce de révision :} Faites une « Fiche Pièges » A5 avec ces 5 erreurs + 1 exemple correct / 1 exemple faux barré. Relisez-la la veille de l'épreuve.
\end{attention}

\newpage

\coursec{Exercices}

\courssub{Je vérifie mes connaissances}

\begin{exercice}[QCM : vocabulaire et structures]
Choisissez la bonne réponse parmi A, B, C. Une seule réponse est correcte.
\begin{enumerate}
    \item Comment dit-on « l'égalité entre les hommes et les femmes » ?
    \begin{itemize}
        \item A. 男女平等
        \item B. 男女不平等
        \item C. 男女一样
    \end{itemize}
    \item Complétez : 妈妈每天\_\_\_\_做家务。 (Maman fait le ménage tous les jours.)
    \begin{itemize}
        \item A. 都
        \item B. 也
        \item C. 还
    \end{itemize}
    \item Quel connecteur exprime la conséquence (« donc », « c'est pourquoi ») ?
    \begin{itemize}
        \item A. 因为
        \item B. 但是
        \item C. 所以
    \end{itemize}
    \item La phrase « 她不但聪明，而且勤劳 » signifie :
    \begin{itemize}
        \item A. Elle est non seulement intelligente, mais aussi travailleuse.
        \item B. Elle est intelligente, mais pas travailleuse.
        \item C. Elle est soit intelligente, soit travailleuse.
    \end{itemize}
\end{enumerate}
\end{exercice}

\begin{exercice}[Vrai ou Faux : compréhension du texte modèle]
Lisez le texte ci-dessous et dites si les affirmations sont Vraies (V) ou Fausses (F). Justifiez en citant un extrait du texte.
\begin{textebox}[Extrait : La femme d'aujourd'hui]
现代妇女的地位提高了。以前，很多妇女只在家里做家务、带孩子。现在，她们不但可以上班、工作，而且可以当领导、当医生。但是，社会上还有人认为女人应该以家庭为主。所以，我们要继续努力，争取真正的男女平等。
\end{textebox}
\begin{enumerate}
    \item Autrefois, les femmes ne travaillaient qu'à la maison.
    \item Aujourd'hui, les femmes peuvent devenir dirigeantes ou médecins.
    \item Tout le monde pense maintenant que les femmes doivent travailler.
    \item Le texte affirme que l'égalité totale est déjà atteinte.
\end{enumerate}
\end{exercice}

\begin{exercice}[Vocabulaire : association caractères / pinyin / sens]
Reliez chaque caractère (ou mot) à son pinyin et à sa traduction française.
\begin{center}
\begin{tabularx}{\linewidth}{|X|X|X|X|}
\hline
\rowcolor{popblueL} Caractères & Pinyin & Nature & Traduction \\ \hline
妇女 & fùnǔ & n. & femme (terme formel) \\ \hline
地位 & dìwèi & n. & statut, position sociale \\ \hline
提高 & tígāo & v. & augmenter, améliorer \\ \hline
家务 & jiāwù & n. & tâches ménagères \\ \hline
领导 & lǐngdǎo & n./v. & dirigeant / diriger \\ \hline
争取 & zhēngqǔ & v. & lutter pour, s'efforcer d'obtenir \\ \hline
平等 & píngděng & adj. & égalitaire, égal \\ \hline
社会 & shèhuì & n. & société \\ \hline
\end{tabularx}
\end{center}
\textbf{Consigne :} Écrivez les correspondances correctes (ex: 1-fùnǔ-femme).
\end{exercice}

\begin{exercice}[Pinyin et tons : discrimination auditive]
Associez chaque syllabe à son ton correct (1er, 2e, 3e, 4e, neutre) puis écrivez le pinyin complet du mot.
\begin{enumerate}
    \item 妇 \_\_\_\_ (mot : 妇女)
    \item 务 \_\_\_\_ (mot : 家务)
    \item 领 \_\_\_\_ (mot : 领导)
    \item 平 \_\_\_\_ (mot : 平等)
\end{enumerate}
\textbf{Rappel :} 1er ton (haut plat \textasciimacron), 2e ton (montant \textasciiacute), 3e ton (bas descendant-remontant \textasciicaron), 4e ton (descendant \textasciigrave).
\end{exercice}

\courssub{Je m'entraîne}

\begin{exercice}[Thème : phrases à traduire en chinois]
Traduisez les phrases suivantes en chinois (caractères + pinyin). Utilisez le vocabulaire de la leçon.
\begin{enumerate}
    \item Le statut des femmes s'est amélioré.
    \item Ma mère fait les tâches ménagères tous les jours.
    \item Les femmes peuvent non seulement travailler, mais aussi être des dirigeantes.
    \item Parce qu'elle est fatiguée, elle ne veut pas cuisiner ce soir.
    \item Bien qu'il y ait des progrès, l'égalité n'est pas encore complète.
\end{enumerate}
\end{exercice}

\begin{exercice}[Transformation : passé / présent / avis]
Réécrivez la phrase de base selon les indications. Base : \zh{女人在家里做饭。} (Les femmes cuisinent à la maison.)
\begin{enumerate}
    \item Au passé (autrefois) : \zh{以前，...}
    \item Au présent (maintenant) : \zh{现在，...}
    \item Négation au présent : \zh{现在，女人不...}
    \item Avec \zh{不但...而且...} (non seulement... mais aussi) : \zh{女人不但...而且...}
\end{enumerate}
\end{exercice}

\begin{exercice}[Texte à trous : connecteurs logiques]
Complétez le texte avec les connecteurs : \zh{因为}、\zh{所以}、\zh{但是}、\zh{而且}、\zh{虽然}。
\begin{textebox}[Texte à compléter]
\_\_\_\_ 现代妇女受教育程度高了，\_\_\_\_ 她们找到好工作容易多了。\_\_\_\_，很多妇女下班后还要做家务，\_\_\_\_ 很累。\_\_\_\_ 社会在进步，我们还是要争取男女平等。
\end{textebox}
\end{exercice}

\begin{exercice}[Remise en ordre : structure de phrase]
Remettez les mots dans l'ordre pour former une phrase correcte. Écrivez la phrase en caractères, pinyin et traduction.
\begin{enumerate}
    \item 妇女 / 地位 / 了 / 提高 / 的 / 现代
    \item 应该 / 家务 / 男人 / 和 / 一起 / 做 / 女人
    \item 很 / 她 / 累 / 所以 / 工作 / 了 / 下班 / 后
    \item 不但 / 可以 / 当 / 医生 / 而且 / 可以 / 当 / 领导 / 妇女
\end{enumerate}
\end{exercice}

\courssub{Je me prépare au Bac}

\begin{exercice}[Compréhension de texte et langue : type Bac]
Lisez attentivement le texte original ci-dessous, puis répondez aux questions en français.
\begin{textebox}[Situation de la femme au Cameroun et en Chine]
喀麦隆和中国的妇女地位都有很大变化。在喀麦隆，以前很多妈妈在家里种地、卖菜、带孩子。现在，越来越多的女孩上大学，当老师、护士、律师。在北京或上海，女人开公司、当经理很平常。但是，两个国家都有共同的问题：女人工作回来，还要洗衣服、做饭、辅导作业。男人分担家务的还不够多。真正的平等，不只是法律上，更要靠家庭里的行动。
\end{textebox}
\textbf{Questions :}
\begin{enumerate}
    \item \textbf{Compréhension globale :} Quel est le thème principal de ce texte ? (2 pts)
    \item \textbf{Détails :} Citez trois métiers exercés aujourd'hui par les femmes au Cameroun selon le texte. (3 pts)
    \item \textbf{Comparaison :} Quel point commun le texte souligne-t-il entre le Cameroun et la Chine concernant la charge domestique ? (2 pts)
    \item \textbf{Vocabulaire en contexte :} Expliquez le sens de \zh{分担} dans la phrase \zh{男人分担家务的还不够多} et donnez un synonyme possible. (2 pts)
    \item \textbf{Grammaire :} Repérez dans le texte la structure exprimant la concession (« bien que... ») et recopiez-la. Traduisez-la. (2 pts)
    \item \textbf{Expression personnelle :} Selon vous, que signifie « \zh{真正的平等，不只是法律上，更要靠家庭里的行动} » ? Donnez un exemple concret. (4 pts)
\end{enumerate}
\textit{Total : 15 points}
\end{exercice}

\begin{exercice}[Thème et Version : type Bac]
\textbf{Partie A : Thème (Français $\rightarrow$ Chinois) — 10 points}
Traduisez le texte suivant en chinois (caractères simplifiés + pinyin). Soignez l'ordre des mots, les mots-outils (\zh{了}、\zh{的}、\zh{地}) et les connecteurs.
\begin{quote}
« Autrefois, beaucoup de femmes restaient à la maison pour s'occuper des enfants. Maintenant, elles peuvent travailler dehors, mais elles font encore beaucoup de tâches ménagères. Je pense que les hommes doivent partager le travail à la maison. Ainsi, les femmes ne seront pas trop fatiguées. »
\end{quote}

\textbf{Partie B : Version (Chinois $\rightarrow$ Français) — 10 points}
Traduisez le texte suivant en français correct et naturel.
\begin{textebox}[Texte à traduire]
以前，我妈妈不识字，只能在农村种地。后来，她努力学习，识了很多字。现在她在医院当护士，工资不低。我爸爸也开始学做饭、洗碗。我们一家人觉得现在的生活比以前幸福多了。
\end{textebox}
\end{exercice}

\begin{exercice}[Expression écrite : sujet type Bac — 15 points]
\textbf{Sujet :} \emph{« Décrivez la situation de la femme dans votre famille ou dans votre quartier. Avez-vous constaté des changements par rapport au passé ? Qu'est-ce qui reste à améliorer pour atteindre l'égalité ? »}

\textbf{Consignes :}
\begin{itemize}
    \item Rédigez un texte cohérent de \textbf{80 à 100 caractères chinois} (environ 8 à 10 phrases).
    \item Utilisez \textbf{au moins 5 mots du vocabulaire de la leçon} (surlignez-les ou mettez-les en gras dans votre copie).
    \item Utilisez \textbf{au moins 3 connecteurs logiques} (\zh{因为/所以}、\zh{但是}、\zh{而且}、\zh{虽然/可是}、\zh{不但/而且}...).
    \item Respectez l'ordre des mots chinois (Sujet + Temps + Lieu + Verbe + Complément).
    \item Soignez l'écriture des caractères (traits, proportions).
\end{itemize}

\textbf{Grille d'évaluation indicative (sur 15 pts) :}
\begin{center}
\begin{tabularx}{\linewidth}{|X|c|}
\hline
\rowcolor{popblueL} Critères & Points \\ \hline
Contenu : pertinence, idées claires, exemples personnels (Cameroun) & 4 \\ \hline
Vocabulaire : richesse, exactitude, 5 mots de la leçon utilisés & 3 \\ \hline
Grammaire : structures correctes, connecteurs (3 min.), mots-outils & 4 \\ \hline
Caractères : écriture lisible, traits corrects, pas de caractères inventés & 2 \\ \hline
Cohérence / Ponctuation / Pinyin (si fourni) : fluidité, logique & 2 \\ \hline
\end{tabularx}
\end{center}
\end{exercice}

\newpage

\coursec{Corrigés des exercices}

\begin{corrige}[Exercice 1 — QCM : vocabulaire et structures]
\begin{enumerate}
    \item \textbf{Réponse A :} \zh{男女平等} (nán nǚ píngděng) « égalité entre hommes et femmes ». B signifie « inégalité » (\zh{不} = négation) ; C (\zh{一样} « pareil ») est correct grammaticalement mais n'est pas le terme consacré.
    \item \textbf{Réponse A :} \zh{妈妈每天都做家务。} (Māma měitiān dōu zuò jiāwù.) \zh{都} « tous/toujours » renforce la régularité avec \zh{每天}. \zh{也} « aussi » et \zh{还} « encore/en plus » ne conviennent pas ici.
    \item \textbf{Réponse C :} \zh{所以} (suǒyǐ) « donc, c'est pourquoi ». \zh{因为} exprime la cause, \zh{但是} l'opposition.
    \item \textbf{Réponse A :} \zh{不但...而且...} (búdàn... érqiě...) = « non seulement... mais aussi... ». Phrase : \zh{她不但聪明，而且勤劳。} (Tā búdàn cōngming, érqiě qínláo.) « Elle est non seulement intelligente, mais aussi travailleuse. »
\end{enumerate}
\end{corrige}

\begin{corrige}[Exercice 2 — Vrai ou Faux : compréhension du texte modèle]
\begin{enumerate}
    \item \textbf{Vrai.} Justification : \zh{以前，很多妇女只在家里做家务、带孩子。} « Autrefois, beaucoup de femmes restaient seulement à la maison pour faire le ménage et élever les enfants. » Le mot-clé est \zh{只} « seulement ».
    \item \textbf{Vrai.} Justification : \zh{她们不但可以上班、工作，而且可以当领导、当医生。} « Elles peuvent non seulement travailler, mais aussi devenir dirigeantes ou médecins. »
    \item \textbf{Faux.} Justification : \zh{社会上还有人认为女人应该以家庭为主。} « Dans la société, certaines personnes pensent encore que la femme doit se consacrer d'abord à la famille. »
    \item \textbf{Faux.} Justification : \zh{我们要继续努力，争取真正的男女平等。} « Nous devons continuer à lutter pour une véritable égalité » : l'égalité n'est donc pas encore atteinte.
\end{enumerate}
\end{corrige}

\begin{corrige}[Exercice 3 — Association caractères / pinyin / sens]
Correspondances correctes :
\begin{enumerate}
    \item \zh{妇女} — fùnǔ — femme (terme formel)
    \item \zh{地位} — dìwèi — statut, position sociale
    \item \zh{提高} — tígāo — augmenter, améliorer
    \item \zh{家务} — jiāwù — tâches ménagères
    \item \zh{领导} — lǐngdǎo — dirigeant / diriger
    \item \zh{争取} — zhēngqǔ — lutter pour, s'efforcer d'obtenir
    \item \zh{平等} — píngděng — égalitaire, égal
    \item \zh{社会} — shèhuì — société
\end{enumerate}
\textbf{Points de vigilance :} ne pas confondre \zh{妇女} fùnǔ (formel, collectif) et \zh{女人} nǚrén (courant) ; attention au 3e ton de \zh{领} lǐng et au 2e ton de \zh{平} píng.
\end{corrige}

\begin{corrige}[Exercice 4 — Pinyin et tons]
\begin{enumerate}
    \item \zh{妇} : 4e ton → fù ; mot complet : \zh{妇女} \textbf{fùnǔ}.
    \item \zh{务} : 4e ton → wù ; mot complet : \zh{家务} \textbf{jiāwù}.
    \item \zh{领} : 3e ton → lǐng ; mot complet : \zh{领导} \textbf{lǐngdǎo} (deux 3e tons : le premier se prononce comme un 2e ton à l'oral, mais s'écrit lǐngdǎo).
    \item \zh{平} : 2e ton → píng ; mot complet : \zh{平等} \textbf{píngděng}.
\end{enumerate}
\textbf{Rappel :} la règle du sandhi des 3e tons (3+3 → 2+3) concerne la prononciation, jamais l'écriture du pinyin.
\end{corrige}

\begin{corrige}[Exercice 5 — Thème : phrases à traduire en chinois]
\begin{enumerate}
    \item \zh{妇女的地位提高了。} \emph{Fùnǔ de dìwèi tígāo le.} « Le statut des femmes s'est amélioré. » — \zh{了} marque le changement d'état accompli.
    \item \zh{我妈妈每天做家务。} \emph{Wǒ māma měitiān zuò jiāwù.} « Ma mère fait les tâches ménagères tous les jours. » — Le complément de temps \zh{每天} se place avant le verbe.
    \item \zh{女人不但可以工作，而且可以当领导。} \emph{Nǚrén búdàn kěyǐ gōngzuò, érqiě kěyǐ dāng lǐngdǎo.} « Les femmes peuvent non seulement travailler, mais aussi être dirigeantes. » — \zh{不但...而且...} encadre les deux prédicats.
    \item \zh{因为她很累，所以今天晚上不想做饭。} \emph{Yīnwèi tā hěn lèi, suǒyǐ jīntiān wǎnshang bù xiǎng zuò fàn.} « Parce qu'elle est fatiguée, elle ne veut pas cuisiner ce soir. » — \zh{因为...所以...} peuvent se combiner en chinois (contrairement au français).
    \item \zh{虽然有进步，但是平等还没有完全实现。} \emph{Suīrán yǒu jìnbù, dànshì píngděng hái méiyǒu wánquán shíxiàn.} « Bien qu'il y ait des progrès, l'égalité n'est pas encore complètement réalisée. » — \zh{虽然...但是...} : concession + opposition.
\end{enumerate}
\end{corrige}

\begin{corrige}[Exercice 6 — Transformation : passé / présent / avis]
\begin{enumerate}
    \item \zh{以前，女人在家里做饭。} \emph{Yǐqián, nǚrén zài jiāli zuò fàn.} « Autrefois, les femmes cuisinaient à la maison. » — Le marqueur de temps en tête suffit ; pas de conjugaison en chinois.
    \item \zh{现在，女人在家里做饭。} \emph{Xiànzài, nǚrén zài jiāli zuò fàn.}
    \item \zh{现在，女人不在家里做饭。} \emph{Xiànzài, nǚrén bú zài jiāli zuò fàn.} « Maintenant, les femmes ne cuisinent pas à la maison. » — \zh{不} se place avant \zh{在}.
    \item \zh{女人不但在家里做饭，而且还在外面工作。} \emph{Nǚrén búdàn zài jiāli zuò fàn, érqiě hái zài wàimiàn gōngzuò.} « Les femmes non seulement cuisinent à la maison, mais travaillent aussi dehors. »
\end{enumerate}
\textbf{Point de vigilance :} ne jamais ajouter \zh{了} systématiquement pour « traduire » le passé français ; \zh{以前} seul situe déjà l'action dans le passé.
\end{corrige}

\begin{corrige}[Exercice 7 — Texte à trous : connecteurs logiques]
Texte complété :
\begin{quote}
\zh{因为现代妇女受教育程度高了，所以她们找到好工作容易多了。但是，很多妇女下班后还要做家务，而且很累。虽然社会在进步，我们还是要争取男女平等。}
\end{quote}
\emph{Yīnwèi xiàndài fùnǔ shòu jiàoyù chéngdù gāo le, suǒyǐ tāmen zhǎodào hǎo gōngzuò róngyì duō le. Dànshì, hěn duō fùnǔ xiàbān hòu hái yào zuò jiāwù, érqiě hěn lèi. Suīrán shèhuì zài jìnbù, wǒmen háishi yào zhēngqǔ nán nǚ píngděng.}

Traduction : « Parce que le niveau d'éducation des femmes modernes s'est élevé, il leur est beaucoup plus facile de trouver un bon emploi. Mais beaucoup de femmes, après le travail, doivent encore faire le ménage, et elles sont très fatiguées. Bien que la société progresse, nous devons encore lutter pour l'égalité hommes-femmes. »

\textbf{Justification :} cause (\zh{因为}) → conséquence (\zh{所以}) → opposition (\zh{但是}) → addition (\zh{而且}) → concession (\zh{虽然}).
\end{corrige}

\begin{corrige}[Exercice 8 — Remise en ordre : structure de phrase]
\begin{enumerate}
    \item \zh{现代妇女的地位提高了。} \emph{Xiàndài fùnǔ de dìwèi tígāo le.} « Le statut des femmes modernes s'est amélioré. » — Schéma : Déterminant (\zh{现代}) + \zh{的} + Sujet + Verbe + \zh{了}.
    \item \zh{男人和女人应该一起做家务。} \emph{Nánrén hé nǚrén yīnggāi yìqǐ zuò jiāwù.} « Les hommes et les femmes doivent faire le ménage ensemble. » — \zh{应该} (modal) avant l'adverbe \zh{一起} et le verbe.
    \item \zh{她下班后工作了，所以很累。} ou mieux : \zh{她下班后还要工作，所以很累。} Avec les mots donnés : \zh{她下班后工作了，所以很累。} \emph{Tā xiàbān hòu gōngzuò le, suǒyǐ hěn lèi.} « Après le travail, elle a encore travaillé, donc elle est très fatiguée. » — Ordre : Sujet + Temps (\zh{下班后}) + Verbe.
    \item \zh{妇女不但可以当医生，而且可以当领导。} \emph{Fùnǔ búdàn kěyǐ dāng yīshēng, érqiě kěyǐ dāng lǐngdǎo.} « Les femmes peuvent non seulement devenir médecins, mais aussi dirigeantes. »
\end{enumerate}
\end{corrige}

\begin{corrige}[Exercice 9 — Compréhension de texte et langue : type Bac]
\textbf{1. Compréhension globale (2 pts) :} Le texte compare l'évolution du statut de la femme au Cameroun et en Chine : progrès réels (éducation, métiers), mais persistance de la double charge domestique ; la vraie égalité exige des actes dans la famille. (1 pt thème + 1 pt mention des deux pays.)

\textbf{2. Détails (3 pts, 1 pt par métier) :} \zh{当老师} enseignante, \zh{护士} infirmière, \zh{律师} avocate.

\textbf{3. Comparaison (2 pts) :} Dans les deux pays, les femmes qui travaillent doivent encore, en rentrant, laver le linge, cuisiner et aider aux devoirs (\zh{洗衣服、做饭、辅导作业}) ; les hommes ne partagent pas assez les tâches (\zh{男人分担家务的还不够多}).

\textbf{4. Vocabulaire en contexte (2 pts) :} \zh{分担} (fēndān) = « partager, se répartir (une charge) ». Synonyme possible : \zh{一起承担} ou \zh{帮忙做} (aider à faire).

\textbf{5. Grammaire (2 pts) :} structure de concession : \zh{但是，两个国家都有共同的问题} (opposition) ; la concession explicite « bien que » est \zh{虽然...但是...} ; dans le texte on relève l'opposition \zh{但是...} : \emph{Dànshì, liǎng ge guójiā dōu yǒu gòngtóng de wèntí.} « Mais les deux pays ont un problème commun. » (Accepter aussi la reformulation avec \zh{虽然}.)

\textbf{6. Expression personnelle (4 pts) :} Sens : l'égalité réelle ne se limite pas aux lois ; elle dépend des gestes quotidiens dans la famille. Exemple concret : à la maison, le père lave la vaisselle et aide les enfants à faire leurs devoirs, comme dans le texte \zh{我爸爸也开始学做饭、洗碗}. (2 pts compréhension + 2 pts exemple.)

\textbf{Points de vigilance :} répondre en français ; citer le texte en caractères pour justifier ; ne pas inventer d'informations absentes du texte.
\end{corrige}

\begin{corrige}[Exercice 10 — Thème et Version : type Bac]
\textbf{Partie A — Thème (proposition de traduction, 10 pts) :}
\begin{quote}
\zh{以前，很多女人在家里带孩子。现在，她们可以在外面工作，但是还要做很多家务。我认为男人应该分担家里的工作。这样，女人就不会太累了。}
\end{quote}
\emph{Yǐqián, hěn duō nǚrén zài jiāli dài háizi. Xiànzài, tāmen kěyǐ zài wàimiàn gōngzuò, dànshì hái yào zuò hěn duō jiāwù. Wǒ rènwéi nánrén yīnggāi fēndān jiāli de gōngzuò. Zhèyàng, nǚrén jiù bú huì tài lèi le.}

\textbf{Barème indicatif :} 2 pts par phrase (ordre des mots 1 pt, lexique et mots-outils 1 pt).
\textbf{Points de vigilance :} \zh{带孩子} « s'occuper des enfants » ; \zh{在外面工作} « travailler dehors » ; \zh{认为} « penser que » ; \zh{这样} « ainsi » ; ne pas oublier \zh{了} de changement d'état final.

\textbf{Partie B — Version (proposition de traduction, 10 pts) :}
\begin{quote}
« Autrefois, ma mère ne savait pas lire et ne pouvait que cultiver la terre à la campagne. Plus tard, elle a étudié avec effort et a appris beaucoup de caractères. Maintenant, elle travaille comme infirmière à l'hôpital, et son salaire n'est pas bas. Mon père, lui aussi, a commencé à apprendre à cuisiner et à laver la vaisselle. Toute notre famille trouve que la vie actuelle est bien plus heureuse qu'avant. »
\end{quote}
\textbf{Barème indicatif :} 2 pts par phrase (fidélité 1 pt, qualité du français 1 pt).
\textbf{Points de vigilance :} \zh{不识字} « ne pas savoir lire, être illettré » ; \zh{农村} « campagne » ; \zh{工资不低} « le salaire n'est pas bas » ; \zh{比以前幸福多了} « bien plus heureux qu'avant » (comparatif \zh{比}).
\end{corrige}

\begin{corrige}[Exercice 11 — Expression écrite : sujet type Bac]
\textbf{Proposition de rédaction modèle (environ 95 caractères) :}
\begin{quote}
\zh{以前，我们家的女人只在家里做饭、洗衣服、带孩子。现在，情况变了：我妈妈在医院当护士，我姐姐在大学学习。她们不但工作，而且工作得很好。但是，妈妈下班后还要做很多家务，所以很累。我认为，男人应该分担家务。虽然进步很大，我们还要继续努力，争取真正的男女平等。}
\end{quote}
\emph{Yǐqián, wǒmen jiā de nǚrén zhǐ zài jiāli zuò fàn, xǐ yīfu, dài háizi. Xiànzài, qíngkuàng biàn le : wǒ māma zài yīyuàn dāng hùshi, wǒ jiějie zài dàxué xuéxí. Tāmen búdàn gōngzuò, érqiě gōngzuò de hěn hǎo. Dànshì, māma xiàbān hòu hái yào zuò hěn duō jiāwù, suǒyǐ hěn lèi. Wǒ rènwéi, nánrén yīnggāi fēndān jiāwù. Suīrán jìnbù hěn dà, wǒmen hái yào jìxù nǔlì, zhēngqǔ zhēnzhèng de nán nǚ píngděng.}

Traduction : « Autrefois, les femmes de notre famille restaient à la maison pour cuisiner, laver le linge et élever les enfants. Maintenant, la situation a changé : ma mère est infirmière à l'hôpital, ma sœur étudie à l'université. Elles travaillent, et elles travaillent bien. Mais après le travail, maman doit encore faire beaucoup de tâches ménagères, donc elle est très fatiguée. Je pense que les hommes doivent partager les tâches. Bien que les progrès soient grands, nous devons continuer nos efforts pour obtenir une véritable égalité hommes-femmes. »

\textbf{Vocabulaire de la leçon utilisé (5 mots minimum) :} \zh{家务}、\zh{地位/情况}、\zh{分担}、\zh{争取}、\zh{平等}、\zh{进步}.
\textbf{Connecteurs utilisés (3 minimum) :} \zh{但是}、\zh{所以}、\zh{不但...而且...}、\zh{虽然}.

\textbf{Barème indicatif (15 pts) :} Contenu 4 pts (passé/présent/reste à faire + exemple familial camerounais) ; Vocabulaire 3 pts (5 mots du thème exacts) ; Grammaire 4 pts (ordre Sujet+Temps+Lieu+Verbe, 3 connecteurs, \zh{了/的/得} corrects) ; Caractères 2 pts (lisibles, traits corrects) ; Cohérence et ponctuation 2 pts.

\textbf{Points de vigilance :}
\begin{itemize}
    \item Ne pas écrire \zh{*我妈妈是护士在医院} : le lieu précède le verbe → \zh{我妈妈在当医院护士} est fautif, dire \zh{我妈妈在医院当护士}.
    \item \zh{工作得很好} avec \zh{得} (complément de manière), pas \zh{的}.
    \item Éviter la double négation française ; \zh{不累} suffit pour « ne pas être fatiguée ».
    \item Compter les caractères : rester entre 80 et 100.
\end{itemize}
\end{corrige}

\newpage

\coursec{Fiche bilan}

\begin{popfigure}
\begin{tikzpicture}[
  mindmap,
  concept color=popblue,
  every node/.style={concept, execute at begin node=\hskip0pt},
  root concept/.style={concept color=popblue, font=\large\bfseries, text width=3.5cm, align=center, inner sep=4pt},
  level 1 concept/.append style={sibling angle=60, level distance=4.5cm, font=\normalsize\bfseries, text width=2.2cm, align=center, inner sep=3pt},
  level 2 concept/.append style={level distance=3cm, font=\small, text width=2cm, align=center, inner sep=2pt},
  Vocab/.style={concept color=popblue},
  Gram/.style={concept color=poporange},
  Conn/.style={concept color=popgreen},
  Ecri/.style={concept color=poppurple},
  Trad/.style={concept color=poppink},
  Prod/.style={concept color=popgold},
]
\node[root concept, align=center]{Leçon 3\\La femme d'aujourd'hui\\当代女性}
  child[Vocab] { node {Vocabulaire clé}
    child { node {女性 平等 就业} }
    child { node {歧视 独立 平衡} }
  }
  child[Gram] { node[align=center]{Grammaire\\Structures}
    child { node {比 越来越} }
    child { node {不但…而且… 虽然…但是…} }
    child { node {被/让/使 了/过/着} }
  }
  child[Conn] { node[align=center]{Connecteurs\\logiques}
    child { node {首先 其次 最后} }
    child { node {此外 因此 然而} }
  }
  child[Ecri] { node[align=center]{Écriture\\Caractères}
    child { node {女 母 好 她 媽} }
    child { node {平 等 就 业 场} }
    child { node {歧 视 独 立 婚} }
  }
  child[Trad] { node {Thème / Version}
    child { node {FR $\to$ ZH : ordre mots} }
    child { node {ZH $\to$ FR : aspects} }
  }
  child[Prod] { node[align=center]{Production\\écrite \& Socio}
    child { node {Paragraphe 80-100 car.} }
    child { node {Exemples CM / Chine} }
  };
\end{tikzpicture}
\legende{Carte mentale de la Leçon 3 : Expression écrite — Situation de la femme}
\end{popfigure}

\begin{bilan}

\textbf{Lexique clé (HSK 3--4)}\par
\begin{center}
\begin{tabularx}{\linewidth}{|X|X|X|X|}
\hline
\rowcolor{popblueL}
\textbf{Caractères} & \textbf{Pinyin} & \textbf{Nature} & \textbf{Traduction} \\
\hline
女性 & nǚxìng & n. & femme, sexe féminin \\
地位 & dìwèi & n. & statut, position sociale \\
平等 & píngděng & adj./n. & égal, égalité \\
就业 & jiùyè & v. & trouver un emploi, travailler \\
职场 & zhíchǎng & n. & milieu professionnel, lieu de travail \\
歧视 & qíshì & v./n. & discriminer, discrimination \\
独立 & dúlì & adj./v. & indépendant, s'émanciper \\
事业 & shìyè & n. & carrière, cause, entreprise \\
家庭 & jiātíng & n. & famille, foyer \\
平衡 & pínghéng & v./n. & équilibrer, équilibre \\
机会 & jīhuì & n. & occasion, opportunité \\
挑战 & tiǎozhàn & v./n. & défier, défi \\
进步 & jìnbù & v./n. & progresser, progrès \\
法律 & fǎlǜ & n. & loi, législation \\
观念 & guānniàn & n. & concept, mentalité, idée \\
社会 & shèhuì & n. & société \\
\hline
\end{tabularx}
\end{center}

\textbf{Structures grammaticales indispensables}\par
\begin{center}
\begin{tabularx}{\linewidth}{|X|X|X|}
\hline
\rowcolor{poporangeL}
\textbf{Structure / Règle} & \textbf{Exemple (Caractères + Pinyin)} & \textbf{Traduction} \\
\hline
\cle{比} : A 比 B + Adj (comparatif) &
现在的女性 \textbf{比} 过去 独立。 &
Les femmes d'aujourd'hui sont \textbf{plus} indépendantes qu'autrefois. \\
\hline
\cle{越来越} + Adj (évolution progressive) &
社会 对 女性 的 看法 \textbf{越来越} 开放。 &
La vision de la société sur les femmes devient \textbf{de plus en plus} ouverte. \\
\hline
\cle{不但…而且…} (coordination progressive) &
她 \textbf{不但} 事业 成功，\textbf{而且} 家庭 幸福。 &
Elle a \textbf{non seulement} réussi sa carrière, \textbf{mais encore} un foyer heureux. \\
\hline
\cle{虽然…，但是/可是…} (concession) &
\textbf{虽然} 进步 很大，\textbf{但是} 仍有 挑战。 &
\textbf{Bien que} le progrès soit grand, \textbf{il reste} des défis. \\
\hline
\cle{被 / 让 / 叫 / 使} (passif / causatif) &
女性 \textbf{被} 重视；\textbf{让/使} 她们 受益。 &
Les femmes \textbf{sont} valorisées ; \textbf{faire en sorte qu'}elles en bénéficient. \\
\hline
Aspect \cle{了} (changement d'état / action terminée) &
现在 法律 \textbf{保护} 了 女性 权益。 &
Maintenant, la loi \textbf{a protégé / protège} les droits des femmes (fait accompli / nouvelle situation). \\
\hline
Aspect \cle{过} (expérience vécue) &
她 \textbf{当过} 经理。 &
Elle \textbf{a été} directrice (une fois dans sa vie). \\
\hline
Aspect \cle{着} (état durable) &
她 \textbf{坚持着} 奋斗。 &
Elle \textbf{persiste} dans la lutte (action continue). \\
\hline
\cle{把} 结构 (disposition de l'objet) &
我们 应该 \textbf{把} 机会 \textbf{给} 女性。 &
Nous devrions \textbf{donner} les opportunités \textbf{aux} femmes. \\
\hline
\end{tabularx}
\end{center}

\textbf{Connecteurs logiques pour structurer le texte}\par
\begin{itemize}
\item \textbf{Énumération :} 首先 (premier) — 其次 (deuxièmement) — 再次 / 然后 (ensuite) — 最后 (enfin).
\item \textbf{Addition :} 此外 (de plus) — 另外 (par ailleurs) — 而且 (qui plus est).
\item \textbf{Conséquence :} 因此 (par conséquent) — 所以 (donc) — 从而 (afin que / de sorte que).
\item \textbf{Opposition / Concession :} 但是 / 可是 (mais) — 然而 (cependant) — 虽然…但是… (bien que… pourtant).
</{itemize}

\textbf{Méthodes express}
\begin{enumerate}
\item \textbf{Thème (FR $\to$ ZH) :} 1) Identifier le Sujet, le Temps, le Lieu. 2) Appliquer l'ordre canonique \cle{Sujet + Temps + Lieu + Verbe + Objet/Complément}. 3) Choisir l'aspect (了/过/着) et la voix (被/让). 4) Vérifier les classificateurs et la place des compléments (degré, durée).
\item \textbf{Version (ZH $\to$ FR) :} 1) Segmenter en groupes (Sujet / Prédicat / Objet). 2) Repérer les mots-outils (了/过/着/被/比/连…都…). 3) Traduire le sens global, pas mot à mot. 4) Adapter le temps français (passé composé, imparfait, présent) selon l'aspect.
\item \textbf{Écriture caractères :} Respecter l'ordre des traits (横 avant 竖, gauche avant droite, haut avant bas, cadre avant contenu, trait central avant ailes). Identifier le \cle{radical} (clé) pour le dictionnaire : 女 (femme), 扌(main), 忄(cœur), 讠(parole).
\item \textbf{Production écrite :} 1) Plan en 3 parties (Situation actuelle / Problèmes restants / Solutions/Avenir). 2) Utiliser 5 connecteurs minimum. 3) Intégrer 8--10 mots du lexique cible. 4) Viser 80--100 caractères. 5) Relire : tons, ordre mots, ponctuation chinoise (，。；！？).
\end{enumerate}

\end{itemize}
\end{bilan}

\begin{aretenir}[Checklist avant le Bac]
$\square$ Je connais par cœur le \textbf{vocabulaire clé} (15 mots : 女性, 平等, 就业, 歧视, 独立, 平衡, 进步, 观念…).
$\square$ J'écris correctement les \textbf{15 caractères cibles} (traits, radicaux, ordre, caractères proches : 末/未, 失/矢, 午/牛).
$\square$ Je maîtrise la structure \textbf{比} pour comparer la situation passée et présente.
$\square$ J'utilise \textbf{越来越 + adjectif} pour décrire l'évolution positive des mentalités.
$\square$ J'articule mon paragraphe avec \textbf{au moins 5 connecteurs} (首先, 此外, 然而, 因此, 最后).
$\square$ Je distingue \textbf{被} (passif formel/écrit) de \textbf{让/叫} (passif oral/causatif) et de \textbf{使} (causatif formel).
$\square$ Je place correctement \textbf{了} (changement d'état / action close), \textbf{过} (expérience), \textbf{着} (duratif).
$\square$ Je respecte l'ordre des mots \textbf{Sujet + Temps + Lieu + Verbe + Objet} en thème (FR $\to$ ZH).
$\square$ Je rends les aspects verbaux chinois par les \textbf{temps français adaptés} en version (ZH $\to$ FR).
$\square$ Je produis un \textbf{texte cohérent de 80--100 caractères} structuré en 3 parties logiques.
$\square$ J'intègre un \textbf{exemple socioculturel précis} (ex: loi camerounaise, femme ministre, congé maternité en Chine).
$\square$ Je relis systématiquement : \textbf{tons, pinyin, ponctuation pleine chasse, radicaux, classificateurs}.
\end{aretenir}

\findecours

\end{document}
